% Relatório técnico: inserção de elementos de interface em malhas 2D.
% Compilar com:  tectonic -X compile relatorio.tex   (ou xelatex/lualatex)
% As figuras de malha em figs/ são geradas por tools/gen.py a partir de execuções reais do código.
\documentclass[11pt,a4paper]{article}

\usepackage[a4paper,margin=2.3cm]{geometry}
\usepackage{amsmath}
\usepackage{libertinus}
\usepackage[brazilian]{babel}

\usepackage{microtype}
\usepackage{xcolor}
\usepackage{tikz}
\usepackage{booktabs,tabularx,multirow,array}
\usepackage{enumitem}
\usepackage{caption,subcaption}
\usepackage{algorithm}
\usepackage[noend]{algpseudocode}
\usepackage{listings}
\usepackage[most]{tcolorbox}
\usepackage{float}
\usepackage[section]{placeins}
\usepackage[hidelinks]{hyperref}
\usepackage[nameinlink,noabbrev,brazilian]{cleveref}

% Shared TikZ styles for the report figures (generated figures in figs/ use these names).
\usetikzlibrary{arrows.meta, patterns, positioning, shapes.geometric, shapes.misc, calc, decorations.pathreplacing,
  decorations.markings, fit, backgrounds}

\definecolor{ceblue}{RGB}{214,240,250}
\definecolor{ieorange}{RGB}{250,200,150}
\definecolor{vblue}{RGB}{52,92,160}
\definecolor{iehl}{RGB}{255,170,40}

\tikzset{
  cefill/.style={fill=ceblue, draw=black!75, line width=0.45pt, line join=round},
  iefill/.style={fill=ieorange, draw=orange!55!black, line width=0.45pt, line join=round},
  holefill/.style={pattern=north east lines, pattern color=black!55},
  seledge/.style={draw=red!80!black, line width=1.3pt},
  curedge/.style={draw=red, line width=2.4pt},
  edgenum/.style={text=red!80!black, font=\tiny\bfseries, inner sep=0.5pt},
  hearrow/.style={-{Stealth[length=3.2pt, width=2.6pt]}, line width=0.55pt},
  helabel/.style={font=\scriptsize, inner sep=0.3pt},
  vreg/.style={circle, fill=vblue, text=white, inner sep=0pt, font=\scriptsize\bfseries},
  vie/.style={circle, fill=vblue, text=iehl, inner sep=0pt, font=\scriptsize\bfseries},
  vcol/.style={circle, fill=vblue, text=iehl, draw=red, line width=1.1pt, inner sep=0pt, font=\scriptsize\bfseries},
  dot/.style={minimum size=3.2pt, inner sep=0pt, line width=0.6pt},
  elemlabel/.style={font=\footnotesize\sffamily\bfseries, text=black!55},
  % hand-drawn diagrams
  box/.style={draw=black!70, rounded corners=2pt, fill=white, align=center, font=\small, inner sep=4pt},
  decision/.style={draw=black!70, diamond, aspect=2.2, fill=yellow!12, align=center, font=\small, inner sep=1pt},
  opbox/.style={draw=black!70, rounded corners=3pt, align=center, font=\small\bfseries, inner sep=5pt, minimum width=2.5cm},
  state/.style={draw=black!70, rounded corners=6pt, align=center, font=\small, inner sep=5pt, minimum width=2.3cm,
    minimum height=1cm},
  flow/.style={-{Stealth[length=5pt]}, line width=0.7pt},
}


\captionsetup{font=small,labelfont=bf,format=hang}
\setlist{itemsep=2pt,topsep=4pt}
\setlength{\parskip}{3pt}
\setlength{\emergencystretch}{3em}
\captionsetup[subfigure]{justification=centering}

% --- pseudo-código em português -------------------------------------------------------------------------------------
\floatname{algorithm}{Algoritmo}
\renewcommand{\listalgorithmname}{Lista de algoritmos}
\algrenewcommand\algorithmicif{\textbf{se}}
\algrenewcommand\algorithmicthen{\textbf{então}}
\algrenewcommand\algorithmicelse{\textbf{senão}}
\algrenewcommand\algorithmicfor{\textbf{para}}
\algrenewcommand\algorithmicforall{\textbf{para cada}}
\algrenewcommand\algorithmicdo{\textbf{faça}}
\algrenewcommand\algorithmicwhile{\textbf{enquanto}}
\algrenewcommand\algorithmicrepeat{\textbf{repita}}
\algrenewcommand\algorithmicuntil{\textbf{até que}}
\algrenewcommand\algorithmicreturn{\textbf{retorne}}
\algrenewcommand\algorithmicprocedure{\textbf{procedimento}}
\algrenewcommand\algorithmicfunction{\textbf{função}}
\algrenewcommand\algorithmicrequire{\textbf{Entrada:}}
\algrenewcommand\algorithmicensure{\textbf{Saída:}}
\algrenewcommand{\algorithmiccomment}[1]{\hfill{\color{black!55}\(\triangleright\) #1}}
\crefname{algorithm}{Algoritmo}{Algoritmos}
\crefname{figure}{Figura}{Figuras}
\crefname{table}{Tabela}{Tabelas}
\crefname{section}{Seção}{Seções}
\crefname{subsection}{Seção}{Seções}
\crefname{appendix}{Apêndice}{Apêndices}

% --- notação ---------------------------------------------------------------------------------------------------------
\newcommand{\code}[1]{\texttt{\small #1}}
\newcommand{\V}{\mathrm{V}}
\newcommand{\Op}{\mathrm{O}}
\newcommand{\nxt}{\mathrm{next}}
\newcommand{\prv}{\mathrm{prev}}
\newcommand{\BORDER}{\textsc{border}}
\newcommand{\COLLAPSED}{\textsc{collapsed}}
\newcommand{\opname}[1]{\textsf{\bfseries #1}}

\lstset{basicstyle=\ttfamily\footnotesize, keywordstyle=\color{blue!60!black}\bfseries, commentstyle=\color{black!55},
  language=C++, frame=single, framerule=0.3pt, rulecolor=\color{black!30}, columns=fullflexible, breaklines=true,
  showstringspaces=false, xleftmargin=4pt, xrightmargin=4pt}

\newtcolorbox{nota}[1][]{colback=blue!3, colframe=blue!40!black, boxrule=0.5pt, arc=2pt, left=6pt, right=6pt,
  top=4pt, bottom=4pt, fonttitle=\bfseries\small, title={#1}}
\newtcolorbox{alerta}[1][]{colback=red!3, colframe=red!50!black, boxrule=0.5pt, arc=2pt, left=6pt, right=6pt,
  top=4pt, bottom=4pt, fonttitle=\bfseries\small, title={#1}}

% Versão do relatório: incrementar a cada nova versão e registrar em historico/ e na tabela de histórico.
\newcommand{\versao}{0.3}

\title{\textbf{Inserção adaptativa de elementos de interface}\\ \textbf{em malhas 2D}\\[4pt]
  \large Operadores topológicos sobre estruturas de dados compactas baseadas em índices}
\author{Jeferson Coelho\\ \small Tecgraf/PUC-Rio\thanks{Texto redigido com o auxílio do Claude Code a partir do
  código do ramo \code{feature/edge-element-insertion} (commit \code{066cbde}), da apresentação
  \code{doc/ppt/Operators\_Template} e do artigo da IBHM~\cite{ibhm}. As figuras de malha foram geradas executando o
  próprio código (\cref{app:repro}). As versões anteriores estão em \code{doc/report/historico/}.}}
\date{Relatório técnico — versão \versao{} — outubro de 2026}

\begin{document}
\maketitle

\begin{abstract}
\noindent Este relatório documenta o primeiro protótipo, em 2D, de um algoritmo para inserir elementos de interface de
espessura nula em malhas de elementos finitos representadas por estruturas de dados topológicas compactas (CHE e
IBHM). O algoritmo recebe um conjunto de arestas e, para cada uma, cria um quadrilátero degenerado e repara a vizinhança
de seus dois vértices aplicando um entre quatro operadores --- \opname{Canônico}, \opname{Split Element},
\opname{Expand Edge} e \opname{Open Hole} ---, escolhido exclusivamente a partir do \emph{estado do vértice}.
Descrevemos o problema, a estrutura IBHM, a arquitetura do algoritmo, cada operador com seus detalhes de implementação e
o algoritmo de travessia, tal como implementados na classe \code{IBHMInterfaceElementBuilder}, que é a base a ser
evoluída. Cada caso é ilustrado com figuras produzidas pelo próprio código. Em seguida, confrontamos a implementação com
seus comentários e com o artigo da IBHM e listamos os defeitos em aberto. Uma verificação
experimental com 800 malhas aleatórias (triângulos, quadriláteros, malhas híbridas e polígonos de seis lados) confirma, em
todas as configurações suportadas, a validade topológica, a contagem exata de vértices novos e a independência da ordem
de inserção (inclusive na inserção em lote). Ela também revela falhas que precisam ser corrigidas: vértices de borda com
mais de uma aresta selecionada, entradas que contêm arestas de elementos de interface e índices de half-edge inválidos.
\end{abstract}

\begin{center}\small
\begin{tabular}{clp{11.2cm}}
\toprule
Versão & Data & Mudanças\\
\midrule
0.1 & 18/10/2026 & Primeira versão, baseada na classe \code{CHEInterfaceElementBuilder}.\\
0.2 & 18/10/2026 & Análise refeita sobre a classe \code{IBHMInterfaceElementBuilder}, base a ser evoluída; a implementação
baseada na CHE deixa de ser considerada (a teoria da CHE foi mantida). Nova figura com os vetores da IBHM antes e depois
de uma inserção; verificação com 800 malhas de triângulos, quadriláteros, híbridas e com polígonos, inclusive inserção em
lote; novos casos-limite (índices inválidos, malha com furo interno); lista de defeitos restrita à IBHM.\\
0.3 & 18/10/2026 & Sai o confronto com a apresentação de projeto, que não será mais mantida. Nova figura com as
hipóteses $x_1$ e $x_3$ do Split Element. Comentários do código corrigidos (B1--B9, B11--B14) e reserva exata de memória
em \code{reserveSpaceForNodes} com \code{addPoint} recebendo um \code{std::vector} (B2, parte do I4); defeitos
atualizados (seção~7.4, antiga 7.5). Ferramentas de teste documentadas e explicadas no apêndice, com um teste novo para
as coordenadas.\\
\bottomrule
\end{tabular}
\end{center}

\newpage
\tableofcontents
\newpage

\section{Introdução}
\label{sec:intro}

\subsection{Elementos de interface}

Em muitos problemas de engenharia o comportamento que mais importa acontece em superfícies de descontinuidade:
fraturas e falhas geológicas, juntas em maciços rochosos, contatos entre materiais distintos, caminhos de propagação de
trincas. No método dos elementos finitos, uma forma clássica de representar essas descontinuidades é inserir, ao longo
delas, \emph{elementos de interface} --- também chamados de elementos de espessura nula ou, quando associados a uma lei
de tração--separação, elementos coesivos~\cite{paulino}.

Em 2D, um elemento de interface linear é um quadrilátero de quatro nós cujos nós vêm aos pares geometricamente
coincidentes (\cref{fig:ie-conceito}). Dois nós pertencem a um lado da descontinuidade e dois ao outro. Como a espessura
é nula, a grandeza cinemática relevante deixa de ser a deformação e passa a ser o \emph{salto de deslocamento} entre os
dois lados, $[\![\mathbf{u}]\!] = \mathbf{u}^{+} - \mathbf{u}^{-}$, que mede abertura e deslizamento da interface.

\begin{figure}[ht]
\centering
\begin{tikzpicture}[scale=1.05, every node/.style={font=\small}]
  % (a) contínuo
  \begin{scope}
    \filldraw[cefill] (0,0) rectangle (3,1); \filldraw[cefill] (0,1) rectangle (3,2);
    \draw[seledge] (0,1) -- (3,1);
    \node[elemlabel] at (1.5,1.5) {U}; \node[elemlabel] at (1.5,0.5) {L};
    \node[vreg, minimum size=0.38cm] at (0,1) {1}; \node[vreg, minimum size=0.38cm] at (3,1) {2};
    \node at (1.5,-0.5) {(a) malha contínua};
  \end{scope}
  % (b) com elemento de interface (espessura exagerada)
  \begin{scope}[xshift=4.2cm]
    \filldraw[cefill] (0,0) rectangle (3,1); \filldraw[iefill] (0,1) rectangle (3,1.35);
    \filldraw[cefill] (0,1.35) rectangle (3,2.35);
    \node[elemlabel] at (1.5,1.85) {U}; \node[elemlabel] at (1.5,0.5) {L};
    \node[font=\scriptsize\bfseries, text=orange!40!black] at (1.5,1.175) {I};
    \node[vie, minimum size=0.36cm] at (0,1) {1}; \node[vie, minimum size=0.36cm] at (3,1) {2};
    \node[vie, minimum size=0.36cm] at (3,1.35) {3}; \node[vie, minimum size=0.36cm] at (0,1.35) {4};
    \draw[decorate, decoration={brace, amplitude=3pt}] (3.3,1.4) -- (3.3,0.95)
      node[midway, right=3pt, align=left, font=\scriptsize] {espessura nula\\(exagerada)};
    \node at (1.5,-0.5) {(b) elemento de interface};
  \end{scope}
  % (c) abertura
  \begin{scope}[xshift=9.6cm]
    \filldraw[cefill] (0,0) rectangle (3,1);
    \filldraw[iefill] (0,1) -- (3,1) -- (3.25,1.55) -- (0.15,1.3) -- cycle;
    \filldraw[cefill] (0.15,1.3) -- (3.25,1.55) -- (3.25,2.55) -- (0.15,2.3) -- cycle;
    \node[elemlabel] at (1.7,1.95) {U}; \node[elemlabel] at (1.5,0.5) {L};
    \draw[-{Stealth[length=4pt]}, red!70!black, line width=0.8pt] (3,1) -- (3.25,1.55);
    \draw[-{Stealth[length=4pt]}, red!70!black, line width=0.8pt] (0,1) -- (0.15,1.3);
    \node[text=red!70!black, font=\scriptsize] at (4.05,1.25) {$[\![\mathbf{u}]\!]=\mathbf{u}^{+}-\mathbf{u}^{-}$};
    \node at (1.6,-0.5) {(c) abertura da interface};
  \end{scope}
\end{tikzpicture}
\caption{O elemento de interface. (a) Dois elementos contínuos compartilham uma aresta (em vermelho). (b) A aresta é
\emph{aberta}: os nós 1 e 2 são duplicados em 4 e 3, e entre U e L surge um quadrilátero de quatro nós com
$\mathbf{x}_1=\mathbf{x}_4$ e $\mathbf{x}_2=\mathbf{x}_3$ (desenhado com espessura exagerada). (c) Durante a simulação os
dois lados podem se separar; o elemento mede o salto de deslocamento.}
\label{fig:ie-conceito}
\end{figure}

Há duas estratégias para colocar esses elementos na malha. Na estratégia \emph{a priori}, todos os elementos de
interface são criados antes da análise, ao longo de todas as superfícies em que uma descontinuidade pode ocorrer. Na
estratégia \emph{adaptativa}, eles são inseridos durante a simulação, apenas onde e quando um critério (de tensão, de
dano etc.) indica que a descontinuidade se formou~\cite{paulino}. A segunda estratégia evita a rigidez artificial e o
custo de milhares de elementos inúteis, mas exige um algoritmo que altere a topologia da malha de forma incremental,
local e eficiente, sobre uma malha que \emph{já contém} elementos de interface.

\subsection{O problema da inserção}
\label{sec:problema}

\begin{nota}[Problema]
\textbf{Dados:} uma malha 2D orientada $\mathcal{M}$ e um conjunto $S$ de arestas internas (cada uma compartilhada por
dois elementos).\\
\textbf{Obter:} uma malha $\mathcal{M}'$ em que cada aresta de $S$ foi \emph{aberta}: os dois elementos que a
compartilhavam passam a se ligar através de um novo elemento de interface de quatro nós, e os vértices são duplicados
apenas o necessário para que a topologia continue consistente.
\end{nota}

A \cref{fig:exemplo-intro} mostra a entrada e a saída para o exemplo usado ao longo de todo o texto: uma malha híbrida
de quadriláteros e triângulos, com 12 arestas selecionadas e numeradas na ordem de inserção. Parece simples, mas há
várias dificuldades:
\begin{enumerate}
\item \textbf{Quantas cópias de cada vértice?} Um vértice interno tocado por uma única aresta selecionada não precisa
ser duplicado geometricamente (a trinca termina nele). Com duas arestas, ele vira dois. Com três ou mais, surge uma
\emph{junção} (em T, em X, em estrela), e os elementos de interface não conseguem fechar a vizinhança sozinhos:
aparece um \emph{buraco} poligonal, visível como as regiões hachuradas na \cref{fig:exemplo-intro}b.
\item \textbf{Quais elementos vão para cada cópia?} A estrela do vértice precisa ser dividida em setores delimitados
pelas arestas selecionadas, e cada setor deve passar a usar uma cópia diferente.
\item \textbf{Borda.} Na borda da malha a abertura não precisa ser fechada: o vértice sempre ganha uma cópia a mais.
\item \textbf{Incrementalidade.} Em uma análise adaptativa as arestas chegam aos poucos e em qualquer ordem. O
resultado não deveria depender dessa ordem.
\item \textbf{Estruturas compactas.} Estruturas baseadas em índices, como a CHE e a IBHM~\cite{che,ibhm}, armazenam a
malha em vetores de inteiros, sem ponteiros. Os próprios autores da IBHM apontam como limitação o uso em remalhamento
dinâmico~\cite{ibhm}. Inserir elementos é exatamente esse tipo de alteração.
\item \textbf{Extensão para 3D.} O objetivo final é inserir elementos em faces de malhas 3D. As decisões tomadas em 2D
devem ser transportáveis.
\end{enumerate}

\begin{figure}[ht]
\centering
\begin{subfigure}[b]{0.42\linewidth}\centering
  \begin{tikzpicture}[scale=1.250, every node/.style={font=\scriptsize}]
\filldraw[cefill] (0.000,0.000) -- (1.000,0.000) -- (1.000,1.000) -- (0.000,1.000) -- cycle;
\filldraw[cefill] (1.000,0.000) -- (2.000,0.000) -- (2.000,1.000) -- (1.000,1.000) -- cycle;
\filldraw[cefill] (2.000,0.000) -- (3.000,0.000) -- (3.000,1.000) -- (2.000,1.000) -- cycle;
\filldraw[cefill] (0.000,1.000) -- (1.000,1.000) -- (1.000,2.000) -- (0.000,2.000) -- cycle;
\filldraw[cefill] (1.000,1.000) -- (2.000,1.000) -- (2.000,2.000) -- (1.000,2.000) -- cycle;
\filldraw[cefill] (2.000,1.000) -- (3.000,1.000) -- (3.000,2.000) -- (2.000,2.000) -- cycle;
\filldraw[cefill] (3.000,1.000) -- (4.000,1.000) -- (3.000,2.000) -- cycle;
\filldraw[cefill] (4.000,1.000) -- (4.000,2.000) -- (3.000,2.000) -- cycle;
\filldraw[cefill] (0.000,2.000) -- (1.000,2.000) -- (1.000,3.000) -- (0.000,3.000) -- cycle;
\filldraw[cefill] (1.000,2.000) -- (2.000,2.000) -- (2.000,3.000) -- (1.000,3.000) -- cycle;
\filldraw[cefill] (2.000,2.000) -- (3.000,2.000) -- (3.000,3.000) -- (2.000,3.000) -- cycle;
\filldraw[cefill] (3.000,2.000) -- (4.000,2.000) -- (4.000,3.000) -- cycle;
\filldraw[cefill] (3.000,2.000) -- (4.000,3.000) -- (3.000,3.000) -- cycle;
\draw[seledge] (1.000,2.000) -- (1.000,3.000);
\node[edgenum] at (0.840,2.500) {1};
\draw[seledge] (1.000,2.000) -- (2.000,2.000);
\node[edgenum] at (1.500,2.160) {2};
\draw[seledge] (2.000,1.000) -- (2.000,0.000);
\node[edgenum] at (2.160,0.500) {3};
\draw[seledge] (3.000,2.000) -- (3.000,1.000);
\node[edgenum] at (3.160,1.500) {4};
\draw[seledge] (3.000,3.000) -- (3.000,2.000);
\node[edgenum] at (3.160,2.500) {5};
\draw[seledge] (2.000,2.000) -- (3.000,2.000);
\node[edgenum] at (2.500,2.160) {6};
\draw[seledge] (2.000,2.000) -- (2.000,1.000);
\node[edgenum] at (2.160,1.500) {7};
\draw[seledge] (1.000,2.000) -- (1.000,1.000);
\node[edgenum] at (1.160,1.500) {8};
\draw[seledge] (0.000,2.000) -- (1.000,2.000);
\node[edgenum] at (0.500,2.160) {9};
\draw[seledge] (3.000,2.000) -- (4.000,2.000);
\node[edgenum] at (3.500,2.160) {10};
\draw[seledge] (3.000,2.000) -- (4.000,1.000);
\node[edgenum] at (3.613,1.613) {11};
\draw[seledge] (3.000,2.000) -- (4.000,3.000);
\node[edgenum] at (3.387,2.613) {12};
\draw[hearrow, black!55] (0.220,0.100) -- (0.780,0.100);
\draw[hearrow, black!55] (0.900,0.220) -- (0.900,0.780);
\draw[hearrow, black!55] (0.780,0.900) -- (0.220,0.900);
\draw[hearrow, black!55] (0.100,0.780) -- (0.100,0.220);
\draw[hearrow, black!55] (1.220,0.100) -- (1.780,0.100);
\draw[hearrow, black!55] (1.900,0.220) -- (1.900,0.780);
\draw[hearrow, black!55] (1.780,0.900) -- (1.220,0.900);
\draw[hearrow, black!55] (1.100,0.780) -- (1.100,0.220);
\draw[hearrow, black!55] (2.220,0.100) -- (2.780,0.100);
\draw[hearrow, black!55] (2.900,0.220) -- (2.900,0.780);
\draw[hearrow, black!55] (2.780,0.900) -- (2.220,0.900);
\draw[hearrow, black!55] (2.100,0.780) -- (2.100,0.220);
\draw[hearrow, black!55] (0.220,1.100) -- (0.780,1.100);
\draw[hearrow, black!55] (0.900,1.220) -- (0.900,1.780);
\draw[hearrow, black!55] (0.780,1.900) -- (0.220,1.900);
\draw[hearrow, black!55] (0.100,1.780) -- (0.100,1.220);
\draw[hearrow, black!55] (1.220,1.100) -- (1.780,1.100);
\draw[hearrow, black!55] (1.900,1.220) -- (1.900,1.780);
\draw[hearrow, black!55] (1.780,1.900) -- (1.220,1.900);
\draw[hearrow, black!55] (1.100,1.780) -- (1.100,1.220);
\draw[hearrow, black!55] (2.220,1.100) -- (2.780,1.100);
\draw[hearrow, black!55] (2.900,1.220) -- (2.900,1.780);
\draw[hearrow, black!55] (2.780,1.900) -- (2.220,1.900);
\draw[hearrow, black!55] (2.100,1.780) -- (2.100,1.220);
\draw[hearrow, black!55] (3.175,1.089) -- (3.735,1.089);
\draw[hearrow, black!55] (3.709,1.149) -- (3.149,1.709);
\draw[hearrow, black!55] (3.089,1.735) -- (3.089,1.175);
\draw[hearrow, black!55] (3.911,1.265) -- (3.911,1.825);
\draw[hearrow, black!55] (3.825,1.911) -- (3.265,1.911);
\draw[hearrow, black!55] (3.291,1.851) -- (3.851,1.291);
\draw[hearrow, black!55] (0.220,2.100) -- (0.780,2.100);
\draw[hearrow, black!55] (0.900,2.220) -- (0.900,2.780);
\draw[hearrow, black!55] (0.780,2.900) -- (0.220,2.900);
\draw[hearrow, black!55] (0.100,2.780) -- (0.100,2.220);
\draw[hearrow, black!55] (1.220,2.100) -- (1.780,2.100);
\draw[hearrow, black!55] (1.900,2.220) -- (1.900,2.780);
\draw[hearrow, black!55] (1.780,2.900) -- (1.220,2.900);
\draw[hearrow, black!55] (1.100,2.780) -- (1.100,2.220);
\draw[hearrow, black!55] (2.220,2.100) -- (2.780,2.100);
\draw[hearrow, black!55] (2.900,2.220) -- (2.900,2.780);
\draw[hearrow, black!55] (2.780,2.900) -- (2.220,2.900);
\draw[hearrow, black!55] (2.100,2.780) -- (2.100,2.220);
\draw[hearrow, black!55] (3.265,2.089) -- (3.825,2.089);
\draw[hearrow, black!55] (3.911,2.175) -- (3.911,2.735);
\draw[hearrow, black!55] (3.851,2.709) -- (3.291,2.149);
\draw[hearrow, black!55] (3.149,2.291) -- (3.709,2.851);
\draw[hearrow, black!55] (3.735,2.911) -- (3.175,2.911);
\draw[hearrow, black!55] (3.089,2.825) -- (3.089,2.265);
\node[vreg, minimum size=0.41cm] at (0.000,0.000) {$A$};
\node[vreg, minimum size=0.41cm] at (1.000,0.000) {$B$};
\node[vreg, minimum size=0.41cm] at (2.000,0.000) {$C$};
\node[vreg, minimum size=0.41cm] at (3.000,0.000) {$D$};
\node[vreg, minimum size=0.41cm] at (0.000,1.000) {$E$};
\node[vreg, minimum size=0.41cm] at (1.000,1.000) {$F$};
\node[vreg, minimum size=0.41cm] at (2.000,1.000) {$G$};
\node[vreg, minimum size=0.41cm] at (3.000,1.000) {$H$};
\node[vreg, minimum size=0.41cm] at (4.000,1.000) {$I$};
\node[vreg, minimum size=0.41cm] at (0.000,2.000) {$J$};
\node[vreg, minimum size=0.41cm] at (1.000,2.000) {$K$};
\node[vreg, minimum size=0.41cm] at (2.000,2.000) {$L$};
\node[vreg, minimum size=0.41cm] at (3.000,2.000) {$M$};
\node[vreg, minimum size=0.41cm] at (4.000,2.000) {$N$};
\node[vreg, minimum size=0.41cm] at (0.000,3.000) {$P$};
\node[vreg, minimum size=0.41cm] at (1.000,3.000) {$Q$};
\node[vreg, minimum size=0.41cm] at (2.000,3.000) {$R$};
\node[vreg, minimum size=0.41cm] at (3.000,3.000) {$S$};
\node[vreg, minimum size=0.41cm] at (4.000,3.000) {$T$};

\end{tikzpicture}

  \caption{Entrada: malha com as arestas selecionadas (em vermelho) e a ordem de inserção.}
\end{subfigure}\hfill
\begin{subfigure}[b]{0.55\linewidth}\centering
  \begin{tikzpicture}[scale=1.900, every node/.style={font=\scriptsize}]
\fill[holefill] (0.880,2.120) -- (1.120,2.120) -- (1.120,1.880) -- (0.880,1.880) -- cycle;
\fill[holefill] (1.880,1.880) -- (2.000,2.170) -- (2.120,1.880) -- cycle;
\fill[holefill] (3.076,1.848) -- (2.880,1.880) -- (2.880,2.120) -- (3.076,2.152) -- (3.152,2.076) -- (3.152,1.924) -- cycle;
\filldraw[cefill] (0.000,0.000) -- (1.000,0.000) -- (1.000,1.000) -- (0.000,1.000) -- cycle;
\filldraw[cefill] (1.000,0.000) -- (1.880,0.120) -- (1.830,1.000) -- (1.000,1.000) -- cycle;
\filldraw[cefill] (2.120,0.120) -- (3.000,0.000) -- (2.830,1.000) -- (2.170,1.000) -- cycle;
\filldraw[cefill] (0.000,1.000) -- (1.000,1.000) -- (0.880,1.880) -- (0.120,1.880) -- cycle;
\filldraw[cefill] (1.000,1.000) -- (1.830,1.000) -- (1.880,1.880) -- (1.120,1.880) -- cycle;
\filldraw[cefill] (2.170,1.000) -- (2.830,1.000) -- (2.880,1.880) -- (2.120,1.880) -- cycle;
\filldraw[cefill] (3.120,1.120) -- (3.848,1.076) -- (3.076,1.848) -- cycle;
\filldraw[cefill] (3.924,1.152) -- (3.880,1.880) -- (3.152,1.924) -- cycle;
\filldraw[cefill] (0.120,2.120) -- (0.880,2.120) -- (0.880,2.880) -- (0.000,3.000) -- cycle;
\filldraw[cefill] (1.120,2.120) -- (2.000,2.170) -- (2.000,3.000) -- (1.120,2.880) -- cycle;
\filldraw[cefill] (2.000,2.170) -- (2.880,2.120) -- (2.880,2.880) -- (2.000,3.000) -- cycle;
\filldraw[cefill] (3.152,2.076) -- (3.880,2.120) -- (3.924,2.848) -- cycle;
\filldraw[cefill] (3.076,2.152) -- (3.848,2.924) -- (3.120,2.880) -- cycle;
\filldraw[iefill] (1.120,2.120) -- (1.120,2.880) -- (0.880,2.880) -- (0.880,2.120) -- cycle;
\filldraw[iefill] (1.120,1.880) -- (1.880,1.880) -- (2.000,2.170) -- (1.120,2.120) -- cycle;
\filldraw[iefill] (1.830,1.000) -- (1.880,0.120) -- (2.120,0.120) -- (2.170,1.000) -- cycle;
\filldraw[iefill] (2.880,1.880) -- (2.830,1.000) -- (3.120,1.120) -- (3.076,1.848) -- cycle;
\filldraw[iefill] (2.880,2.880) -- (2.880,2.120) -- (3.076,2.152) -- (3.120,2.880) -- cycle;
\filldraw[iefill] (2.120,1.880) -- (2.880,1.880) -- (2.880,2.120) -- (2.000,2.170) -- cycle;
\filldraw[iefill] (1.880,1.880) -- (1.830,1.000) -- (2.170,1.000) -- (2.120,1.880) -- cycle;
\filldraw[iefill] (0.880,1.880) -- (0.947,1.046) -- (1.050,1.049) -- (1.120,1.880) -- cycle;
\filldraw[iefill] (0.120,1.880) -- (0.880,1.880) -- (0.880,2.120) -- (0.120,2.120) -- cycle;
\filldraw[iefill] (3.152,1.924) -- (3.880,1.880) -- (3.880,2.120) -- (3.152,2.076) -- cycle;
\filldraw[iefill] (3.076,1.848) -- (3.848,1.076) -- (3.924,1.152) -- (3.152,1.924) -- cycle;
\filldraw[iefill] (3.152,2.076) -- (3.924,2.848) -- (3.848,2.924) -- (3.076,2.152) -- cycle;
\draw[red, line width=1.1pt] (0.947,1.046) -- (1.050,1.049);
\node[vreg, dot] at (0.000,0.000) {};
\node[vreg, dot] at (1.000,0.000) {};
\node[vie, dot] at (2.120,0.120) {};
\node[vreg, dot] at (3.000,0.000) {};
\node[vreg, dot] at (0.000,1.000) {};
\node[vcol, dot] at (1.000,1.000) {};
\node[vie, dot] at (1.830,1.000) {};
\node[vie, dot] at (3.120,1.120) {};
\node[vie, dot] at (3.924,1.152) {};
\node[vie, dot] at (0.120,1.880) {};
\node[vie, dot] at (0.880,2.120) {};
\node[vie, dot] at (2.000,2.170) {};
\node[vie, dot] at (2.880,2.120) {};
\node[vie, dot] at (3.880,2.120) {};
\node[vreg, dot] at (0.000,3.000) {};
\node[vie, dot] at (0.880,2.880) {};
\node[vreg, dot] at (2.000,3.000) {};
\node[vie, dot] at (2.880,2.880) {};
\node[vie, dot] at (3.848,2.924) {};
\node[vie, dot] at (1.120,2.880) {};
\node[vie, dot] at (1.120,2.120) {};
\node[vie, dot] at (1.880,0.120) {};
\node[vie, dot] at (2.830,1.000) {};
\node[vie, dot] at (3.120,2.880) {};
\node[vie, dot] at (3.076,1.848) {};
\node[vie, dot] at (1.880,1.880) {};
\node[vie, dot] at (2.880,1.880) {};
\node[vie, dot] at (2.120,1.880) {};
\node[vie, dot] at (2.170,1.000) {};
\node[vie, dot] at (1.120,1.880) {};
\node[vie, dot] at (0.120,2.120) {};
\node[vie, dot] at (0.880,1.880) {};
\node[vie, dot] at (3.152,2.076) {};
\node[vie, dot] at (3.880,1.880) {};
\node[vie, dot] at (3.152,1.924) {};
\node[vie, dot] at (3.848,1.076) {};
\node[vie, dot] at (3.076,2.152) {};
\node[vie, dot] at (3.924,2.848) {};

\end{tikzpicture}

  \caption{Saída produzida pelo código: 12 elementos de interface (laranja), 19 vértices novos e três buracos
  (hachurados).}
\end{subfigure}
\caption{O exemplo da apresentação de projeto~\cite{slides}, executado pelo código. Na saída, cada cópia de um vértice
foi deslocada um pouco na direção dos elementos que a usam, para que vértices coincidentes fiquem visíveis (``vista
explodida''). A malha é a do arquivo \code{malha2Hyb.txt}.}
\label{fig:exemplo-intro}
\end{figure}

\subsection{Proposta e contribuições deste relatório}

O algoritmo estudado aqui~\cite{slides} parte de uma observação simples, que dá nome a um dos slides de projeto: embora
o problema seja \emph{abrir arestas e criar elementos}, as decisões difíceis acontecem nos \emph{vértices}. Por isso,
a inserção de cada aresta é decomposta em duas aplicações de operadores locais, uma em cada extremidade, e o operador
é escolhido olhando apenas para o estado daquele vértice. Há quatro operadores (\opname{Canônico},
\opname{Split Element}, \opname{Expand Edge} e \opname{Open Hole}), e a ordem das arestas é livre.

Este relatório tem quatro objetivos:
\begin{itemize}
\item explicar o problema e a estrutura de dados IBHM de forma acessível, inclusive o que acontece com seus vetores a cada inserção (\cref{sec:intro,sec:estruturas});
\item descrever a arquitetura do algoritmo, cada operador e o algoritmo de travessia exatamente como estão
implementados na classe \code{IBHMInterfaceElementBuilder}, a base a ser evoluída (\cref{sec:arquitetura,sec:operadores,sec:travessia});
\item verificar experimentalmente o comportamento do código (\cref{sec:verificacao});
\item confrontar a documentação existente (apresentação, comentários do código e artigo da IBHM) com a implementação,
apontando imprecisões e defeitos (\cref{sec:documentacao}).
\end{itemize}
A \cref{sec:conclusao} reúne próximos passos, inclusive observações para a extensão 3D.

\section{Estruturas de dados: da CHE à IBHM}
\label{sec:estruturas}

O algoritmo opera inteiramente sobre a topologia da malha. Esta seção apresenta o vocabulário necessário: o conceito
de \emph{half-edge}, a estrutura compacta CHE, sua extensão para malhas híbridas (IBHM) e como as duas estão
implementadas no repositório.

\subsection{Half-edges}

Em uma malha 2D orientada, cada elemento é um polígono cujos vértices são listados no sentido anti-horário. Cada lado
do elemento, visto com a orientação induzida pelo elemento, é uma \emph{half-edge}~\cite{mantyla}. Uma aresta interna
é, portanto, representada por duas half-edges com sentidos opostos, uma em cada elemento vizinho; uma aresta de borda
tem apenas uma. Cada half-edge $h$ guarda o vértice de onde parte, $\V[h]$. A partir dela definem-se três consultas
(\cref{fig:halfedge}):
\begin{itemize}
\item $\nxt(h)$ e $\prv(h)$: a half-edge seguinte e a anterior no mesmo elemento;
\item $\Op[h]$: a half-edge \emph{oposta}, isto é, a mesma aresta vista pelo elemento vizinho (ou \BORDER{} se $h$
está na borda).
\end{itemize}
Note que $h$ parte de $\V[h]$ e chega em $\V[\nxt(h)]$, e que $\Op[h]$ vai de $\V[\nxt(h)]$ para $\V[h]$.

\begin{figure}[ht]
\centering
\begin{subfigure}[b]{0.44\linewidth}\centering
\begin{tikzpicture}[scale=1.25, every node/.style={font=\small}]
  \filldraw[cefill] (0,0) -- (2,0) -- (2,2) -- (0,2) -- cycle;
  \filldraw[cefill, fill=ceblue!60] (2,0) -- (4,0) -- (4,2) -- (2,2) -- cycle;
  \node[elemlabel] at (0.8,1.05) {$c_0$}; \node[elemlabel] at (3.3,1.05) {$c_1$};
  \draw[hearrow, red!70!black, line width=0.9pt] (1.82,0.35) -- (1.82,1.65) node[midway, left=1pt] {$h$};
  \draw[hearrow, violet] (1.65,1.82) -- (0.35,1.82) node[midway, below=1pt] {$\nxt(h)$};
  \draw[hearrow, green!45!black] (0.35,0.18) -- (1.65,0.18) node[midway, above=1pt] {$\prv(h)$};
  \draw[hearrow, blue!70!black, line width=0.9pt] (2.18,1.65) -- (2.18,0.35) node[midway, right=1pt] {$\Op[h]$};
  \node[vreg, minimum size=0.36cm, label={[font=\scriptsize]below right:{$v=\V[h]$}}] at (2,0) {};
  \node[vreg, minimum size=0.36cm] at (2,2) {}; \node[vreg, dot] at (0,0) {}; \node[vreg, dot] at (0,2) {};
  \node[vreg, dot] at (4,0) {}; \node[vreg, dot] at (4,2) {};
\end{tikzpicture}
\caption{Consultas básicas.}
\label{fig:halfedge}
\end{subfigure}\hfill
\begin{subfigure}[b]{0.52\linewidth}\centering
\begin{tikzpicture}[scale=1.2, every node/.style={font=\small}]
  \foreach \i in {0,...,5} { \coordinate (p\i) at ({60*\i}:2.0); }
  \coordinate (v) at (0,0);
  \foreach \i/\j in {0/1,1/2,2/3,3/4,4/5,5/0} { \filldraw[cefill] (v) -- (p\i) -- (p\j) -- cycle; }
  \filldraw[cefill, fill=ceblue!55] (v) -- (p1) -- (p2) -- cycle;
  \node[elemlabel, font=\scriptsize\sffamily\bfseries] at (90:0.62) {$T_1$};
  \node[elemlabel, font=\scriptsize\sffamily\bfseries] at (150:0.62) {$T_2$};
  \node[elemlabel, font=\scriptsize\sffamily\bfseries] at (30:0.62) {$T_0$};
  % h e prev(h) em T1
  \draw[hearrow, red!70!black, line width=0.9pt] ($(v)!0.2!(p1)+(150:0.1)$) -- ($(v)!0.8!(p1)+(150:0.1)$);
  \node[text=red!70!black] at (73:1.05) {$h$};
  \draw[hearrow, violet] ($(p2)!0.2!(v)+(30:0.1)$) -- ($(p2)!0.8!(v)+(30:0.1)$);
  \node[text=violet, font=\scriptsize] at (98:1.55) {$\prv(h)$};
  % O[prev(h)] em T2
  \draw[hearrow, green!45!black, line width=0.9pt] ($(v)!0.2!(p2)+(210:0.1)$) -- ($(v)!0.8!(p2)+(210:0.1)$);
  \node[text=green!45!black, font=\scriptsize] at (142:1.3) {$\Op[\prv(h)]$};
  % O[h] e next(O[h]) em T0
  \draw[hearrow, blue!70!black] ($(p1)!0.2!(v)+(-30:0.1)$) -- ($(p1)!0.8!(v)+(-30:0.1)$);
  \node[text=blue!70!black, font=\scriptsize] at (44:1.5) {$\Op[h]$};
  \draw[hearrow, orange!80!black, line width=0.9pt] ($(v)!0.2!(p0)+(90:0.1)$) -- ($(v)!0.8!(p0)+(90:0.1)$);
  \node[text=orange!80!black, font=\scriptsize] at (13:1.25) {$\nxt(\Op[h])$};
  \node[vreg, minimum size=0.36cm] at (v) {$v$};
  \foreach \i in {0,...,5} { \node[vreg, dot] at (p\i) {}; }
  \draw[-{Stealth[length=5pt]}, green!45!black, line width=0.8pt] (68:2.3) arc (68:118:2.3)
    node[pos=0.55, above=1pt, font=\scriptsize] {anti-horário};
  \draw[-{Stealth[length=5pt]}, orange!80!black, line width=0.8pt] (52:2.3) arc (52:8:2.3)
    node[midway, right=1pt, font=\scriptsize] {horário};
\end{tikzpicture}
\caption{Girando em torno de um vértice $v$.}
\label{fig:rotacao}
\end{subfigure}
\caption{(a) Uma half-edge $h$ do elemento $c_0$ e suas consultas. (b) As duas maneiras de girar em torno de um vértice
usando apenas $\nxt$, $\prv$ e $\Op$: a regra $h\leftarrow\Op[\prv(h)]$ passa ao elemento vizinho no sentido
anti-horário (usada por \code{reindexElementsCCW}), e a regra $h\leftarrow\nxt(\Op[h])$ passa ao vizinho no sentido
horário (usada por \code{isBorder} e pelo segundo laço de \code{openHole}). Em ambos os casos a nova half-edge também
parte de $v$.}
\end{figure}

\subsection{CHE: half-edges implícitas}

A \emph{Compact Half-Edge} (CHE)~\cite{che} dispensa objetos ``half-edge''. Se todos os elementos têm $k$ vértices,
basta guardar a lista de elementos como um vetor de inteiros $\V[\,]$ com $k$ entradas por elemento: a posição $h$ desse
vetor \emph{é} a half-edge, e o valor armazenado é o vértice de origem. As demais consultas saem por aritmética:
\[
  \mathrm{elem}(h)=\lfloor h/k\rfloor,\qquad
  \nxt(h)=k\,\mathrm{elem}(h)+(h+1)\bmod k,\qquad
  \prv(h)=k\,\mathrm{elem}(h)+(h+k-1)\bmod k .
\]
O único vetor adicional é o das opostas, $\Op[\,]$, com uma entrada por half-edge. Toda a adjacência da malha cabe em
dois vetores de inteiros do mesmo tamanho.

\subsection{IBHM: malhas híbridas}

A IBHM (\emph{Index-Based data structure for Hybrid Meshes})~\cite{ibhm} estende a CHE (e a CHF~\cite{chf}, em 3D) para
malhas com elementos de tipos diferentes. Em 2D, a mudança essencial é um vetor de \emph{offsets} $Of[\,]$ que diz onde
cada elemento começa em $\V[\,]$. Assim, o elemento $c$ ocupa as posições $Of[c],\dots,Of[c]+n_v(c)-1$, e:
\begin{itemize}
\item $\mathrm{elem}(h)$ é obtido por \emph{busca binária} em $Of[\,]$, em $O(\log n_c)$ (em vez de uma divisão);
\item $\nxt(h)=Of[c]+\big((h-Of[c]+1)\bmod n_v(c)\big)$ e analogamente para $\prv$;
\item $\Op[\,]$ é idêntico ao da CHE;
\item um vetor extra $VH[\,]$ associa a cada vértice uma half-edge que parte dele (a de borda, se o vértice estiver na
borda), o que permite iniciar consultas de vizinhança e testar borda em $O(1)$.
\end{itemize}
A \cref{fig:ibhm} mostra os vetores para a malha do arquivo \code{meshIBHM.txt}, que reproduz (com índices a partir de
zero) o exemplo da Figura~4 do artigo: um quadrado com um furo quadrado, formado por três quadriláteros e três
triângulos.

\begin{figure}[ht]
\centering
\begin{tikzpicture}[scale=2.500, every node/.style={font=\scriptsize}]
\fill[holefill] (2.000,2.000) -- (2.000,1.000) -- (1.000,1.000) -- (1.000,2.000) -- cycle;
\filldraw[cefill] (2.000,2.000) -- (3.000,3.000) -- (1.500,3.000) -- cycle;
\filldraw[cefill] (3.000,3.000) -- (2.000,2.000) -- (2.000,1.000) -- (3.000,0.000) -- cycle;
\filldraw[cefill] (3.000,0.000) -- (2.000,1.000) -- (1.000,1.000) -- (0.000,0.000) -- cycle;
\filldraw[cefill] (0.000,0.000) -- (1.000,1.000) -- (1.000,2.000) -- (0.000,3.000) -- cycle;
\filldraw[cefill] (0.000,3.000) -- (1.000,2.000) -- (1.500,3.000) -- cycle;
\filldraw[cefill] (1.000,2.000) -- (2.000,2.000) -- (1.500,3.000) -- cycle;
\node[elemlabel] at (2.167,2.667) {$c_{0}$};
\node[elemlabel] at (2.500,1.500) {$c_{1}$};
\node[elemlabel] at (1.500,0.500) {$c_{2}$};
\node[elemlabel] at (0.500,1.500) {$c_{3}$};
\node[elemlabel] at (0.833,2.667) {$c_{4}$};
\node[elemlabel] at (1.500,2.333) {$c_{5}$};
\draw[hearrow, black!70] (2.134,2.308) -- (2.634,2.808);
\node[helabel, text=black!70] at (2.279,2.610) {$0$};
\draw[hearrow, black!70] (2.593,2.874) -- (1.843,2.874);
\node[helabel, text=black!70] at (2.190,2.760) {$1$};
\draw[hearrow, black!70] (1.746,2.798) -- (1.996,2.298);
\node[helabel, text=black!70] at (1.979,2.592) {$2$};
\draw[hearrow, black!70] (2.750,2.620) -- (2.250,2.120);
\node[helabel, text=black!70] at (2.500,2.253) {$3$};
\draw[hearrow, black!70] (2.130,1.750) -- (2.130,1.250);
\node[helabel, text=black!70] at (2.247,1.500) {$4$};
\draw[hearrow, black!70] (2.250,0.880) -- (2.750,0.380);
\node[helabel, text=black!70] at (2.500,0.747) {$5$};
\draw[hearrow, black!70] (2.870,0.750) -- (2.870,2.250);
\node[helabel, text=black!70] at (2.753,1.500) {$6$};
\draw[hearrow, black!70] (2.620,0.250) -- (2.120,0.750);
\node[helabel, text=black!70] at (2.253,0.500) {$7$};
\draw[hearrow, black!70] (1.750,0.870) -- (1.250,0.870);
\node[helabel, text=black!70] at (1.500,0.753) {$8$};
\draw[hearrow, black!70] (0.880,0.750) -- (0.380,0.250);
\node[helabel, text=black!70] at (0.747,0.500) {$9$};
\draw[hearrow, black!70] (0.750,0.130) -- (2.250,0.130);
\node[helabel, text=black!70] at (1.500,0.247) {$10$};
\draw[hearrow, black!70] (0.250,0.380) -- (0.750,0.880);
\node[helabel, text=black!70] at (0.500,0.747) {$11$};
\draw[hearrow, black!70] (0.870,1.250) -- (0.870,1.750);
\node[helabel, text=black!70] at (0.753,1.500) {$12$};
\draw[hearrow, black!70] (0.750,2.120) -- (0.250,2.620);
\node[helabel, text=black!70] at (0.500,2.253) {$13$};
\draw[hearrow, black!70] (0.130,2.250) -- (0.130,0.750);
\node[helabel, text=black!70] at (0.247,1.500) {$14$};
\draw[hearrow, black!70] (0.366,2.808) -- (0.866,2.308);
\node[helabel, text=black!70] at (0.721,2.610) {$15$};
\draw[hearrow, black!70] (1.004,2.298) -- (1.254,2.798);
\node[helabel, text=black!70] at (1.021,2.592) {$16$};
\draw[hearrow, black!70] (1.157,2.874) -- (0.407,2.874);
\node[helabel, text=black!70] at (0.810,2.760) {$17$};
\draw[hearrow, black!70] (1.300,2.070) -- (1.700,2.070);
\node[helabel, text=black!70] at (1.500,2.161) {$18$};
\draw[hearrow, black!70] (1.792,2.261) -- (1.592,2.661);
\node[helabel, text=black!70] at (1.616,2.411) {$19$};
\draw[hearrow, black!70] (1.408,2.661) -- (1.208,2.261);
\node[helabel, text=black!70] at (1.384,2.411) {$20$};
\node[vreg, minimum size=0.29cm] at (3.000,3.000) {$0$};
\node[vreg, minimum size=0.29cm] at (0.000,3.000) {$1$};
\node[vreg, minimum size=0.29cm] at (0.000,0.000) {$2$};
\node[vreg, minimum size=0.29cm] at (3.000,0.000) {$3$};
\node[vreg, minimum size=0.29cm] at (2.000,2.000) {$4$};
\node[vreg, minimum size=0.29cm] at (1.000,2.000) {$5$};
\node[vreg, minimum size=0.29cm] at (1.000,1.000) {$6$};
\node[vreg, minimum size=0.29cm] at (2.000,1.000) {$7$};
\node[vreg, minimum size=0.29cm] at (1.500,3.000) {$8$};

\end{tikzpicture}

\par\medskip
{\footnotesize\setlength{\tabcolsep}{2.2pt}\begin{tabular}{lccccccccccccccccccccc}
\textit{he} & 0 & 1 & 2 & 3 & 4 & 5 & 6 & 7 & 8 & 9 & 10 & 11 & 12 & 13 & 14 & 15 & 16 & 17 & 18 & 19 & 20 \\ \hline
$V[\,]$ & 4 & 0 & 8 & 0 & 4 & 7 & 3 & 3 & 7 & 6 & 2 & 2 & 6 & 5 & 1 & 1 & 5 & 8 & 5 & 4 & 8 \\
$O[\,]$ & 3 & -- & 19 & 0 & -- & 7 & -- & 5 & -- & 11 & -- & 9 & -- & 15 & -- & 13 & 20 & -- & -- & 2 & 16 \\
\end{tabular}
\par\medskip $Of[\,] = [0,\,3,\,7,\,11,\,15,\,18,\,21]$
}
\par\smallskip
{\small $\V[h]$: vértice de origem de $h$;\quad $\Op[h]$: oposta de $h$ (``--'' = \BORDER);\quad $Of[c]$: primeira
half-edge do elemento $c$.}
\caption{A malha de \code{meshIBHM.txt} com o índice de cada half-edge e os vetores da IBHM tal como construídos pela
classe \code{IBHM}. O último elemento ($c_5$) começa num vértice diferente do usado no artigo, por isso suas três
half-edges têm outra numeração; o restante coincide com a Figura~4 de~\cite{ibhm} subtraindo-se 1 dos índices. O furo
central é delimitado por half-edges de borda.}
\label{fig:ibhm}
\end{figure}

\subsection{A classe \code{IBHM} do repositório}
\label{sec:impl-estruturas}

A classe \code{IBHM} (\code{core/TopologyDataStructures/IBHM.\{h,cpp\}}) segue as ideias acima, com algumas escolhas
próprias que importam para o algoritmo de inserção (\cref{tab:estruturas}). As coordenadas são guardadas em
\code{float}, e as malhas são lidas no formato de \code{malha2Hyb.txt}: número de vértices e de elementos, coordenadas e,
para cada elemento, o número de vértices seguido dos índices.

\begin{table}[ht]
\centering\small
\begin{tabularx}{\linewidth}{>{\raggedright}p{3.4cm}>{\raggedright}X>{\raggedright\arraybackslash}X}
\toprule
Aspecto & Artigo da IBHM~\cite{ibhm} & Classe \code{IBHM}\\
\midrule
Índices & a partir de 1; $\Op[h]=0$ na borda & a partir de 0; $\Op[h]=\BORDER=-1$ (como \code{unsigned}: \code{0xFFFFFFFF})\\
Estado extra de $\Op$ & --- & $\COLLAPSED=-2$ marca uma aresta \emph{colapsada} (de comprimento zero) de um
elemento de interface\\
Offsets & $n_c$ entradas & $n_c+1$ entradas (sentinela final), de modo que $n_v(c)=Of[c+1]-Of[c]$\\
$\mathrm{elem}(h)$ & busca binária & busca binária com \code{std::upper\_bound} nas $n_c+1$ primeiras entradas de
$Of[\,]$ (\code{offsetIterator}); $\nxt$ e $\prv$ também fazem essa busca\\
Construção de $\Op$ & mapa indexado pelo par ordenado de vértices & para cada vértice, lista das half-edges que partem
dele; procura-se a half-edge reversa (\code{buildOppositesLinear})\\
Vetor $VH[\,]$ & sim & \textbf{não}; o teste de borda (\code{isBorder}) gira em torno do vértice\\
Crescimento da malha & não previsto (``limitação'' declarada) & \code{reserveSpaceForElements} aumenta os vetores de
elementos (\code{resize}; as posições novas de $\Op[\,]$ começam com \BORDER) e \code{commitElement} valida cada novo
elemento; \code{numberOfHalfEdges} devolve $Of[n_c]$, o número de half-edges válidas (sem o espaço reservado);
\code{reserveSpaceForNodes} reserva capacidade exata para os pontos (\code{reserve}); \code{addPoint} acrescenta um ponto
com as coordenadas dadas num \code{std::vector} (vazio = origem) e \code{duplicatePoint} acrescenta a cópia de um vértice
existente, com as mesmas coordenadas\\
\bottomrule
\end{tabularx}
\caption{Diferenças entre a IBHM do artigo e a classe \code{IBHM} do repositório.}
\label{tab:estruturas}
\end{table}

Dois pontos merecem destaque. Primeiro, o estado \COLLAPSED{} é uma extensão criada para o algoritmo de inserção:
ele permite que um elemento de interface tenha um lado de comprimento zero sem vizinho, e as consultas que giram em torno
de vértices (\code{isBorder}) sabem pular esse lado. Segundo, a combinação ``contar $\to$ reservar $\to$ validar'' é a
resposta da implementação à limitação apontada no artigo para alterações dinâmicas: o algoritmo de inserção calcula
antes quantos vértices e elementos novos serão necessários (\cref{sec:contagem}), reserva o espaço de uma vez e só
então altera a malha, sem realocações durante as operações.

A interface usada pelo construtor de elementos de interface é pequena: \code{heVertexIndex} ($\V$), \code{heOpposite}
($\Op$), \code{heNext}, \code{hePrevious}, \code{heElement}, \code{isBorder}, \code{setElementVertex} (altera $\V[h]$),
\code{setOpposite} (altera $\Op$ nos dois sentidos quando a oposta é uma half-edge válida, ou em um só sentido quando é
\BORDER{} ou \COLLAPSED), \code{duplicatePoint}, \code{numberOfHalfEdges}, \code{nextAvailableHE} e
\code{commitElement}. Nenhuma coordenada entra nas decisões: o algoritmo é puramente topológico, e as cópias criadas por
\code{duplicatePoint} apenas herdam as coordenadas do vértice de origem.

\subsection{Os vetores da IBHM durante uma inserção}
\label{sec:ibhm-insercao}

A \cref{fig:insercao-vetores} mostra, numa malha híbrida pequena, o que acontece com os vetores quando uma aresta é
aberta. A aresta $4$--$7$ é compartilhada pelos triângulos que contêm as half-edges $a=9$ ($4\to7$) e $b=19$ ($7\to4$);
o vértice 4 é interno e o vértice 7 está na borda.
\begin{enumerate}
\item \textbf{Reserva.} Antes de qualquer alteração, \code{reserveSpaceForElements(4, 1)} acrescenta 4 posições a
$\V[\,]$ e a $\Op[\,]$ (estas com \BORDER) e uma entrada a $Of[\,]$, e \code{reserveSpaceForNodes(1)} reserva capacidade
exata para um vértice, sem mudar o tamanho do vetor de coordenadas (a pré-contagem previu um vértice novo, para a
borda).
\item \textbf{Validação antecipada.} \code{nextAvailableHE()} devolve $Of[n_c]=20$, a primeira half-edge livre, e
\code{commitElement(4)} já declara o elemento 6 válido, gravando $Of[7]=24$. Na IBHM isso precisa acontecer \emph{antes}
dos operadores, porque $\nxt$ e $\prv$ das half-edges 20--23 dependem do offset do elemento.
\item \textbf{Operador em 4} (Canônico, interno): $\V[23]=\V[20]=4$ e $\Op[23]=\COLLAPSED$ (marcado como \textsc{c}).
\item \textbf{Operador em 7} (Canônico, borda): cria o vértice 9, cópia de 7 (\code{duplicatePoint(7)}, mesmas
coordenadas); $\V[21]=9$; a rotação a partir de $b$ troca $\V[19]$ de 7 para 9 e para logo na borda; $\V[22]=7$ e
$\Op[21]=\BORDER$.
\item \textbf{Arestas do meio}: $\Op[20]\leftrightarrow19$ e $\Op[22]\leftrightarrow9$. As antigas opostas $\Op[9]=19$ e
$\Op[19]=9$ são sobrescritas.
\end{enumerate}
Ao todo, uma inserção altera poucas posições dos vetores antigos (em vermelho na figura) e preenche as quatro posições
novas (em laranja). As posições reservadas de $\Op[\,]$ começam com \BORDER: até a v0.3 começavam com 0, um índice
válido de half-edge, e um índice do espaço reservado passado como entrada era aceito como aresta (defeito I3, corrigido;
\cref{sec:defeitos}). No fluxo normal todas são sobrescritas pelos operadores e pela ligação das arestas do meio.

\begin{figure}[ht]
\centering
\begin{subfigure}[b]{0.48\linewidth}\centering\begin{tikzpicture}[scale=2.100, every node/.style={font=\scriptsize}]
\filldraw[cefill] (0.000,0.000) -- (1.000,0.000) -- (1.000,1.000) -- (0.000,1.000) -- cycle;
\filldraw[cefill] (1.000,0.000) -- (2.000,0.000) -- (2.000,1.000) -- (1.000,1.000) -- cycle;
\filldraw[cefill] (0.000,1.000) -- (1.000,1.000) -- (1.000,2.000) -- cycle;
\filldraw[cefill] (0.000,1.000) -- (1.000,2.000) -- (0.000,2.000) -- cycle;
\filldraw[cefill] (1.000,1.000) -- (2.000,1.000) -- (2.000,2.000) -- cycle;
\filldraw[cefill] (1.000,1.000) -- (2.000,2.000) -- (1.000,2.000) -- cycle;
\draw[seledge] (1.000,1.000) -- (1.000,2.000);
\draw[hearrow, black!70] (0.250,0.110) -- (0.750,0.110);
\node[helabel, text=black!70] at (0.500,0.198) {$0$};
\draw[hearrow, black!70] (0.890,0.250) -- (0.890,0.750);
\node[helabel, text=black!70] at (0.802,0.500) {$1$};
\draw[hearrow, black!70] (0.750,0.890) -- (0.250,0.890);
\node[helabel, text=black!70] at (0.500,0.802) {$2$};
\draw[hearrow, black!70] (0.110,0.750) -- (0.110,0.250);
\node[helabel, text=black!70] at (0.198,0.500) {$3$};
\draw[hearrow, black!70] (1.250,0.110) -- (1.750,0.110);
\node[helabel, text=black!70] at (1.500,0.198) {$4$};
\draw[hearrow, black!70] (1.890,0.250) -- (1.890,0.750);
\node[helabel, text=black!70] at (1.802,0.500) {$5$};
\draw[hearrow, black!70] (1.750,0.890) -- (1.250,0.890);
\node[helabel, text=black!70] at (1.500,0.802) {$6$};
\draw[hearrow, black!70] (1.110,0.750) -- (1.110,0.250);
\node[helabel, text=black!70] at (1.198,0.500) {$7$};
\draw[hearrow, black!70] (0.299,1.098) -- (0.799,1.098);
\node[helabel, text=black!70] at (0.589,1.177) {$8$};
\draw[hearrow, red!70!black] (0.902,1.201) -- (0.902,1.701);
\node[helabel, text=red!70!black] at (0.823,1.411) {$9$};
\draw[hearrow, black!70] (0.825,1.675) -- (0.325,1.175);
\node[helabel, text=black!70] at (0.635,1.365) {$10$};
\draw[hearrow, black!70] (0.175,1.325) -- (0.675,1.825);
\node[helabel, text=black!70] at (0.365,1.635) {$11$};
\draw[hearrow, black!70] (0.701,1.902) -- (0.201,1.902);
\node[helabel, text=black!70] at (0.411,1.823) {$12$};
\draw[hearrow, black!70] (0.098,1.799) -- (0.098,1.299);
\node[helabel, text=black!70] at (0.177,1.589) {$13$};
\draw[hearrow, black!70] (1.299,1.098) -- (1.799,1.098);
\node[helabel, text=black!70] at (1.589,1.177) {$14$};
\draw[hearrow, black!70] (1.902,1.201) -- (1.902,1.701);
\node[helabel, text=black!70] at (1.823,1.411) {$15$};
\draw[hearrow, black!70] (1.825,1.675) -- (1.325,1.175);
\node[helabel, text=black!70] at (1.635,1.365) {$16$};
\draw[hearrow, black!70] (1.175,1.325) -- (1.675,1.825);
\node[helabel, text=black!70] at (1.365,1.635) {$17$};
\draw[hearrow, black!70] (1.701,1.902) -- (1.201,1.902);
\node[helabel, text=black!70] at (1.411,1.823) {$18$};
\draw[hearrow, red!70!black] (1.098,1.799) -- (1.098,1.299);
\node[helabel, text=red!70!black] at (1.177,1.589) {$19$};
\node[vreg, minimum size=0.26cm] at (0.000,0.000) {$0$};
\node[vreg, minimum size=0.26cm] at (1.000,0.000) {$1$};
\node[vreg, minimum size=0.26cm] at (2.000,0.000) {$2$};
\node[vreg, minimum size=0.26cm] at (0.000,1.000) {$3$};
\node[vreg, minimum size=0.26cm] at (1.000,1.000) {$4$};
\node[vreg, minimum size=0.26cm] at (2.000,1.000) {$5$};
\node[vreg, minimum size=0.26cm] at (0.000,2.000) {$6$};
\node[vreg, minimum size=0.26cm] at (1.000,2.000) {$7$};
\node[vreg, minimum size=0.26cm] at (2.000,2.000) {$8$};

\end{tikzpicture}
\caption{Antes: a aresta $4$--$7$.}\end{subfigure}\hfill
\begin{subfigure}[b]{0.48\linewidth}\centering\begin{tikzpicture}[scale=2.100, every node/.style={font=\scriptsize}]
\filldraw[cefill] (0.000,0.000) -- (1.000,0.000) -- (1.000,1.000) -- (0.000,1.000) -- cycle;
\filldraw[cefill] (1.000,0.000) -- (2.000,0.000) -- (2.000,1.000) -- (1.000,1.000) -- cycle;
\filldraw[cefill] (0.000,1.000) -- (1.000,1.000) -- (0.859,1.859) -- cycle;
\filldraw[cefill] (0.000,1.000) -- (0.859,1.859) -- (0.000,2.000) -- cycle;
\filldraw[cefill] (1.000,1.000) -- (2.000,1.000) -- (2.000,2.000) -- cycle;
\filldraw[cefill] (1.000,1.000) -- (2.000,2.000) -- (1.141,1.859) -- cycle;
\filldraw[iefill] (1.052,1.085) -- (1.141,1.859) -- (0.859,1.859) -- (0.920,1.060) -- cycle;
\draw[red, line width=1.1pt] (0.920,1.060) -- (1.052,1.085);
\draw[hearrow, black!70] (0.250,0.110) -- (0.750,0.110);
\node[helabel, text=black!70] at (0.500,0.198) {$0$};
\draw[hearrow, black!70] (0.890,0.250) -- (0.890,0.750);
\node[helabel, text=black!70] at (0.802,0.500) {$1$};
\draw[hearrow, black!70] (0.750,0.890) -- (0.250,0.890);
\node[helabel, text=black!70] at (0.500,0.802) {$2$};
\draw[hearrow, black!70] (0.110,0.750) -- (0.110,0.250);
\node[helabel, text=black!70] at (0.198,0.500) {$3$};
\draw[hearrow, black!70] (1.250,0.110) -- (1.750,0.110);
\node[helabel, text=black!70] at (1.500,0.198) {$4$};
\draw[hearrow, black!70] (1.890,0.250) -- (1.890,0.750);
\node[helabel, text=black!70] at (1.802,0.500) {$5$};
\draw[hearrow, black!70] (1.750,0.890) -- (1.250,0.890);
\node[helabel, text=black!70] at (1.500,0.802) {$6$};
\draw[hearrow, black!70] (1.110,0.750) -- (1.110,0.250);
\node[helabel, text=black!70] at (1.198,0.500) {$7$};
\draw[hearrow, black!70] (0.292,1.102) -- (0.792,1.102);
\node[helabel, text=black!70] at (0.576,1.183) {$8$};
\draw[hearrow, black!70] (0.865,1.169) -- (0.794,1.598);
\node[helabel, text=black!70] at (0.750,1.346) {$9$};
\draw[hearrow, black!70] (0.730,1.580) -- (0.300,1.150);
\node[helabel, text=black!70] at (0.583,1.313) {$10$};
\draw[hearrow, black!70] (0.150,1.300) -- (0.580,1.730);
\node[helabel, text=black!70] at (0.313,1.583) {$11$};
\draw[hearrow, black!70] (0.598,1.794) -- (0.169,1.865);
\node[helabel, text=black!70] at (0.346,1.750) {$12$};
\draw[hearrow, black!70] (0.102,1.792) -- (0.102,1.292);
\node[helabel, text=black!70] at (0.183,1.576) {$13$};
\draw[hearrow, black!70] (1.299,1.098) -- (1.799,1.098);
\node[helabel, text=black!70] at (1.589,1.177) {$14$};
\draw[hearrow, black!70] (1.902,1.201) -- (1.902,1.701);
\node[helabel, text=black!70] at (1.823,1.411) {$15$};
\draw[hearrow, black!70] (1.825,1.675) -- (1.325,1.175);
\node[helabel, text=black!70] at (1.635,1.365) {$16$};
\draw[hearrow, black!70] (1.196,1.304) -- (1.696,1.804);
\node[helabel, text=black!70] at (1.403,1.597) {$17$};
\draw[hearrow, black!70] (1.728,1.871) -- (1.299,1.800);
\node[helabel, text=black!70] at (1.467,1.761) {$18$};
\draw[hearrow, black!70] (1.200,1.701) -- (1.129,1.272);
\node[helabel, text=black!70] at (1.239,1.533) {$19$};
\draw[hearrow, orange!50!black] (1.132,1.255) -- (1.182,1.688);
\node[helabel, text=orange!50!black] at (1.337,1.222) {$20$};
\draw[hearrow, orange!50!black] (0.812,1.683) -- (0.847,1.236);
\node[helabel, text=orange!50!black] at (0.629,1.709) {$22$};
\node[vreg, minimum size=0.26cm] at (0.000,0.000) {$0$};
\node[vreg, minimum size=0.26cm] at (1.000,0.000) {$1$};
\node[vreg, minimum size=0.26cm] at (2.000,0.000) {$2$};
\node[vreg, minimum size=0.26cm] at (0.000,1.000) {$3$};
\node[vcol, minimum size=0.26cm] at (1.000,1.000) {$4$};
\node[vreg, minimum size=0.26cm] at (2.000,1.000) {$5$};
\node[vreg, minimum size=0.26cm] at (0.000,2.000) {$6$};
\node[vie, minimum size=0.26cm] at (0.859,1.859) {$7$};
\node[vreg, minimum size=0.26cm] at (2.000,2.000) {$8$};
\node[vie, minimum size=0.26cm] at (1.141,1.859) {$9$};

\end{tikzpicture}
\caption{Depois: elemento $6=[4,9,7,4]$.}\end{subfigure}
\par\medskip
{\scriptsize\setlength{\tabcolsep}{2.1pt}\begin{tabular}{lcccccccccccccccccccccccc}
\textit{he} & 0 & 1 & 2 & 3 & 4 & 5 & 6 & 7 & 8 & 9 & 10 & 11 & 12 & 13 & 14 & 15 & 16 & 17 & 18 & 19 & 20 & 21 & 22 & 23 \\ \hline
$\V$ antes & 0 & 1 & 4 & 3 & 1 & 2 & 5 & 4 & 3 & 4 & 7 & 3 & 7 & 6 & 4 & 5 & 8 & 4 & 8 & 7 &  &  &  &  \\
$\V$ depois & 0 & 1 & 4 & 3 & 1 & 2 & 5 & 4 & 3 & 4 & 7 & 3 & 7 & 6 & 4 & 5 & 8 & 4 & 8 & \textcolor{red!75!black}{\textbf{9}} & \textcolor{orange!50!black}{\textbf{4}} & \textcolor{orange!50!black}{\textbf{9}} & \textcolor{orange!50!black}{\textbf{7}} & \textcolor{orange!50!black}{\textbf{4}} \\ \hline
$\Op$ antes & -- & 7 & 8 & -- & -- & -- & 14 & 1 & 2 & 19 & 11 & 10 & -- & -- & 6 & -- & 17 & 16 & -- & 9 &  &  &  &  \\
$\Op$ depois & -- & 7 & 8 & -- & -- & -- & 14 & 1 & 2 & \textcolor{red!75!black}{\textbf{22}} & 11 & 10 & -- & -- & 6 & -- & 17 & 16 & -- & \textcolor{red!75!black}{\textbf{20}} & \textcolor{orange!50!black}{\textbf{19}} & \textcolor{orange!50!black}{\textbf{--}} & \textcolor{orange!50!black}{\textbf{9}} & \textcolor{orange!50!black}{\textbf{\textsc{c}}} \\
\end{tabular}
\par\smallskip $Of$ antes $= [0,\,4,\,8,\,11,\,14,\,17,\,20]$\qquad $Of$ depois $= [0,\,4,\,8,\,11,\,14,\,17,\,20,\,24]$
}
\caption{Uma inserção vista pelos vetores da IBHM. Na tabela, vermelho indica entradas antigas que mudaram e laranja as
posições do novo elemento (half-edges 20--23); ``--'' é \BORDER{} e \textsc{c} é \COLLAPSED. O novo elemento é
$[F,G,H,E]=[4,9,7,4]$: colapsado em 4 (interno) e aberto em 7 (borda), onde surgiu o vértice 9.}
\label{fig:insercao-vetores}
\end{figure}

\section{Arquitetura do algoritmo}
\label{sec:arquitetura}

\subsection{Princípios de projeto}

A apresentação de projeto~\cite{slides} lista como objetivos ``um conjunto simples e pequeno de operadores'', resolver
a implementação ``no nível dos operadores'' e um algoritmo que percorra as arestas de entrada aplicando-os. A
implementação materializa esses objetivos nas seguintes características:

\begin{enumerate}[label=\textbf{P\arabic*.}, leftmargin=2.2em]
\item \textbf{Baseado em operadores.} Inserir uma aresta $AB$ é sempre: criar o elemento, aplicar \emph{um} operador em
$A$, aplicar \emph{um} operador em $B$ e ligar as duas arestas ``do meio'' (\cref{sec:gabarito}). Há apenas quatro
operadores.
\item \textbf{Decisão centrada no vértice.} O operador aplicado em um vértice depende só do estado desse vértice
(pertence a algum elemento de interface? tem uma aresta colapsada? essa aresta compartilha um elemento com a nova
aresta?). Não há análise de combinações de arestas nem dependência da ordem de chegada.
\item \textbf{Localidade.} Cada operador altera somente a estrela de um vértice; o custo é $O(1)$ ou proporcional à
valência.
\item \textbf{Puramente topológico.} Nenhuma coordenada, ângulo ou produto vetorial é usado nas decisões: só $\V$,
$\Op$, $\nxt$ e $\prv$. As cópias de um vértice apenas herdam as coordenadas do original. A apresentação justifica
explicitamente essa escolha pela futura extensão a 3D.
\item \textbf{Duplicação adiada.} O primeiro contato com um vértice interno \emph{não} cria vértice novo: o elemento de
interface ganha uma aresta \emph{colapsada} (de comprimento zero) naquele vértice. A cópia só é criada quando, e se,
uma segunda aresta chegar. É isso que faz uma ponta de trinca ficar fechada e que torna a contagem de vértices exata.
\item \textbf{Incremental.} O construtor guarda o estado necessário entre chamadas; novas arestas podem ser inseridas a
qualquer momento sobre uma malha que já tem elementos de interface (inserção adaptativa).
\item \textbf{Memória previsível.} O número de vértices novos é calculado antes de alterar a malha, e todo o espaço é
reservado de uma vez (\cref{sec:contagem}), com tamanho exato.
\end{enumerate}
Cada um desses princípios, quando pode ser verificado pelo resultado, tem um teste automatizado
(\cref{tab:principios}, \cref{sec:testes}).

\subsection{Componentes}

A classe \code{IBHMInterfaceElementBuilder} (\code{core/InterfaceElementsBuilder}) recebe um ponteiro para a malha
\code{IBHM} e mantém três estruturas auxiliares, que formam uma ``segunda camada'' de informação sobre a malha:

\begin{table}[ht]
\centering\small
\begin{tabularx}{\linewidth}{>{\ttfamily\small}p{4.2cm}X}
\toprule
\normalfont Estado & Significado\\
\midrule
\_inInterfaceElement[v] & o vértice $v$ pertence a algum elemento de interface (por vértice).\\
\_isInterfaceElement[e] & o elemento $e$ é um elemento de interface (por elemento).\\
\_collapsed & mapa $v\mapsto c$: vértices que têm uma aresta colapsada e a half-edge colapsada $c$ (que vai de $v$
para $v$).\\
\bottomrule
\end{tabularx}
\caption{Estado mantido pelo construtor entre inserções.}
\label{tab:estado}
\end{table}

Como esse estado não está guardado na malha, deve existir um único construtor por malha durante toda a sessão de
edição; os comentários da classe registram essa restrição, ainda sem um mecanismo que a garanta (há um \code{@todo}).
A classe declara como amiga a classe \code{BuilderInspector} (\code{tests/support}), que dá aos testes acesso de leitura
a esse estado e ao operador escolhido; a biblioteca não a usa. Os métodos se dividem em quatro grupos:

\begin{itemize}
\item \textbf{travessia}: \code{insertInterfaceElements} valida e percorre as arestas de entrada (\cref{sec:travessia});
\item \textbf{decisão}: \code{retrieveOperator} escolhe o operador, com o auxílio de \code{retrieveAvailableVertex} para
distinguir \opname{Split Element} de \opname{Expand Edge};
\item \textbf{operadores}: \code{canonical}, \code{splitElement}, \code{expandEdge} e \code{openHole};
\item \textbf{apoio}: \code{reindexElementsCCW} (troca o vértice de um leque de elementos) e
\code{computeNumberOfNewVertices} (pré-contagem).
\end{itemize}

\subsection{O gabarito do elemento de interface}
\label{sec:gabarito}

Todo elemento de interface é criado com o mesmo gabarito (\cref{fig:gabarito}). Seja $a$ a half-edge selecionada, que
vai de $A$ para $B$ no elemento $U$, e $b=\Op[a]$, que vai de $B$ para $A$ no elemento $L$. O novo elemento recebe as
quatro próximas half-edges livres, $h_0,\dots,h_3$, cujos vértices de origem são $F,G,H,E$:
\begin{itemize}
\item $h_0: F\to G$ e $h_2: H\to E$ são as arestas ``do meio'', paralelas a $AB$. Ao final, $h_0$ vira oposta de $b$ e
$h_2$ vira oposta de $a$; $a$ e $b$ deixam de ser opostas entre si.
\item $h_3: E\to F$ é o lado curto em $A$; $E$ e $F$ são cópias (topológicas) de $A$, $E$ do lado de $U$ e $F$ do lado
de $L$.
\item $h_1: G\to H$ é o lado curto em $B$; $G$ é a cópia de $B$ do lado de $L$, e $H$, do lado de $U$.
\end{itemize}
O operador do vértice $A$ é chamado com o par $(a, h_3)$ e o do vértice $B$ com $(b, h_1)$. Em ambos os casos vale a
mesma convenção: \emph{a origem de \code{elementHE} fica do lado do elemento que contém \code{edgeHE}, e
$\nxt(\text{\code{elementHE}})$ fica do outro lado}. Cada operador, portanto, só precisa decidir dois vértices,
$\V[\text{\code{elementHE}}]$ e $\V[\nxt(\text{\code{elementHE}})]$, reindexar o que for preciso na estrela do vértice
e definir $\Op[\text{\code{elementHE}}]$. As arestas do meio são sempre tratadas da mesma forma pela travessia, qualquer
que seja o operador (slide ``What happens in the middle?'').

\begin{figure}[ht]
\centering
\begin{tikzpicture}[scale=1.15, every node/.style={font=\small}]
  \filldraw[cefill] (0,1.6) rectangle (6,3.0); \filldraw[cefill] (0,-1.4) rectangle (6,0);
  \filldraw[iefill] (0,0) rectangle (6,1.6);
  \node[elemlabel] at (3,2.65) {U}; \node[elemlabel] at (3,-1.0) {L};
  \node[font=\footnotesize\bfseries, text=orange!40!black] at (3,0.8) {elemento de interface};
  % half-edges da interface
  \draw[hearrow, line width=0.8pt] (0.6,0.25) -- (5.4,0.25) node[pos=0.12, above=0pt] {$h_0$};
  \draw[hearrow, line width=0.8pt] (5.75,0.3) -- (5.75,1.3) node[midway, left=1pt] {$h_1$};
  \draw[hearrow, line width=0.8pt] (5.4,1.35) -- (0.6,1.35) node[pos=0.12, below=0pt] {$h_2$};
  \draw[hearrow, line width=0.8pt] (0.25,1.3) -- (0.25,0.3) node[midway, right=1pt] {$h_3$};
  % a e b
  \draw[hearrow, red!70!black, line width=0.9pt] (0.6,1.85) -- (5.4,1.85) node[pos=0.12, above=0pt] {$a$};
  \draw[hearrow, blue!70!black, line width=0.9pt] (5.4,-0.25) -- (0.6,-0.25) node[pos=0.12, below=0pt] {$b$};
  % opostas
  \draw[{Stealth[length=3pt]}-{Stealth[length=3pt]}, dashed, black!60] (3,1.42) -- (3,1.78)
    node[right=0pt, font=\scriptsize, pos=0.5] {$\Op$};
  \draw[{Stealth[length=3pt]}-{Stealth[length=3pt]}, dashed, black!60] (3,-0.18) -- (3,0.18)
    node[right=0pt, font=\scriptsize, pos=0.5] {$\Op$};
  % vértices
  \node[vie, minimum size=0.42cm] at (0,0) {$F$}; \node[vie, minimum size=0.42cm] at (6,0) {$G$};
  \node[vie, minimum size=0.42cm] at (6,1.6) {$H$}; \node[vie, minimum size=0.42cm] at (0,1.6) {$E$};
  \node[align=right, font=\scriptsize, anchor=east] at (-0.35,0.8) {cópias de $A$\\ $E$: lado de $U$\\ $F$: lado de $L$};
  \node[align=left, font=\scriptsize, anchor=west] at (6.35,0.8) {cópias de $B$\\ $H$: lado de $U$\\ $G$: lado de $L$};
  \node[font=\scriptsize, text=black!65, anchor=north] at (0,-1.45) {operador em $A$: $(a,h_3)$};
  \node[font=\scriptsize, text=black!65, anchor=north] at (6,-1.45) {operador em $B$: $(b,h_1)$};
\end{tikzpicture}
\caption{Gabarito de todo elemento de interface ($FGHE$, sentido anti-horário), desenhado com espessura exagerada. As
arestas do meio $h_0$ e $h_2$ são ligadas a $b$ e $a$ pela travessia; os lados curtos $h_3$ (em $A$) e $h_1$ (em $B$)
são responsabilidade dos operadores. Quando $A$ e $B$ são internos e ainda não pertencem a elemento de interface, o
elemento resultante é $[A,B,B,A]$, com $h_1$ e $h_3$ colapsadas.}
\label{fig:gabarito}
\end{figure}

\subsection{Estados de um vértice e escolha do operador}
\label{sec:estados}

Do ponto de vista de um vértice original, a sequência de inserções o faz passar por uma pequena máquina de estados
(\cref{fig:estados}). Um vértice interno começa \emph{regular}; a primeira aresta o deixa \emph{colapsado} (um elemento
de interface, nenhuma cópia nova); a segunda o \emph{divide} em duas cópias, que ficam unidas pelos lados curtos dos
dois elementos de interface; a terceira e as seguintes abrem e aumentam um \emph{buraco}. Um vértice de borda é
dividido já na primeira aresta, porque a abertura escapa para o exterior; cada aresta seguinte acrescenta uma cópia pelo
Open Hole, com o buraco aberto para o exterior (\cref{sec:falha-borda}). Os testes automatizados verificam essa máquina
de estados em cada vértice, depois de cada inserção (\cref{sec:testes}).

\begin{figure}[ht]
\centering
\begin{tikzpicture}[node distance=1.0cm and 1.55cm, every node/.style={font=\small}]
  \node[state, fill=ceblue] (R) {\textbf{Regular}\\ \scriptsize 0 EI};
  \node[state, fill=ieorange!60, right=of R] (C) {\textbf{Colapsado}\\ \scriptsize 1 EI, 1 cópia};
  \node[state, fill=ieorange!60, right=of C] (D) {\textbf{Dividido}\\ \scriptsize 2 EI, 2 cópias};
  \node[state, fill=ieorange!60, right=of D] (H) {\textbf{Buraco}\\ \scriptsize $k\ge3$ EI, $k$ cópias};
  \draw[flow] (R) -- node[above, font=\scriptsize\bfseries] {Canônico} (C);
  \draw[flow] (C) -- node[above, font=\scriptsize\bfseries, align=center] {Split /\\Expand} (D);
  \draw[flow] (D) -- node[above, font=\scriptsize\bfseries] {Open Hole} (H);
  \draw[flow] (H) edge[loop below, looseness=5] node[below, font=\scriptsize\bfseries] {Open Hole} (H);
  \node[state, fill=ceblue, below=1.5cm of R] (Rb) {\textbf{Regular}\\ \scriptsize na borda};
  \node[state, fill=ieorange!60, right=of Rb] (Db) {\textbf{Aberto}\\ \scriptsize 1 EI, 2 cópias};
  \node[state, fill=ieorange!60, right=of Db] (Fb) {\textbf{Aberto}\\ \scriptsize $k\ge2$ EI, $k+1$ cópias};
  \draw[flow] (Rb) -- node[above, font=\scriptsize\bfseries, align=center] {Canônico\\(borda)} (Db);
  \draw[flow] (Db) -- node[above, font=\scriptsize\bfseries] {Open Hole} (Fb);
  \draw[flow] (Fb) edge[loop below, looseness=5] node[below, font=\scriptsize\bfseries] {Open Hole} (Fb);
  \node[anchor=west, font=\scriptsize, text=black!60] at ($(R.north west)+(0,0.35)$) {vértice interno};
  \node[anchor=west, font=\scriptsize, text=black!60] at ($(Rb.north west)+(0,0.35)$) {vértice de borda};
\end{tikzpicture}
\caption{Estados de um vértice original ao longo das inserções (EI = elemento de interface). Os rótulos das
transições são os operadores aplicados. Com $k$ arestas selecionadas, um vértice interno termina com $\max(1,k)$ cópias
e um de borda com $k+1$. Os lados curtos dos elementos de interface em volta do vértice são: uma aresta colapsada ($k=1$,
interno), dois lados colados um no outro ($k=2$, interno) ou $k$ lados de borda (buraco com $k\ge3$, interno, ou
abertura para o exterior, borda).}
\label{fig:estados}
\end{figure}

A implementação não guarda esse estado explicitamente; ele é deduzido em \code{retrieveOperator} a partir das três
estruturas da \cref{tab:estado}, olhando para o vértice $v$ de onde parte a half-edge recebida (\cref{fig:decisao}):

\begin{figure}[ht]
\centering
\begin{tikzpicture}[node distance=0.75cm and 1.1cm, every node/.style={font=\small}]
  \node[box] (in) {half-edge $he$ que parte de $v$};
  \node[decision, below=of in] (q1) {$v$ pertence a\\algum EI?};
  \node[opbox, fill=green!12, left=1.4cm of q1] (can) {Canônico\\[-1pt]\scriptsize\normalfont borda? novo vértice\\[-2pt]
    \scriptsize\normalfont : aresta colapsada};
  \node[decision, below=of q1] (q2) {$v$ tem aresta\\colapsada?};
  \node[opbox, fill=red!10, left=1.4cm of q2] (hole) {Open Hole};
  \node[decision, below=of q2] (q3) {hipótese $x_1$\\ou $x_3$ vale?};
  \node[opbox, fill=blue!10, left=1.4cm of q3] (split) {Split Element};
  \node[opbox, fill=violet!10, right=1.4cm of q3] (exp) {Expand Edge};
  \draw[flow] (in) -- (q1);
  \draw[flow] (q1) -- node[above, font=\scriptsize] {não} (can);
  \draw[flow] (q1) -- node[right, font=\scriptsize] {sim} (q2);
  \draw[flow] (q2) -- node[above, font=\scriptsize] {não} (hole);
  \draw[flow] (q2) -- node[right, font=\scriptsize] {sim} (q3);
  \draw[flow] (q3) -- node[above, font=\scriptsize] {sim} (split);
  \draw[flow] (q3) -- node[above, font=\scriptsize] {não} (exp);
  \node[font=\scriptsize\ttfamily, text=black!55, right=0.3cm of q1] {\_inInterfaceElement[v]};
  \node[font=\scriptsize\ttfamily, text=black!55, right=0.3cm of q2] {\_collapsed.find(v)};
  \node[font=\scriptsize\ttfamily, text=black!55, below=0.15cm of q3] {retrieveAvailableVertex};
\end{tikzpicture}
\caption{Árvore de decisão de \code{retrieveOperator}. As hipóteses $x_1$ e $x_3$ verificam se o elemento de interface
que tem a aresta colapsada em $v$ é vizinho de um dos dois elementos que contêm a aresta nova
(\cref{sec:split}).}
\label{fig:decisao}
\end{figure}

\begin{table}[ht]
\centering\small
\setlength{\tabcolsep}{4pt}
\begin{tabularx}{\linewidth}{lXcccc}
\toprule
Operador & Condição em $v=\V[he]$ & novos & $\V[\text{eHE}]$, $\V[\nxt(\text{eHE})]$ & $\Op[\text{eHE}]$ & custo\\
\midrule
Canônico (interno) & $v$ fora de EI, interno & 0 & $v$, $v$ & \COLLAPSED & $O(d\log n_c)$\textsuperscript{*}\\
Canônico (borda) & $v$ fora de EI, na borda & 1 ($v'$) & $v'$, $v$ & \BORDER & $O(d\log n_c)$\\
Split Element ($x_1$) & colapsado, $x_1$ vale & 1 ($n$) & $n$, $v$ & $c$ & $O(\log n_c)$\\
Split Element ($x_3$) & colapsado, $x_3$ vale & 1 ($n$) & $v$, $n$ & $c$ & $O(\log n_c)$\\
Expand Edge & colapsado, nenhuma hipótese & 1 ($n$) & $v$, $n$ & $c$ & $O(d\log n_c)$\\
Open Hole & em EI, sem colapsada & 1 ($n$) & $n$, $v$ & \BORDER & $O(d\log n_c)$\\
\bottomrule
\end{tabularx}
\caption{Resumo dos operadores. ``eHE'' é a half-edge \code{elementHE} do novo elemento ($h_3$ em $A$, $h_1$ em $B$);
$c$ é a half-edge colapsada registrada em \code{\_collapsed[v]}, que deixa de ser colapsada e passa a ser a oposta de
eHE. $d$ é a valência do vértice e $n_c$ o número de elementos: na IBHM cada $\nxt$/$\prv$ faz uma busca binária em
$Of[\,]$. \textsuperscript{*}O custo do Canônico interno vem apenas do teste \code{isBorder}, que gira em torno do
vértice.}
\label{tab:operadores}
\end{table}

\section{Os operadores}
\label{sec:operadores}

Esta seção descreve os quatro operadores tal como implementados. Cada um recebe a half-edge da aresta nova que parte do
vértice $v$ (chamada \code{edgeHE} ou \code{he} no código) e a half-edge do lado curto correspondente do novo elemento
(\code{elementHE}, $h_3$ ou $h_1$). Nos pseudocódigos, $\V[h]\leftarrow w$ corresponde a \code{setElementVertex(h, w)},
``$\Op[x]\leftrightarrow y$'' a \code{setOpposite(x, y)} com $y$ uma half-edge válida (liga nos dois sentidos) e
\code{duplicatePoint}$(v)$ cria um vértice novo, cópia de $v$, com as mesmas coordenadas.

As figuras desta seção foram geradas executando o código em malhas pequenas (\cref{app:repro}). Convenções visuais:
elementos contínuos em azul, elementos de interface em laranja, aresta a inserir em vermelho; vértices com aresta
colapsada têm contorno vermelho; vértices que pertencem a algum elemento de interface têm rótulo laranja; as cópias de
um vértice $X$ são chamadas $X_1, X_2,\dots$ na ordem de criação e são desenhadas ligeiramente afastadas umas das
outras. Arestas colapsadas (comprimento zero) aparecem como um pequeno segmento vermelho, com uma espessura
\emph{apenas visual}.

Todos os operadores usam a rotina de apoio \code{reindexElementsCCW} (\cref{alg:reindex}), que gira no sentido
anti-horário a partir de uma half-edge trocando o vértice de origem até encontrar uma half-edge de parada, a borda ou o
ponto de partida.

\begin{algorithm}[ht]
\caption{\code{reindexElementsCCW}$(h, \mathit{stop}, w)$}
\label{alg:reindex}
\begin{algorithmic}[1]
\State $s\gets h$
\Repeat
  \State $\V[h]\gets w$ \Comment{o canto de $h$ passa a usar o vértice $w$}
  \State $h\gets \Op[\prv(h)]$ \Comment{próximo elemento no sentido anti-horário}
\Until{$h=\mathit{stop}$ \textbf{ou} $h=\BORDER$ \textbf{ou} $h=s$}
\end{algorithmic}
\end{algorithm}

\subsection{Canônico}
\label{sec:canonico}

\textbf{Quando:} o vértice $v$ ainda não pertence a nenhum elemento de interface. É a ``inicialização'' do vértice.

\textbf{O que faz:} se $v$ é interno, nada é duplicado: o lado curto do novo elemento vira uma aresta colapsada de $v$
para $v$, com oposta \COLLAPSED, e ela é registrada em \code{\_collapsed[v]}. A trinca termina em $v$ e o elemento de
interface fica ``fechado'' ali. Se $v$ está na borda, a abertura chega ao exterior: cria-se um vértice $v'$, todos os
elementos do lado de \code{edgeHE} até a borda passam a usá-lo (rotação anti-horária), e o lado curto do novo elemento
fica na borda (\cref{alg:canonical,fig:canonico}).

\begin{algorithm}[ht]
\caption{\code{canonical}$(e, x)$ \quad ($e$ = \code{edgeHE}, $x$ = \code{elementHE})}
\label{alg:canonical}
\begin{algorithmic}[1]
\State $v\gets\V[e]$
\If{\code{isBorder}$(e)$} \Comment{gira em torno de $v$ procurando uma half-edge sem oposta}
  \State $v'\gets$ \code{duplicatePoint}$(v)$
  \State $\V[x]\gets v'$
  \State \code{reindexElementsCCW}$(e, \BORDER, v')$ \Comment{lado de $e$, até a borda}
  \State $\Op[x]\gets\BORDER$; \quad \code{\_inInterfaceElement}$[v']\gets$ \textbf{true}
\Else
  \State $\V[x]\gets v$; \quad $\Op[x]\gets\COLLAPSED$; \quad \code{\_collapsed}$[v]\gets x$
\EndIf
\State $\V[\nxt(x)]\gets v$; \quad \code{\_inInterfaceElement}$[v]\gets$ \textbf{true}
\end{algorithmic}
\end{algorithm}

\begin{figure}[ht]
\centering
\begin{subfigure}[b]{0.28\linewidth}\centering\begin{tikzpicture}[scale=1.400, every node/.style={font=\scriptsize}]
\filldraw[cefill] (0.000,0.000) -- (1.000,0.000) -- (1.000,1.000) -- (0.000,1.000) -- cycle;
\filldraw[cefill] (1.000,0.000) -- (2.000,0.000) -- (2.000,1.000) -- (1.000,1.000) -- cycle;
\filldraw[cefill] (2.000,0.000) -- (3.000,0.000) -- (3.000,1.000) -- (2.000,1.000) -- cycle;
\filldraw[cefill] (0.000,1.000) -- (1.000,1.000) -- (1.000,2.000) -- (0.000,2.000) -- cycle;
\filldraw[cefill] (1.000,1.000) -- (2.000,1.000) -- (2.000,2.000) -- (1.000,2.000) -- cycle;
\filldraw[cefill] (2.000,1.000) -- (3.000,1.000) -- (3.000,2.000) -- (2.000,2.000) -- cycle;
\node[elemlabel] at (1.500,1.500) {U};
\node[elemlabel] at (1.500,0.500) {L};
\draw[seledge] (1.000,1.000) -- (2.000,1.000);
\draw[hearrow, red!70!black] (1.220,1.140) -- (1.780,1.140);
\node[helabel, text=red!70!black] at (1.500,1.266) {$a$};
\draw[hearrow, blue!70!black] (1.780,0.860) -- (1.220,0.860);
\node[helabel, text=blue!70!black] at (1.500,0.734) {$b$};
\node[vreg, dot] at (0.000,0.000) {};
\node[vreg, dot] at (1.000,0.000) {};
\node[vreg, dot] at (2.000,0.000) {};
\node[vreg, dot] at (3.000,0.000) {};
\node[vreg, dot] at (0.000,1.000) {};
\node[vreg, minimum size=0.31cm] at (1.000,1.000) {$A$};
\node[vreg, minimum size=0.31cm] at (2.000,1.000) {$B$};
\node[vreg, dot] at (3.000,1.000) {};
\node[vreg, dot] at (0.000,2.000) {};
\node[vreg, dot] at (1.000,2.000) {};
\node[vreg, dot] at (2.000,2.000) {};
\node[vreg, dot] at (3.000,2.000) {};

\end{tikzpicture}
\caption{Antes: $A$ e $B$ internos.}\end{subfigure}\hfill
\begin{subfigure}[b]{0.28\linewidth}\centering\begin{tikzpicture}[scale=1.400, every node/.style={font=\scriptsize}]
\filldraw[cefill] (0.000,0.000) -- (1.000,0.000) -- (1.000,1.000) -- (0.000,1.000) -- cycle;
\filldraw[cefill] (1.000,0.000) -- (2.000,0.000) -- (2.000,1.000) -- (1.000,1.000) -- cycle;
\filldraw[cefill] (2.000,0.000) -- (3.000,0.000) -- (3.000,1.000) -- (2.000,1.000) -- cycle;
\filldraw[cefill] (0.000,1.000) -- (1.000,1.000) -- (1.000,2.000) -- (0.000,2.000) -- cycle;
\filldraw[cefill] (1.000,1.000) -- (2.000,1.000) -- (2.000,2.000) -- (1.000,2.000) -- cycle;
\filldraw[cefill] (2.000,1.000) -- (3.000,1.000) -- (3.000,2.000) -- (2.000,2.000) -- cycle;
\filldraw[iefill] (1.085,0.915) -- (1.915,0.915) -- (1.915,1.085) -- (1.085,1.085) -- cycle;
\draw[red, line width=1.1pt] (1.915,0.915) -- (1.915,1.085);
\draw[red, line width=1.1pt] (1.085,1.085) -- (1.085,0.915);
\node[elemlabel] at (1.500,1.500) {U};
\node[elemlabel] at (1.500,0.500) {L};
\draw[hearrow, red!70!black] (1.220,1.220) -- (1.780,1.220);
\node[helabel, text=red!70!black] at (1.276,1.352) {$a$};
\draw[hearrow, blue!70!black] (1.780,0.780) -- (1.220,0.780);
\node[helabel, text=blue!70!black] at (1.724,0.648) {$b$};
\draw[hearrow, black] (1.285,1.076) -- (1.285,0.924);
\node[helabel, text=black] at (0.865,1.200) {$h_3$};
\draw[hearrow, black] (1.715,0.924) -- (1.715,1.076);
\node[helabel, text=black] at (2.135,1.200) {$h_1$};
\node[vreg, dot] at (0.000,0.000) {};
\node[vreg, dot] at (1.000,0.000) {};
\node[vreg, dot] at (2.000,0.000) {};
\node[vreg, dot] at (3.000,0.000) {};
\node[vreg, dot] at (0.000,1.000) {};
\node[vcol, minimum size=0.31cm] at (1.000,1.000) {$A$};
\node[vcol, minimum size=0.31cm] at (2.000,1.000) {$B$};
\node[vreg, dot] at (3.000,1.000) {};
\node[vreg, dot] at (0.000,2.000) {};
\node[vreg, dot] at (1.000,2.000) {};
\node[vreg, dot] at (2.000,2.000) {};
\node[vreg, dot] at (3.000,2.000) {};

\end{tikzpicture}
\caption{Depois: elemento $[A,B,B,A]$.}\end{subfigure}\hfill
\begin{subfigure}[b]{0.2\linewidth}\centering\begin{tikzpicture}[scale=1.400, every node/.style={font=\scriptsize}]
\filldraw[cefill] (0.000,0.000) -- (1.000,0.000) -- (1.000,1.000) -- (0.000,1.000) -- cycle;
\filldraw[cefill] (1.000,0.000) -- (2.000,0.000) -- (2.000,1.000) -- (1.000,1.000) -- cycle;
\filldraw[cefill] (0.000,1.000) -- (1.000,1.000) -- (1.000,2.000) -- (0.000,2.000) -- cycle;
\filldraw[cefill] (1.000,1.000) -- (2.000,1.000) -- (2.000,2.000) -- (1.000,2.000) -- cycle;
\draw[seledge] (1.000,1.000) -- (1.000,2.000);
\draw[hearrow, red!70!black] (0.860,1.220) -- (0.860,1.780);
\node[helabel, text=red!70!black] at (0.734,1.500) {$a$};
\draw[hearrow, blue!70!black] (1.140,1.780) -- (1.140,1.220);
\node[helabel, text=blue!70!black] at (1.266,1.500) {$b$};
\node[vreg, dot] at (0.000,0.000) {};
\node[vreg, dot] at (1.000,0.000) {};
\node[vreg, dot] at (2.000,0.000) {};
\node[vreg, dot] at (0.000,1.000) {};
\node[vreg, minimum size=0.31cm] at (1.000,1.000) {$A$};
\node[vreg, dot] at (2.000,1.000) {};
\node[vreg, dot] at (0.000,2.000) {};
\node[vreg, minimum size=0.31cm] at (1.000,2.000) {$B$};
\node[vreg, dot] at (2.000,2.000) {};

\end{tikzpicture}
\caption{Antes: $B$ na borda.}\end{subfigure}\hfill
\begin{subfigure}[b]{0.2\linewidth}\centering\begin{tikzpicture}[scale=1.400, every node/.style={font=\scriptsize}]
\filldraw[cefill] (0.000,0.000) -- (1.000,0.000) -- (1.000,1.000) -- (0.000,1.000) -- cycle;
\filldraw[cefill] (1.000,0.000) -- (2.000,0.000) -- (2.000,1.000) -- (1.000,1.000) -- cycle;
\filldraw[cefill] (0.000,1.000) -- (1.000,1.000) -- (0.844,1.844) -- (0.000,2.000) -- cycle;
\filldraw[cefill] (1.000,1.000) -- (2.000,1.000) -- (2.000,2.000) -- (1.156,1.844) -- cycle;
\filldraw[iefill] (1.067,1.057) -- (1.156,1.844) -- (0.844,1.844) -- (0.933,1.057) -- cycle;
\draw[red, line width=1.1pt] (0.933,1.057) -- (1.067,1.057);
\draw[hearrow, red!70!black] (0.826,1.198) -- (0.739,1.670);
\node[helabel, text=red!70!black] at (0.657,1.445) {$a$};
\draw[hearrow, blue!70!black] (1.261,1.670) -- (1.174,1.198);
\node[helabel, text=blue!70!black] at (1.343,1.445) {$b$};
\node[vreg, dot] at (0.000,0.000) {};
\node[vreg, dot] at (1.000,0.000) {};
\node[vreg, dot] at (2.000,0.000) {};
\node[vreg, dot] at (0.000,1.000) {};
\node[vcol, minimum size=0.31cm] at (1.000,1.000) {$A$};
\node[vreg, dot] at (2.000,1.000) {};
\node[vreg, dot] at (0.000,2.000) {};
\node[vie, minimum size=0.31cm] at (0.844,1.844) {$B$};
\node[vreg, dot] at (2.000,2.000) {};
\node[vie, minimum size=0.31cm] at (1.156,1.844) {$B_{1}$};

\end{tikzpicture}
\caption{Depois.}\end{subfigure}
\caption{Operador Canônico. (a--b) Com os dois vértices internos, o elemento de interface tem as duas extremidades
colapsadas ($h_3$ em $A$ e $h_1$ em $B$, em vermelho) e nenhum vértice é criado: geometricamente ele é só a aresta $AB$.
(c--d) Com $B$ na borda, $B$ é duplicado em $B_1$; o operador em $B$ recebe $b$, e os elementos do lado de $b$ (à
direita) passam a usar a cópia; o lado curto em $B$ é uma aresta de borda. Em $A$ (interno) a extremidade continua colapsada.}
\label{fig:canonico}
\end{figure}

\subsection{Split Element}
\label{sec:split}

\textbf{Quando:} $v$ tem uma aresta colapsada $c$ (pertence a exatamente um elemento de interface $I$, criado pelo
Canônico) \emph{e} $I$ é vizinho de um dos dois elementos que contêm a aresta nova. Esse elemento comum é o
\emph{elemento compartilhado} $S$.

\textbf{Como o código testa} (\code{retrieveAvailableVertex}, \cref{alg:retrieve}). Seja $he$ a half-edge nova que parte
de $v$ e $b=\Op[he]$. Há duas posições possíveis para $S$:
\begin{itemize}
\item \textbf{hipótese $x_1$}: $S$ é o elemento de $he$. Então a aresta de $S$ que chega em $v$, $\prv(he)$, é vizinha
de $I$, e sua oposta $x_1=\Op[\prv(he)]$ é o lado longo de $I$ que \emph{sai} de $v$ logo depois de $c$, isto é,
$\prv(x_1)=c$;
\item \textbf{hipótese $x_3$}: $S$ é o elemento de $b$. Então a aresta de $S$ que sai de $v$, $\nxt(b)$, é vizinha de
$I$, e $x_3=\Op[\nxt(b)]$ é o lado longo de $I$ que \emph{chega} em $v$ logo antes de $c$, isto é, $\nxt(x_3)=c$.
\end{itemize}
Qual hipótese ocorre não depende de qual half-edge da aresta o usuário escolheu, e sim da posição da aresta nova em
relação a $I$: em $S$, se a aresta nova vem depois da aresta de $I$ (no sentido anti-horário de $S$) vale $x_1$; se vem
antes, vale $x_3$. A \cref{fig:hipoteses} mostra, em detalhe, as consultas feitas em cada caso e também o caso em que
nenhuma das duas vale; a \cref{fig:split} mostra o resultado numa malha.

% Base comum das três figuras: vértice v na origem, elemento de interface I à esquerda (colapsado em v) e quatro
% elementos contínuos em volta. \hipS{x}{y} rotula o elemento compartilhado.
\newcommand{\hipbase}[1]{%
  \filldraw[cefill] (-2.8,0.6) -- (0,0) -- (0,2.4) -- (-2.8,2.4) -- cycle;
  \filldraw[cefill] (-2.8,-2.4) -- (0,-2.4) -- (0,0) -- (-2.8,-0.6) -- cycle;
  \filldraw[cefill] (0,0) -- (2.6,0) -- (2.6,2.4) -- (0,2.4) -- cycle;
  \filldraw[cefill] (0,-2.4) -- (2.6,-2.4) -- (2.6,0) -- (0,0) -- cycle;
  \filldraw[iefill] (-2.8,-0.6) -- (0,0) -- (-2.8,0.6) -- cycle;
  \draw[white, line width=1.2pt] (-2.8,-0.55) -- (-2.8,0.55);
  \draw[orange!55!black, dashed, line width=0.45pt] (-2.8,-0.6) -- (-2.8,0.6);
  \node[font=\footnotesize\bfseries, text=orange!40!black] at (-2.35,0) {$I$};
  #1
  \node[vcol, minimum size=0.32cm] at (0,0) {};
  \node[font=\scriptsize] at (0.28,-0.22) {$v$};
}
\newcommand{\hipc}[2]{% rótulo da aresta colapsada c, com uma linha até v
  \draw[red!70!black, line width=0.4pt] (0.05,0.03*#2) -- (#1,0.5*#2);
  \node[font=\scriptsize, text=red!70!black, anchor=west, align=left] at (#1,0.5*#2) {$c$ (colapsada,\\ de $v$ para $v$)};
}

\begin{figure}[ht]
\centering
\begin{subfigure}[b]{0.325\linewidth}\centering
\begin{tikzpicture}[scale=0.88, every node/.style={font=\scriptsize}]
  \hipbase{\draw[curedge] (0,0) -- (0,2.4);}
  \node[elemlabel] at (-1.5,1.65) {S};
  \draw[hearrow, red!70!black, line width=0.8pt] (-0.17,0.45) -- (-0.17,2.0) node[midway, left=0pt] {$he$};
  \draw[hearrow, blue!70!black] (0.17,2.0) -- (0.17,0.45) node[midway, right=0pt] {$b$};
  \draw[hearrow, violet, line width=0.8pt] (-2.4,0.68) -- (-0.5,0.28) node[pos=0.45, above=1pt, sloped] {$\prv(he)$};
  \draw[hearrow, green!45!black, line width=0.8pt] (-0.9,0.07) -- (-2.5,0.42);
  \node[text=green!45!black] at (-1.75,0.03) {$x_1$};
  \hipc{0.75}{-1}
\end{tikzpicture}
\caption{Hipótese $x_1$: $S$ é o elemento de $he$.\\ $x_1=\Op[\prv(he)]$ e $\prv(x_1)=c$.}
\end{subfigure}\hfill
\begin{subfigure}[b]{0.325\linewidth}\centering
\begin{tikzpicture}[scale=0.88, every node/.style={font=\scriptsize}]
  \hipbase{\draw[curedge] (0,0) -- (0,-2.4);}
  \node[elemlabel] at (-1.5,-1.65) {S};
  \draw[hearrow, red!70!black, line width=0.8pt] (0.17,-0.45) -- (0.17,-2.0) node[midway, right=0pt] {$he$};
  \draw[hearrow, blue!70!black] (-0.17,-2.0) -- (-0.17,-0.45) node[midway, left=0pt] {$b$};
  \draw[hearrow, violet, line width=0.8pt] (-0.5,-0.28) -- (-2.4,-0.68) node[pos=0.55, below=1pt, sloped] {$\nxt(b)$};
  \draw[hearrow, green!45!black, line width=0.8pt] (-2.5,-0.42) -- (-0.9,-0.07);
  \node[text=green!45!black] at (-1.75,-0.03) {$x_3$};
  \hipc{0.75}{1}
\end{tikzpicture}
\caption{Hipótese $x_3$: $S$ é o elemento de $b$.\\ $x_3=\Op[\nxt(b)]$ e $\nxt(x_3)=c$.}
\end{subfigure}\hfill
\begin{subfigure}[b]{0.325\linewidth}\centering
\begin{tikzpicture}[scale=0.88, every node/.style={font=\scriptsize}]
  \hipbase{\draw[curedge] (0,0) -- (2.6,0);}
  \draw[hearrow, red!70!black, line width=0.8pt] (0.45,0.17) -- (2.2,0.17) node[midway, above=0pt] {$he$};
  \draw[hearrow, blue!70!black] (2.2,-0.17) -- (0.45,-0.17) node[midway, below=0pt] {$b$};
  \draw[hearrow, violet, line width=0.8pt] (0.17,2.0) -- (0.17,0.45) node[pos=0.3, right=0pt] {$\prv(he)$};
  \draw[hearrow, black!50, dashed] (-0.17,0.45) -- (-0.17,2.0) node[pos=0.7, left=0pt, text=black!60] {$\Op[\prv(he)]$};
  \draw[hearrow, violet, line width=0.8pt] (0.17,-0.45) -- (0.17,-2.0) node[pos=0.7, right=0pt] {$\nxt(b)$};
  \draw[hearrow, black!50, dashed] (-0.17,-2.0) -- (-0.17,-0.45) node[pos=0.3, left=0pt, text=black!60] {$\Op[\nxt(b)]$};
\end{tikzpicture}
\caption{Nenhuma: os opostos caem fora de $I$.\\ $\Rightarrow$ Expand Edge.}
\end{subfigure}
\caption{As consultas de \code{retrieveAvailableVertex} em torno de um vértice $v$ que tem uma aresta colapsada $c$ (o
elemento de interface $I$, à esquerda, termina em $v$; a aresta nova está em vermelho grosso). (a) A aresta nova sobe:
o elemento de $he$ tem uma aresta, $\prv(he)$, colada ao lado longo de $I$ que sai de $v$ logo depois de $c$. (b) A
aresta nova desce: agora é o elemento de $b$ que encosta em $I$, pela aresta $\nxt(b)$, colada ao lado longo de $I$ que
chega em $v$ logo antes de $c$. (c) A aresta nova continua a trinca: os dois testes levam a elementos contínuos, e o
operador escolhido é o Expand Edge. Em (a) e (b), o vértice novo vai para $S$.}
\label{fig:hipoteses}
\end{figure}

\textbf{O que faz} (\cref{alg:split}): cria um vértice $n$ e o coloca ``dentro'' de $S$: só dois cantos mudam, o canto de
$S$ em $v$ e o canto de $I$ adjacente a $S$. A antiga aresta colapsada $c$ deixa de ter comprimento zero (passa a ligar
$v$ e $n$) e vira a fronteira entre $I$ e o novo elemento, sendo ligada a \code{elementHE}. O custo é um número
constante de passos ($O(\log n_c)$ na IBHM, por causa das consultas $\nxt$/$\prv$).

\begin{algorithm}[ht]
\caption{\code{retrieveAvailableVertex}$(he)$ e \code{splitElement}$(he, \mathit{av}, \mathit{sh}, x)$}
\label{alg:retrieve}\label{alg:split}
\begin{algorithmic}[1]
\Function{retrieveAvailableVertex}{$he$}
  \State $v\gets\V[he]$; \quad $b\gets\Op[he]$; \quad $c\gets$ \code{\_collapsed}$[v]$
  \State $x_1\gets\Op[\prv(he)]$
  \If{$x_1\ne\BORDER$ \textbf{e} $\prv(x_1)=c$} \Return $(\mathit{av}=x_1,\ \mathit{sh}=he)$ \Comment{hipótese $x_1$}
  \EndIf
  \State $x_3\gets\Op[\nxt(b)]$
  \If{$x_3\ne\BORDER$ \textbf{e} $\nxt(x_3)=c$} \Return $(\mathit{av}=c,\ \mathit{sh}=\nxt(b))$ \Comment{hipótese $x_3$}
  \EndIf
  \State \Return $(\BORDER,\BORDER)$ \Comment{nenhum elemento compartilhado $\Rightarrow$ Expand Edge}
\EndFunction
\Statex
\Procedure{splitElement}{$he, \mathit{av}, \mathit{sh}, x$}
  \State $v\gets\V[he]$; \quad $n\gets$ \code{duplicatePoint}$(v)$
  \State $\V[\mathit{av}]\gets n$ \Comment{canto do elemento de interface antigo}
  \State $\V[\mathit{sh}]\gets n$ \Comment{canto do elemento compartilhado}
  \If{$he=\mathit{sh}$} \Comment{hipótese $x_1$}
    \State $\V[x]\gets n$; \quad $\V[\nxt(x)]\gets v$; \quad $\Op[x]\leftrightarrow\prv(\mathit{av})$ \Comment{$\prv(\mathit{av})=c$}
  \Else \Comment{hipótese $x_3$}
    \State $\V[x]\gets v$; \quad $\V[\nxt(x)]\gets n$; \quad $\Op[x]\leftrightarrow \mathit{av}$ \Comment{$\mathit{av}=c$}
  \EndIf
  \State marca $v$ e $n$ como vértices de EI; remove $v$ de \code{\_collapsed}
\EndProcedure
\end{algorithmic}
\end{algorithm}

\begin{figure}[ht]
\centering
\begin{subfigure}[b]{0.31\linewidth}\centering\begin{tikzpicture}[scale=2.000, every node/.style={font=\scriptsize}]
\clip (0.450,0.550) rectangle (2.800,3.450);
\filldraw[cefill] (0.000,0.000) -- (1.000,0.000) -- (1.000,1.000) -- (0.000,1.000) -- cycle;
\filldraw[cefill] (1.000,0.000) -- (2.000,0.000) -- (2.000,1.000) -- (1.000,1.000) -- cycle;
\filldraw[cefill] (2.000,0.000) -- (3.000,0.000) -- (3.000,1.000) -- (2.000,1.000) -- cycle;
\filldraw[cefill] (0.000,1.000) -- (1.000,1.000) -- (1.000,2.000) -- (0.000,2.000) -- cycle;
\filldraw[cefill] (1.000,1.000) -- (2.000,1.000) -- (2.000,2.000) -- (1.000,2.000) -- cycle;
\filldraw[cefill] (2.000,1.000) -- (3.000,1.000) -- (3.000,2.000) -- (2.000,2.000) -- cycle;
\filldraw[cefill] (0.000,2.000) -- (1.000,2.000) -- (1.000,3.000) -- (0.000,3.000) -- cycle;
\filldraw[cefill] (1.000,2.000) -- (2.000,2.000) -- (2.000,3.000) -- (1.000,3.000) -- cycle;
\filldraw[cefill] (2.000,2.000) -- (3.000,2.000) -- (3.000,3.000) -- (2.000,3.000) -- cycle;
\filldraw[cefill] (0.000,3.000) -- (1.000,3.000) -- (1.000,4.000) -- (0.000,4.000) -- cycle;
\filldraw[cefill] (1.000,3.000) -- (2.000,3.000) -- (2.000,4.000) -- (1.000,4.000) -- cycle;
\filldraw[cefill] (2.000,3.000) -- (3.000,3.000) -- (3.000,4.000) -- (2.000,4.000) -- cycle;
\filldraw[iefill] (1.062,1.938) -- (1.938,1.938) -- (1.938,2.062) -- (1.062,2.062) -- cycle;
\draw[red, line width=1.1pt] (1.938,1.938) -- (1.938,2.062);
\draw[red, line width=1.1pt] (1.062,2.062) -- (1.062,1.938);
\node[elemlabel] at (1.500,2.500) {S};
\draw[seledge] (2.000,2.000) -- (2.000,3.000);
\draw[hearrow, red!70!black] (1.860,2.220) -- (1.860,2.780);
\node[helabel, text=red!70!black] at (1.692,2.500) {$he$};
\draw[hearrow, blue!70!black] (2.140,2.780) -- (2.140,2.220);
\node[helabel, text=blue!70!black] at (2.266,2.500) {$b$};
\draw[hearrow, violet] (1.220,2.140) -- (1.780,2.140);
\node[helabel, text=violet] at (1.500,2.364) {$\mathrm{prev}(he)$};
\draw[hearrow, green!45!black] (1.745,2.034) -- (1.255,2.034);
\node[helabel, text=green!45!black] at (1.500,1.961) {$x_1$};
\node[vreg, dot] at (0.000,0.000) {};
\node[vreg, dot] at (1.000,0.000) {};
\node[vreg, dot] at (2.000,0.000) {};
\node[vreg, dot] at (3.000,0.000) {};
\node[vreg, dot] at (0.000,1.000) {};
\node[vreg, dot] at (1.000,1.000) {};
\node[vreg, minimum size=0.29cm] at (2.000,1.000) {$D$};
\node[vreg, dot] at (3.000,1.000) {};
\node[vreg, dot] at (0.000,2.000) {};
\node[vcol, minimum size=0.29cm] at (1.000,2.000) {$A$};
\node[vcol, minimum size=0.29cm] at (2.000,2.000) {$B$};
\node[vreg, dot] at (3.000,2.000) {};
\node[vreg, dot] at (0.000,3.000) {};
\node[vreg, dot] at (1.000,3.000) {};
\node[vreg, minimum size=0.29cm] at (2.000,3.000) {$C$};
\node[vreg, dot] at (3.000,3.000) {};
\node[vreg, dot] at (0.000,4.000) {};
\node[vreg, dot] at (1.000,4.000) {};
\node[vreg, dot] at (2.000,4.000) {};
\node[vreg, dot] at (3.000,4.000) {};

\end{tikzpicture}
\caption{$x_1$: antes.}\end{subfigure}\hspace{0.06\linewidth}
\begin{subfigure}[b]{0.31\linewidth}\centering\begin{tikzpicture}[scale=2.000, every node/.style={font=\scriptsize}]
\clip (0.450,0.550) rectangle (2.800,3.450);
\filldraw[cefill] (0.000,0.000) -- (1.000,0.000) -- (1.000,1.000) -- (0.000,1.000) -- cycle;
\filldraw[cefill] (1.000,0.000) -- (2.000,0.000) -- (2.000,1.000) -- (1.000,1.000) -- cycle;
\filldraw[cefill] (2.000,0.000) -- (3.000,0.000) -- (3.000,1.000) -- (2.000,1.000) -- cycle;
\filldraw[cefill] (0.000,1.000) -- (1.000,1.000) -- (1.000,2.000) -- (0.000,2.000) -- cycle;
\filldraw[cefill] (1.000,1.000) -- (2.000,1.000) -- (2.156,1.844) -- (1.000,2.000) -- cycle;
\filldraw[cefill] (2.000,1.000) -- (3.000,1.000) -- (3.000,2.000) -- (2.156,1.844) -- cycle;
\filldraw[cefill] (0.000,2.000) -- (1.000,2.000) -- (1.000,3.000) -- (0.000,3.000) -- cycle;
\filldraw[cefill] (1.000,2.000) -- (1.844,2.156) -- (2.000,3.000) -- (1.000,3.000) -- cycle;
\filldraw[cefill] (2.156,1.844) -- (3.000,2.000) -- (3.000,3.000) -- (2.000,3.000) -- cycle;
\filldraw[cefill] (0.000,3.000) -- (1.000,3.000) -- (1.000,4.000) -- (0.000,4.000) -- cycle;
\filldraw[cefill] (1.000,3.000) -- (2.000,3.000) -- (2.000,4.000) -- (1.000,4.000) -- cycle;
\filldraw[cefill] (2.000,3.000) -- (3.000,3.000) -- (3.000,4.000) -- (2.000,4.000) -- cycle;
\filldraw[iefill] (1.062,1.938) -- (2.156,1.844) -- (1.844,2.156) -- (1.057,2.067) -- cycle;
\filldraw[iefill] (2.156,1.844) -- (2.062,2.938) -- (1.933,2.943) -- (1.844,2.156) -- cycle;
\draw[red, line width=1.1pt] (1.057,2.067) -- (1.062,1.938);
\draw[red, line width=1.1pt] (2.062,2.938) -- (1.933,2.943);
\node[elemlabel] at (1.461,2.539) {S};
\node[elemlabel] at (1.999,2.470) {};
\node[vreg, dot] at (0.000,0.000) {};
\node[vreg, dot] at (1.000,0.000) {};
\node[vreg, dot] at (2.000,0.000) {};
\node[vreg, dot] at (3.000,0.000) {};
\node[vreg, dot] at (0.000,1.000) {};
\node[vreg, dot] at (1.000,1.000) {};
\node[vreg, minimum size=0.29cm] at (2.000,1.000) {$D$};
\node[vreg, dot] at (3.000,1.000) {};
\node[vreg, dot] at (0.000,2.000) {};
\node[vcol, minimum size=0.29cm] at (1.000,2.000) {$A$};
\node[vie, minimum size=0.29cm] at (2.156,1.844) {$B$};
\node[vreg, dot] at (3.000,2.000) {};
\node[vreg, dot] at (0.000,3.000) {};
\node[vreg, dot] at (1.000,3.000) {};
\node[vcol, minimum size=0.29cm] at (2.000,3.000) {$C$};
\node[vreg, dot] at (3.000,3.000) {};
\node[vreg, dot] at (0.000,4.000) {};
\node[vreg, dot] at (1.000,4.000) {};
\node[vreg, dot] at (2.000,4.000) {};
\node[vreg, dot] at (3.000,4.000) {};
\node[vie, minimum size=0.29cm] at (1.844,2.156) {$B_{1}$};

\end{tikzpicture}
\caption{$x_1$: depois.}\end{subfigure}\\[6pt]
\begin{subfigure}[b]{0.31\linewidth}\centering\begin{tikzpicture}[scale=2.000, every node/.style={font=\scriptsize}]
\clip (0.450,0.550) rectangle (2.800,3.450);
\filldraw[cefill] (0.000,0.000) -- (1.000,0.000) -- (1.000,1.000) -- (0.000,1.000) -- cycle;
\filldraw[cefill] (1.000,0.000) -- (2.000,0.000) -- (2.000,1.000) -- (1.000,1.000) -- cycle;
\filldraw[cefill] (2.000,0.000) -- (3.000,0.000) -- (3.000,1.000) -- (2.000,1.000) -- cycle;
\filldraw[cefill] (0.000,1.000) -- (1.000,1.000) -- (1.000,2.000) -- (0.000,2.000) -- cycle;
\filldraw[cefill] (1.000,1.000) -- (2.000,1.000) -- (2.000,2.000) -- (1.000,2.000) -- cycle;
\filldraw[cefill] (2.000,1.000) -- (3.000,1.000) -- (3.000,2.000) -- (2.000,2.000) -- cycle;
\filldraw[cefill] (0.000,2.000) -- (1.000,2.000) -- (1.000,3.000) -- (0.000,3.000) -- cycle;
\filldraw[cefill] (1.000,2.000) -- (2.000,2.000) -- (2.000,3.000) -- (1.000,3.000) -- cycle;
\filldraw[cefill] (2.000,2.000) -- (3.000,2.000) -- (3.000,3.000) -- (2.000,3.000) -- cycle;
\filldraw[cefill] (0.000,3.000) -- (1.000,3.000) -- (1.000,4.000) -- (0.000,4.000) -- cycle;
\filldraw[cefill] (1.000,3.000) -- (2.000,3.000) -- (2.000,4.000) -- (1.000,4.000) -- cycle;
\filldraw[cefill] (2.000,3.000) -- (3.000,3.000) -- (3.000,4.000) -- (2.000,4.000) -- cycle;
\filldraw[iefill] (1.062,1.938) -- (1.938,1.938) -- (1.938,2.062) -- (1.062,2.062) -- cycle;
\draw[red, line width=1.1pt] (1.938,1.938) -- (1.938,2.062);
\draw[red, line width=1.1pt] (1.062,2.062) -- (1.062,1.938);
\node[elemlabel] at (1.500,1.500) {S};
\draw[seledge] (2.000,2.000) -- (2.000,1.000);
\draw[hearrow, red!70!black] (2.140,1.780) -- (2.140,1.220);
\node[helabel, text=red!70!black] at (2.308,1.500) {$he$};
\draw[hearrow, blue!70!black] (1.860,1.220) -- (1.860,1.780);
\node[helabel, text=blue!70!black] at (1.734,1.500) {$b$};
\draw[hearrow, violet] (1.780,1.860) -- (1.220,1.860);
\node[helabel, text=violet] at (1.500,1.636) {$\mathrm{next}(b)$};
\draw[hearrow, green!45!black] (1.255,1.966) -- (1.745,1.966);
\node[helabel, text=green!45!black] at (1.500,2.039) {$x_3$};
\node[vreg, dot] at (0.000,0.000) {};
\node[vreg, dot] at (1.000,0.000) {};
\node[vreg, dot] at (2.000,0.000) {};
\node[vreg, dot] at (3.000,0.000) {};
\node[vreg, dot] at (0.000,1.000) {};
\node[vreg, dot] at (1.000,1.000) {};
\node[vreg, minimum size=0.29cm] at (2.000,1.000) {$D$};
\node[vreg, dot] at (3.000,1.000) {};
\node[vreg, dot] at (0.000,2.000) {};
\node[vcol, minimum size=0.29cm] at (1.000,2.000) {$A$};
\node[vcol, minimum size=0.29cm] at (2.000,2.000) {$B$};
\node[vreg, dot] at (3.000,2.000) {};
\node[vreg, dot] at (0.000,3.000) {};
\node[vreg, dot] at (1.000,3.000) {};
\node[vreg, minimum size=0.29cm] at (2.000,3.000) {$C$};
\node[vreg, dot] at (3.000,3.000) {};
\node[vreg, dot] at (0.000,4.000) {};
\node[vreg, dot] at (1.000,4.000) {};
\node[vreg, dot] at (2.000,4.000) {};
\node[vreg, dot] at (3.000,4.000) {};

\end{tikzpicture}
\caption{$x_3$: antes.}\end{subfigure}\hspace{0.06\linewidth}
\begin{subfigure}[b]{0.31\linewidth}\centering\begin{tikzpicture}[scale=2.000, every node/.style={font=\scriptsize}]
\clip (0.450,0.550) rectangle (2.800,3.450);
\filldraw[cefill] (0.000,0.000) -- (1.000,0.000) -- (1.000,1.000) -- (0.000,1.000) -- cycle;
\filldraw[cefill] (1.000,0.000) -- (2.000,0.000) -- (2.000,1.000) -- (1.000,1.000) -- cycle;
\filldraw[cefill] (2.000,0.000) -- (3.000,0.000) -- (3.000,1.000) -- (2.000,1.000) -- cycle;
\filldraw[cefill] (0.000,1.000) -- (1.000,1.000) -- (1.000,2.000) -- (0.000,2.000) -- cycle;
\filldraw[cefill] (1.000,1.000) -- (2.000,1.000) -- (1.844,1.844) -- (1.000,2.000) -- cycle;
\filldraw[cefill] (2.000,1.000) -- (3.000,1.000) -- (3.000,2.000) -- (2.156,2.156) -- cycle;
\filldraw[cefill] (0.000,2.000) -- (1.000,2.000) -- (1.000,3.000) -- (0.000,3.000) -- cycle;
\filldraw[cefill] (1.000,2.000) -- (2.156,2.156) -- (2.000,3.000) -- (1.000,3.000) -- cycle;
\filldraw[cefill] (2.156,2.156) -- (3.000,2.000) -- (3.000,3.000) -- (2.000,3.000) -- cycle;
\filldraw[cefill] (0.000,3.000) -- (1.000,3.000) -- (1.000,4.000) -- (0.000,4.000) -- cycle;
\filldraw[cefill] (1.000,3.000) -- (2.000,3.000) -- (2.000,4.000) -- (1.000,4.000) -- cycle;
\filldraw[cefill] (2.000,3.000) -- (3.000,3.000) -- (3.000,4.000) -- (2.000,4.000) -- cycle;
\filldraw[iefill] (1.057,1.933) -- (1.844,1.844) -- (2.156,2.156) -- (1.062,2.062) -- cycle;
\filldraw[iefill] (1.844,1.844) -- (1.933,1.057) -- (2.062,1.062) -- (2.156,2.156) -- cycle;
\draw[red, line width=1.1pt] (1.062,2.062) -- (1.057,1.933);
\draw[red, line width=1.1pt] (1.933,1.057) -- (2.062,1.062);
\node[elemlabel] at (1.461,1.461) {S};
\node[elemlabel] at (1.999,1.530) {};
\node[vreg, dot] at (0.000,0.000) {};
\node[vreg, dot] at (1.000,0.000) {};
\node[vreg, dot] at (2.000,0.000) {};
\node[vreg, dot] at (3.000,0.000) {};
\node[vreg, dot] at (0.000,1.000) {};
\node[vreg, dot] at (1.000,1.000) {};
\node[vcol, minimum size=0.29cm] at (2.000,1.000) {$D$};
\node[vreg, dot] at (3.000,1.000) {};
\node[vreg, dot] at (0.000,2.000) {};
\node[vcol, minimum size=0.29cm] at (1.000,2.000) {$A$};
\node[vie, minimum size=0.29cm] at (2.156,2.156) {$B$};
\node[vreg, dot] at (3.000,2.000) {};
\node[vreg, dot] at (0.000,3.000) {};
\node[vreg, dot] at (1.000,3.000) {};
\node[vreg, minimum size=0.29cm] at (2.000,3.000) {$C$};
\node[vreg, dot] at (3.000,3.000) {};
\node[vreg, dot] at (0.000,4.000) {};
\node[vreg, dot] at (1.000,4.000) {};
\node[vreg, dot] at (2.000,4.000) {};
\node[vreg, dot] at (3.000,4.000) {};
\node[vie, minimum size=0.29cm] at (1.844,1.844) {$B_{1}$};

\end{tikzpicture}
\caption{$x_3$: depois.}\end{subfigure}
\caption{Split Element em $B$. Nos dois casos já existe um elemento de interface em $AB$, colapsado em $A$ e em $B$.
(a--b) A aresta nova $BC$ sobe: o elemento compartilhado $S$ (acima de $AB$) contém $he$, e o lado de $I$ que sai de $B$
é $x_1$. (c--d) A aresta nova $BD$ desce: agora $S$ (abaixo de $AB$) contém $b$, e o lado de $I$ que chega em $B$ é
$x_3$. Em ambos, só $S$ e um canto de $I$ passam a usar a nova cópia $B_1$, e a antiga aresta colapsada em $B$ vira o
lado comum entre os dois elementos de interface. Note que $C$ e $D$, internos, ficam colapsados.}
\label{fig:split}
\end{figure}

\subsection{Expand Edge}
\label{sec:expand}

\textbf{Quando:} $v$ tem uma aresta colapsada $c$, mas nenhuma das duas hipóteses vale: o elemento de interface $I$ não
toca os elementos da aresta nova. É o caso típico de uma trinca que continua ``em frente'' através de $v$
(\cref{fig:expand}).

\textbf{O que faz} (\cref{alg:expand}): a estrela de $v$ precisa ser dividida em dois setores, delimitados pela aresta
nova de um lado e por $I$ do outro. O operador cria $n$ e gira no sentido anti-horário \emph{a partir da aresta
colapsada} $c$: primeiro troca o canto de $c$ em $I$, depois sai de $I$ pelo seu lado longo $\prv(c)$ e percorre os
elementos daquele setor até alcançar $he$, que \emph{não} é alterado. Assim, a parada natural da rotação é justamente
cruzar a aresta nova. A aresta $c$ deixa de ser colapsada e é ligada a \code{elementHE}.

\begin{algorithm}[ht]
\caption{\code{expandEdge}$(he, x)$}
\label{alg:expand}
\begin{algorithmic}[1]
\State $v\gets\V[he]$; \quad $n\gets$ \code{duplicatePoint}$(v)$; \quad $c\gets$ \code{\_collapsed}$[v]$
\State \code{reindexElementsCCW}$(c, he, n)$ \Comment{de $c$ até $he$ (exclusive)}
\State $\Op[c]\leftrightarrow x$
\State $\V[x]\gets v$; \quad $\V[\nxt(x)]\gets n$
\State marca $v$ e $n$ como vértices de EI; remove $v$ de \code{\_collapsed}
\end{algorithmic}
\end{algorithm}

\begin{figure}[ht]
\centering
\begin{subfigure}[b]{0.32\linewidth}\centering\resizebox{\linewidth}{!}{\begin{tikzpicture}[scale=1.700, every node/.style={font=\scriptsize}]
\clip (0.450,0.300) rectangle (3.550,1.700);
\filldraw[cefill] (0.000,0.000) -- (1.000,0.000) -- (1.000,1.000) -- (0.000,1.000) -- cycle;
\filldraw[cefill] (1.000,0.000) -- (2.000,0.000) -- (2.000,1.000) -- (1.000,1.000) -- cycle;
\filldraw[cefill] (2.000,0.000) -- (3.000,0.000) -- (3.000,1.000) -- (2.000,1.000) -- cycle;
\filldraw[cefill] (3.000,0.000) -- (4.000,0.000) -- (4.000,1.000) -- (3.000,1.000) -- cycle;
\filldraw[cefill] (0.000,1.000) -- (1.000,1.000) -- (1.000,2.000) -- (0.000,2.000) -- cycle;
\filldraw[cefill] (1.000,1.000) -- (2.000,1.000) -- (2.000,2.000) -- (1.000,2.000) -- cycle;
\filldraw[cefill] (2.000,1.000) -- (3.000,1.000) -- (3.000,2.000) -- (2.000,2.000) -- cycle;
\filldraw[cefill] (3.000,1.000) -- (4.000,1.000) -- (4.000,2.000) -- (3.000,2.000) -- cycle;
\filldraw[iefill] (1.057,0.943) -- (1.943,0.943) -- (1.943,1.057) -- (1.057,1.057) -- cycle;
\draw[red, line width=1.1pt] (1.943,0.943) -- (1.943,1.057);
\draw[red, line width=1.1pt] (1.057,1.057) -- (1.057,0.943);
\draw[seledge] (2.000,1.000) -- (3.000,1.000);
\draw[hearrow, red!70!black] (2.220,1.140) -- (2.780,1.140);
\node[helabel, text=red!70!black] at (2.500,1.266) {$he$};
\draw[hearrow, orange!80!black] (1.803,0.968) -- (1.803,1.032);
\node[helabel, text=orange!80!black] at (1.677,1.000) {$c$};
\draw[hearrow, blue!70!black] (2.780,0.860) -- (2.220,0.860);
\node[helabel, text=blue!70!black] at (2.500,0.734) {$b$};
\node[vreg, dot] at (0.000,0.000) {};
\node[vreg, dot] at (1.000,0.000) {};
\node[vreg, dot] at (2.000,0.000) {};
\node[vreg, dot] at (3.000,0.000) {};
\node[vreg, dot] at (4.000,0.000) {};
\node[vreg, dot] at (0.000,1.000) {};
\node[vcol, minimum size=0.29cm] at (1.000,1.000) {$A$};
\node[vcol, minimum size=0.29cm] at (2.000,1.000) {$B$};
\node[vreg, minimum size=0.29cm] at (3.000,1.000) {$C$};
\node[vreg, dot] at (4.000,1.000) {};
\node[vreg, dot] at (0.000,2.000) {};
\node[vreg, dot] at (1.000,2.000) {};
\node[vreg, dot] at (2.000,2.000) {};
\node[vreg, dot] at (3.000,2.000) {};
\node[vreg, dot] at (4.000,2.000) {};

\end{tikzpicture}
}\caption{Antes.}\end{subfigure}\hfill
\begin{subfigure}[b]{0.32\linewidth}\centering\resizebox{\linewidth}{!}{\begin{tikzpicture}[scale=1.700, every node/.style={font=\scriptsize}]
\clip (0.450,0.300) rectangle (3.550,1.700);
\filldraw[cefill] (0.000,0.000) -- (1.000,0.000) -- (1.000,1.000) -- (0.000,1.000) -- cycle;
\filldraw[cefill] (1.000,0.000) -- (2.000,0.000) -- (2.000,1.000) -- (1.000,1.000) -- cycle;
\filldraw[cefill] (2.000,0.000) -- (3.000,0.000) -- (3.000,1.000) -- (2.000,1.000) -- cycle;
\filldraw[cefill] (3.000,0.000) -- (4.000,0.000) -- (4.000,1.000) -- (3.000,1.000) -- cycle;
\filldraw[cefill] (0.000,1.000) -- (1.000,1.000) -- (1.000,2.000) -- (0.000,2.000) -- cycle;
\filldraw[cefill] (1.000,1.000) -- (2.000,1.000) -- (2.000,2.000) -- (1.000,2.000) -- cycle;
\filldraw[cefill] (2.000,1.000) -- (3.000,1.000) -- (3.000,2.000) -- (2.000,2.000) -- cycle;
\filldraw[cefill] (3.000,1.000) -- (4.000,1.000) -- (4.000,2.000) -- (3.000,2.000) -- cycle;
\filldraw[iefill] (1.057,0.943) -- (1.943,0.943) -- (1.943,1.057) -- (1.057,1.057) -- cycle;
\draw[red, line width=1.1pt] (1.943,0.943) -- (1.943,1.057);
\draw[red, line width=1.1pt] (1.057,1.057) -- (1.057,0.943);
\draw[seledge] (2.000,1.000) -- (3.000,1.000);
\draw[hearrow, orange!80!black] (1.803,0.968) -- (1.803,1.032);
\node[helabel, text=orange!80!black] at (1.583,1.200) {$1$};
\draw[hearrow, orange!80!black] (1.780,0.860) -- (1.220,0.860);
\node[helabel, text=orange!80!black] at (1.500,0.664) {$2$};
\draw[hearrow, orange!80!black] (2.140,0.780) -- (2.140,0.220);
\node[helabel, text=orange!80!black] at (2.336,0.500) {$3$};
\draw[hearrow, red!70!black] (2.220,1.140) -- (2.780,1.140);
\node[helabel, text=red!70!black] at (2.500,1.322) {$he=\mathrm{stop}$};
\node[vreg, dot] at (0.000,0.000) {};
\node[vreg, dot] at (1.000,0.000) {};
\node[vreg, dot] at (2.000,0.000) {};
\node[vreg, dot] at (3.000,0.000) {};
\node[vreg, dot] at (4.000,0.000) {};
\node[vreg, dot] at (0.000,1.000) {};
\node[vcol, minimum size=0.29cm] at (1.000,1.000) {$A$};
\node[vcol, minimum size=0.29cm] at (2.000,1.000) {$B$};
\node[vreg, minimum size=0.29cm] at (3.000,1.000) {$C$};
\node[vreg, dot] at (4.000,1.000) {};
\node[vreg, dot] at (0.000,2.000) {};
\node[vreg, dot] at (1.000,2.000) {};
\node[vreg, dot] at (2.000,2.000) {};
\node[vreg, dot] at (3.000,2.000) {};
\node[vreg, dot] at (4.000,2.000) {};

\end{tikzpicture}
}\caption{Rotação a partir de $c$.}\end{subfigure}\hfill
\begin{subfigure}[b]{0.32\linewidth}\centering\resizebox{\linewidth}{!}{\begin{tikzpicture}[scale=1.700, every node/.style={font=\scriptsize}]
\clip (0.450,0.300) rectangle (3.550,1.700);
\filldraw[cefill] (0.000,0.000) -- (1.000,0.000) -- (1.000,1.000) -- (0.000,1.000) -- cycle;
\filldraw[cefill] (1.000,0.000) -- (2.000,0.000) -- (2.000,0.800) -- (1.000,1.000) -- cycle;
\filldraw[cefill] (2.000,0.000) -- (3.000,0.000) -- (3.000,1.000) -- (2.000,0.800) -- cycle;
\filldraw[cefill] (3.000,0.000) -- (4.000,0.000) -- (4.000,1.000) -- (3.000,1.000) -- cycle;
\filldraw[cefill] (0.000,1.000) -- (1.000,1.000) -- (1.000,2.000) -- (0.000,2.000) -- cycle;
\filldraw[cefill] (1.000,1.000) -- (2.000,1.200) -- (2.000,2.000) -- (1.000,2.000) -- cycle;
\filldraw[cefill] (2.000,1.200) -- (3.000,1.000) -- (3.000,2.000) -- (2.000,2.000) -- cycle;
\filldraw[cefill] (3.000,1.000) -- (4.000,1.000) -- (4.000,2.000) -- (3.000,2.000) -- cycle;
\filldraw[iefill] (1.054,0.941) -- (2.000,0.800) -- (2.000,1.200) -- (1.054,1.059) -- cycle;
\filldraw[iefill] (2.000,0.800) -- (2.946,0.941) -- (2.946,1.059) -- (2.000,1.200) -- cycle;
\draw[red, line width=1.1pt] (1.054,1.059) -- (1.054,0.941);
\draw[red, line width=1.1pt] (2.946,0.941) -- (2.946,1.059);
\node[vreg, dot] at (0.000,0.000) {};
\node[vreg, dot] at (1.000,0.000) {};
\node[vreg, dot] at (2.000,0.000) {};
\node[vreg, dot] at (3.000,0.000) {};
\node[vreg, dot] at (4.000,0.000) {};
\node[vreg, dot] at (0.000,1.000) {};
\node[vcol, minimum size=0.29cm] at (1.000,1.000) {$A$};
\node[vie, minimum size=0.29cm] at (2.000,1.200) {$B$};
\node[vcol, minimum size=0.29cm] at (3.000,1.000) {$C$};
\node[vreg, dot] at (4.000,1.000) {};
\node[vreg, dot] at (0.000,2.000) {};
\node[vreg, dot] at (1.000,2.000) {};
\node[vreg, dot] at (2.000,2.000) {};
\node[vreg, dot] at (3.000,2.000) {};
\node[vreg, dot] at (4.000,2.000) {};
\node[vie, minimum size=0.29cm] at (2.000,0.800) {$B_{1}$};

\end{tikzpicture}
}\caption{Depois.}\end{subfigure}
\caption{Expand Edge em $B$: a aresta $BC$ continua a trinca $AB$, e o elemento de interface em $AB$ não toca os
elementos de $BC$. (b) Ordem em que \code{reindexElementsCCW} visita os cantos: 1 é a aresta colapsada $c$ (em $I$),
2 e 3 são os elementos abaixo da trinca; a rotação para ao encontrar $he$ (a half-edge de $BC$ que parte de $B$). (c)
Os cantos visitados passam a usar $B_1$; $A$ e $C$ continuam colapsados.}
\label{fig:expand}
\end{figure}

\subsection{Open Hole}
\label{sec:hole}

\textbf{Quando:} $v$ já pertence a um elemento de interface, mas não tem aresta colapsada. Isso acontece a partir da
terceira aresta em um vértice interno: depois do Split/Expand o vértice original já virou duas cópias, $v$ e $w$, e os
dois elementos de interface $J$ e $K$ que passam por ele estão colados pelos seus lados curtos (um vai de $v$ para $w$, o
outro de $w$ para $v$).

\textbf{O que faz} (\cref{alg:hole}): com três ou mais elementos de interface chegando ao mesmo ponto, os lados curtos não
conseguem mais fechar a vizinhança, e aparece um buraco. O operador cria $n$ e faz duas travessias a partir de $he$:
\begin{enumerate}
\item no sentido \emph{anti-horário}, troca $v$ por $n$ em todos os elementos até encontrar o primeiro elemento de
interface $J$; troca também o canto de $J$ e ``descola'' o lado curto de $J$ que chega a esse canto ($\Op\gets\BORDER$);
\item no sentido \emph{horário} (regra $h\leftarrow\nxt(\Op[h])$), apenas procura o primeiro elemento de interface $K$ e
descola o lado curto de $K$ que sai de $v$.
\end{enumerate}
O lado curto do novo elemento também fica na borda. Na primeira vez, os lados descolados de $J$ e $K$ eram justamente o
par que unia as duas cópias, e o resultado é um buraco triangular $(v,n,w)$; nas vezes seguintes o buraco já existe, os
lados encontrados já são de borda e o buraco ganha um vértice (\cref{fig:hole}).

\textbf{Vértice de borda.} O Open Hole também é o operador de um vértice de borda (externa ou de furo) que já foi aberto,
porque esse vértice está em um elemento de interface e não tem aresta colapsada. A estrela de um vértice de borda vai de
borda a borda, e por isso cada laço pode terminar em uma de duas ``paredes'': um elemento de interface, como num vértice
interno, ou a própria borda da malha. Quando um laço chega à borda, não há canto de elemento de interface a mover nem
lado curto a descolar: o buraco está aberto para o exterior, e aquele lado faz o mesmo que o Canônico de borda (o setor
entre a aresta nova e a borda passa para a cópia nova). Ao menos um dos lados sempre encontra um elemento de interface,
já que o vértice pertence a um; uma asserção verifica isso. Esse caso era o defeito I1 até a v0.4
(\cref{sec:falha-borda}).

\begin{algorithm}[ht]
\caption{\code{openHole}$(he, x)$}
\label{alg:hole}
\begin{algorithmic}[1]
\State $v\gets\V[he]$; \quad $n\gets$ \code{duplicatePoint}$(v)$; \quad $h\gets he$
\Repeat \Comment{laço 1: anti-horário, reindexando}
  \State $\V[h]\gets n$; \quad $h\gets\Op[\prv(h)]$
\Until{$h=\BORDER$ \textbf{ou} \code{\_isInterfaceElement}$[\mathrm{elem}(h)]$}
\State $\mathit{achouAH}\gets(h\ne\BORDER)$
\If{$\mathit{achouAH}$} $\V[h]\gets n$; \quad $\Op[\prv(h)]\gets\BORDER$ \Comment{canto e lado curto de $J$}
\EndIf
\State $h\gets he$
\Repeat \Comment{laço 2: horário, só procura}
  \State $o\gets\Op[h]$; \quad $h\gets$ (\textbf{se} $o=\BORDER$ \textbf{então} \BORDER{} \textbf{senão} $\nxt(o)$)
\Until{$h=\BORDER$ \textbf{ou} \code{\_isInterfaceElement}$[\mathrm{elem}(h)]$}
\State $\mathit{achouH}\gets(h\ne\BORDER)$
\If{$\mathit{achouH}$} $\Op[h]\gets\BORDER$ \Comment{lado curto de $K$}
\EndIf
\State \textbf{assert}$(\mathit{achouAH}$ \textbf{ou} $\mathit{achouH})$ \Comment{o vértice está em algum EI}
\State $\V[x]\gets n$; \quad $\V[\nxt(x)]\gets v$; \quad $\Op[x]\gets\BORDER$
\State marca $v$ e $n$ como vértices de EI
\end{algorithmic}
\end{algorithm}

\begin{figure}[ht]
\centering
\begin{subfigure}[b]{0.24\linewidth}\centering\resizebox{\linewidth}{!}{\begin{tikzpicture}[scale=1.900, every node/.style={font=\scriptsize}]
\clip (0.850,0.850) rectangle (3.150,3.150);
\filldraw[cefill] (0.000,0.000) -- (1.000,0.000) -- (1.000,1.000) -- (0.000,1.000) -- cycle;
\filldraw[cefill] (1.000,0.000) -- (2.000,0.000) -- (2.000,1.000) -- (1.000,1.000) -- cycle;
\filldraw[cefill] (2.000,0.000) -- (3.000,0.000) -- (3.000,1.000) -- (2.000,1.000) -- cycle;
\filldraw[cefill] (3.000,0.000) -- (4.000,0.000) -- (4.000,1.000) -- (3.000,1.000) -- cycle;
\filldraw[cefill] (0.000,1.000) -- (1.000,1.000) -- (1.000,2.000) -- (0.000,2.000) -- cycle;
\filldraw[cefill] (1.000,1.000) -- (2.000,1.000) -- (2.000,2.000) -- (1.000,2.000) -- cycle;
\filldraw[cefill] (2.000,1.000) -- (3.000,1.000) -- (3.000,2.000) -- (2.000,2.000) -- cycle;
\filldraw[cefill] (3.000,1.000) -- (4.000,1.000) -- (4.000,2.000) -- (3.000,2.000) -- cycle;
\filldraw[cefill] (0.000,2.000) -- (1.000,2.000) -- (1.000,3.000) -- (0.000,3.000) -- cycle;
\filldraw[cefill] (1.000,2.000) -- (2.000,2.000) -- (2.000,3.000) -- (1.000,3.000) -- cycle;
\filldraw[cefill] (2.000,2.000) -- (3.000,2.000) -- (3.000,3.000) -- (2.000,3.000) -- cycle;
\filldraw[cefill] (3.000,2.000) -- (4.000,2.000) -- (4.000,3.000) -- (3.000,3.000) -- cycle;
\filldraw[cefill] (0.000,3.000) -- (1.000,3.000) -- (1.000,4.000) -- (0.000,4.000) -- cycle;
\filldraw[cefill] (1.000,3.000) -- (2.000,3.000) -- (2.000,4.000) -- (1.000,4.000) -- cycle;
\filldraw[cefill] (2.000,3.000) -- (3.000,3.000) -- (3.000,4.000) -- (2.000,4.000) -- cycle;
\filldraw[cefill] (3.000,3.000) -- (4.000,3.000) -- (4.000,4.000) -- (3.000,4.000) -- cycle;
\filldraw[iefill] (1.074,1.926) -- (1.926,1.926) -- (1.926,2.074) -- (1.074,2.074) -- cycle;
\draw[red, line width=1.1pt] (1.926,1.926) -- (1.926,2.074);
\draw[red, line width=1.1pt] (1.074,2.074) -- (1.074,1.926);
\draw[curedge] (2.000,2.000) -- (3.000,2.000);
\node[edgenum] at (2.500,2.160) {2};
\draw[seledge] (2.000,2.000) -- (2.000,3.000);
\node[edgenum] at (1.840,2.500) {3};
\draw[seledge] (2.000,1.000) -- (2.000,2.000);
\node[edgenum] at (1.840,1.500) {4};
\node[vreg, font=\tiny\bfseries, dot] at (0.000,0.000) {};
\node[vreg, font=\tiny\bfseries, dot] at (1.000,0.000) {};
\node[vreg, font=\tiny\bfseries, dot] at (2.000,0.000) {};
\node[vreg, font=\tiny\bfseries, dot] at (3.000,0.000) {};
\node[vreg, font=\tiny\bfseries, dot] at (4.000,0.000) {};
\node[vreg, font=\tiny\bfseries, dot] at (0.000,1.000) {};
\node[vreg, font=\tiny\bfseries, dot] at (1.000,1.000) {};
\node[vreg, font=\tiny\bfseries, dot] at (2.000,1.000) {};
\node[vreg, font=\tiny\bfseries, dot] at (3.000,1.000) {};
\node[vreg, font=\tiny\bfseries, dot] at (4.000,1.000) {};
\node[vreg, font=\tiny\bfseries, dot] at (0.000,2.000) {};
\node[vcol, font=\tiny\bfseries, dot] at (1.000,2.000) {};
\node[vcol, font=\tiny\bfseries, minimum size=0.24cm] at (2.000,2.000) {$V$};
\node[vreg, font=\tiny\bfseries, dot] at (3.000,2.000) {};
\node[vreg, font=\tiny\bfseries, dot] at (4.000,2.000) {};
\node[vreg, font=\tiny\bfseries, dot] at (0.000,3.000) {};
\node[vreg, font=\tiny\bfseries, dot] at (1.000,3.000) {};
\node[vreg, font=\tiny\bfseries, dot] at (2.000,3.000) {};
\node[vreg, font=\tiny\bfseries, dot] at (3.000,3.000) {};
\node[vreg, font=\tiny\bfseries, dot] at (4.000,3.000) {};
\node[vreg, font=\tiny\bfseries, dot] at (0.000,4.000) {};
\node[vreg, font=\tiny\bfseries, dot] at (1.000,4.000) {};
\node[vreg, font=\tiny\bfseries, dot] at (2.000,4.000) {};
\node[vreg, font=\tiny\bfseries, dot] at (3.000,4.000) {};
\node[vreg, font=\tiny\bfseries, dot] at (4.000,4.000) {};

\end{tikzpicture}
}\caption{Aresta 1: Canônico.}\end{subfigure}\hfill
\begin{subfigure}[b]{0.24\linewidth}\centering\resizebox{\linewidth}{!}{\begin{tikzpicture}[scale=1.900, every node/.style={font=\scriptsize}]
\clip (0.850,0.850) rectangle (3.150,3.150);
\filldraw[cefill] (0.000,0.000) -- (1.000,0.000) -- (1.000,1.000) -- (0.000,1.000) -- cycle;
\filldraw[cefill] (1.000,0.000) -- (2.000,0.000) -- (2.000,1.000) -- (1.000,1.000) -- cycle;
\filldraw[cefill] (2.000,0.000) -- (3.000,0.000) -- (3.000,1.000) -- (2.000,1.000) -- cycle;
\filldraw[cefill] (3.000,0.000) -- (4.000,0.000) -- (4.000,1.000) -- (3.000,1.000) -- cycle;
\filldraw[cefill] (0.000,1.000) -- (1.000,1.000) -- (1.000,2.000) -- (0.000,2.000) -- cycle;
\filldraw[cefill] (1.000,1.000) -- (2.000,1.000) -- (2.000,1.740) -- (1.000,2.000) -- cycle;
\filldraw[cefill] (2.000,1.000) -- (3.000,1.000) -- (3.000,2.000) -- (2.000,1.740) -- cycle;
\filldraw[cefill] (3.000,1.000) -- (4.000,1.000) -- (4.000,2.000) -- (3.000,2.000) -- cycle;
\filldraw[cefill] (0.000,2.000) -- (1.000,2.000) -- (1.000,3.000) -- (0.000,3.000) -- cycle;
\filldraw[cefill] (1.000,2.000) -- (2.000,2.260) -- (2.000,3.000) -- (1.000,3.000) -- cycle;
\filldraw[cefill] (2.000,2.260) -- (3.000,2.000) -- (3.000,3.000) -- (2.000,3.000) -- cycle;
\filldraw[cefill] (3.000,2.000) -- (4.000,2.000) -- (4.000,3.000) -- (3.000,3.000) -- cycle;
\filldraw[cefill] (0.000,3.000) -- (1.000,3.000) -- (1.000,4.000) -- (0.000,4.000) -- cycle;
\filldraw[cefill] (1.000,3.000) -- (2.000,3.000) -- (2.000,4.000) -- (1.000,4.000) -- cycle;
\filldraw[cefill] (2.000,3.000) -- (3.000,3.000) -- (3.000,4.000) -- (2.000,4.000) -- cycle;
\filldraw[cefill] (3.000,3.000) -- (4.000,3.000) -- (4.000,4.000) -- (3.000,4.000) -- cycle;
\filldraw[iefill] (1.069,1.922) -- (2.000,1.740) -- (2.000,2.260) -- (1.069,2.078) -- cycle;
\filldraw[iefill] (2.000,1.740) -- (2.931,1.922) -- (2.931,2.078) -- (2.000,2.260) -- cycle;
\draw[red, line width=1.1pt] (1.069,2.078) -- (1.069,1.922);
\draw[red, line width=1.1pt] (2.931,1.922) -- (2.931,2.078);
\draw[curedge] (2.000,2.260) -- (2.000,3.000);
\node[edgenum] at (1.840,2.630) {3};
\draw[seledge] (2.000,1.000) -- (2.000,1.740);
\node[edgenum] at (1.840,1.370) {4};
\node[vreg, font=\tiny\bfseries, dot] at (0.000,0.000) {};
\node[vreg, font=\tiny\bfseries, dot] at (1.000,0.000) {};
\node[vreg, font=\tiny\bfseries, dot] at (2.000,0.000) {};
\node[vreg, font=\tiny\bfseries, dot] at (3.000,0.000) {};
\node[vreg, font=\tiny\bfseries, dot] at (4.000,0.000) {};
\node[vreg, font=\tiny\bfseries, dot] at (0.000,1.000) {};
\node[vreg, font=\tiny\bfseries, dot] at (1.000,1.000) {};
\node[vreg, font=\tiny\bfseries, dot] at (2.000,1.000) {};
\node[vreg, font=\tiny\bfseries, dot] at (3.000,1.000) {};
\node[vreg, font=\tiny\bfseries, dot] at (4.000,1.000) {};
\node[vreg, font=\tiny\bfseries, dot] at (0.000,2.000) {};
\node[vcol, font=\tiny\bfseries, dot] at (1.000,2.000) {};
\node[vie, font=\tiny\bfseries, minimum size=0.24cm] at (2.000,2.260) {$V$};
\node[vcol, font=\tiny\bfseries, dot] at (3.000,2.000) {};
\node[vreg, font=\tiny\bfseries, dot] at (4.000,2.000) {};
\node[vreg, font=\tiny\bfseries, dot] at (0.000,3.000) {};
\node[vreg, font=\tiny\bfseries, dot] at (1.000,3.000) {};
\node[vreg, font=\tiny\bfseries, dot] at (2.000,3.000) {};
\node[vreg, font=\tiny\bfseries, dot] at (3.000,3.000) {};
\node[vreg, font=\tiny\bfseries, dot] at (4.000,3.000) {};
\node[vreg, font=\tiny\bfseries, dot] at (0.000,4.000) {};
\node[vreg, font=\tiny\bfseries, dot] at (1.000,4.000) {};
\node[vreg, font=\tiny\bfseries, dot] at (2.000,4.000) {};
\node[vreg, font=\tiny\bfseries, dot] at (3.000,4.000) {};
\node[vreg, font=\tiny\bfseries, dot] at (4.000,4.000) {};
\node[vie, font=\tiny\bfseries, minimum size=0.24cm] at (2.000,1.740) {$V_{1}$};

\end{tikzpicture}
}\caption{Aresta 2: Expand.}\end{subfigure}\hfill
\begin{subfigure}[b]{0.24\linewidth}\centering\resizebox{\linewidth}{!}{\begin{tikzpicture}[scale=1.900, every node/.style={font=\scriptsize}]
\clip (0.850,0.850) rectangle (3.150,3.150);
\fill[holefill] (2.000,1.740) -- (1.816,2.184) -- (2.184,2.184) -- cycle;
\filldraw[cefill] (0.000,0.000) -- (1.000,0.000) -- (1.000,1.000) -- (0.000,1.000) -- cycle;
\filldraw[cefill] (1.000,0.000) -- (2.000,0.000) -- (2.000,1.000) -- (1.000,1.000) -- cycle;
\filldraw[cefill] (2.000,0.000) -- (3.000,0.000) -- (3.000,1.000) -- (2.000,1.000) -- cycle;
\filldraw[cefill] (3.000,0.000) -- (4.000,0.000) -- (4.000,1.000) -- (3.000,1.000) -- cycle;
\filldraw[cefill] (0.000,1.000) -- (1.000,1.000) -- (1.000,2.000) -- (0.000,2.000) -- cycle;
\filldraw[cefill] (1.000,1.000) -- (2.000,1.000) -- (2.000,1.740) -- (1.000,2.000) -- cycle;
\filldraw[cefill] (2.000,1.000) -- (3.000,1.000) -- (3.000,2.000) -- (2.000,1.740) -- cycle;
\filldraw[cefill] (3.000,1.000) -- (4.000,1.000) -- (4.000,2.000) -- (3.000,2.000) -- cycle;
\filldraw[cefill] (0.000,2.000) -- (1.000,2.000) -- (1.000,3.000) -- (0.000,3.000) -- cycle;
\filldraw[cefill] (1.000,2.000) -- (1.816,2.184) -- (2.000,3.000) -- (1.000,3.000) -- cycle;
\filldraw[cefill] (2.184,2.184) -- (3.000,2.000) -- (3.000,3.000) -- (2.000,3.000) -- cycle;
\filldraw[cefill] (3.000,2.000) -- (4.000,2.000) -- (4.000,3.000) -- (3.000,3.000) -- cycle;
\filldraw[cefill] (0.000,3.000) -- (1.000,3.000) -- (1.000,4.000) -- (0.000,4.000) -- cycle;
\filldraw[cefill] (1.000,3.000) -- (2.000,3.000) -- (2.000,4.000) -- (1.000,4.000) -- cycle;
\filldraw[cefill] (2.000,3.000) -- (3.000,3.000) -- (3.000,4.000) -- (2.000,4.000) -- cycle;
\filldraw[cefill] (3.000,3.000) -- (4.000,3.000) -- (4.000,4.000) -- (3.000,4.000) -- cycle;
\filldraw[iefill] (1.069,1.922) -- (2.000,1.740) -- (1.816,2.184) -- (1.066,2.080) -- cycle;
\filldraw[iefill] (2.000,1.740) -- (2.931,1.922) -- (2.934,2.080) -- (2.184,2.184) -- cycle;
\filldraw[iefill] (2.184,2.184) -- (2.080,2.934) -- (1.920,2.934) -- (1.816,2.184) -- cycle;
\draw[red, line width=1.1pt] (1.066,2.080) -- (1.069,1.922);
\draw[red, line width=1.1pt] (2.931,1.922) -- (2.934,2.080);
\draw[red, line width=1.1pt] (2.080,2.934) -- (1.920,2.934);
\draw[curedge] (2.000,1.000) -- (2.000,1.740);
\node[edgenum] at (1.840,1.370) {4};
\node[vreg, font=\tiny\bfseries, dot] at (0.000,0.000) {};
\node[vreg, font=\tiny\bfseries, dot] at (1.000,0.000) {};
\node[vreg, font=\tiny\bfseries, dot] at (2.000,0.000) {};
\node[vreg, font=\tiny\bfseries, dot] at (3.000,0.000) {};
\node[vreg, font=\tiny\bfseries, dot] at (4.000,0.000) {};
\node[vreg, font=\tiny\bfseries, dot] at (0.000,1.000) {};
\node[vreg, font=\tiny\bfseries, dot] at (1.000,1.000) {};
\node[vreg, font=\tiny\bfseries, dot] at (2.000,1.000) {};
\node[vreg, font=\tiny\bfseries, dot] at (3.000,1.000) {};
\node[vreg, font=\tiny\bfseries, dot] at (4.000,1.000) {};
\node[vreg, font=\tiny\bfseries, dot] at (0.000,2.000) {};
\node[vcol, font=\tiny\bfseries, dot] at (1.000,2.000) {};
\node[vie, font=\tiny\bfseries, minimum size=0.24cm] at (2.184,2.184) {$V$};
\node[vcol, font=\tiny\bfseries, dot] at (3.000,2.000) {};
\node[vreg, font=\tiny\bfseries, dot] at (4.000,2.000) {};
\node[vreg, font=\tiny\bfseries, dot] at (0.000,3.000) {};
\node[vreg, font=\tiny\bfseries, dot] at (1.000,3.000) {};
\node[vcol, font=\tiny\bfseries, dot] at (2.000,3.000) {};
\node[vreg, font=\tiny\bfseries, dot] at (3.000,3.000) {};
\node[vreg, font=\tiny\bfseries, dot] at (4.000,3.000) {};
\node[vreg, font=\tiny\bfseries, dot] at (0.000,4.000) {};
\node[vreg, font=\tiny\bfseries, dot] at (1.000,4.000) {};
\node[vreg, font=\tiny\bfseries, dot] at (2.000,4.000) {};
\node[vreg, font=\tiny\bfseries, dot] at (3.000,4.000) {};
\node[vreg, font=\tiny\bfseries, dot] at (4.000,4.000) {};
\node[vie, font=\tiny\bfseries, minimum size=0.24cm] at (2.000,1.740) {$V_{1}$};
\node[vie, font=\tiny\bfseries, minimum size=0.24cm] at (1.816,2.184) {$V_{2}$};

\end{tikzpicture}
}\caption{Aresta 3: Open Hole.}\end{subfigure}\hfill
\begin{subfigure}[b]{0.24\linewidth}\centering\resizebox{\linewidth}{!}{\begin{tikzpicture}[scale=1.900, every node/.style={font=\scriptsize}]
\clip (0.850,0.850) rectangle (3.150,3.150);
\fill[holefill] (1.816,1.816) -- (1.816,2.184) -- (2.184,2.184) -- (2.184,1.816) -- cycle;
\filldraw[cefill] (0.000,0.000) -- (1.000,0.000) -- (1.000,1.000) -- (0.000,1.000) -- cycle;
\filldraw[cefill] (1.000,0.000) -- (2.000,0.000) -- (2.000,1.000) -- (1.000,1.000) -- cycle;
\filldraw[cefill] (2.000,0.000) -- (3.000,0.000) -- (3.000,1.000) -- (2.000,1.000) -- cycle;
\filldraw[cefill] (3.000,0.000) -- (4.000,0.000) -- (4.000,1.000) -- (3.000,1.000) -- cycle;
\filldraw[cefill] (0.000,1.000) -- (1.000,1.000) -- (1.000,2.000) -- (0.000,2.000) -- cycle;
\filldraw[cefill] (1.000,1.000) -- (2.000,1.000) -- (1.816,1.816) -- (1.000,2.000) -- cycle;
\filldraw[cefill] (2.000,1.000) -- (3.000,1.000) -- (3.000,2.000) -- (2.184,1.816) -- cycle;
\filldraw[cefill] (3.000,1.000) -- (4.000,1.000) -- (4.000,2.000) -- (3.000,2.000) -- cycle;
\filldraw[cefill] (0.000,2.000) -- (1.000,2.000) -- (1.000,3.000) -- (0.000,3.000) -- cycle;
\filldraw[cefill] (1.000,2.000) -- (1.816,2.184) -- (2.000,3.000) -- (1.000,3.000) -- cycle;
\filldraw[cefill] (2.184,2.184) -- (3.000,2.000) -- (3.000,3.000) -- (2.000,3.000) -- cycle;
\filldraw[cefill] (3.000,2.000) -- (4.000,2.000) -- (4.000,3.000) -- (3.000,3.000) -- cycle;
\filldraw[cefill] (0.000,3.000) -- (1.000,3.000) -- (1.000,4.000) -- (0.000,4.000) -- cycle;
\filldraw[cefill] (1.000,3.000) -- (2.000,3.000) -- (2.000,4.000) -- (1.000,4.000) -- cycle;
\filldraw[cefill] (2.000,3.000) -- (3.000,3.000) -- (3.000,4.000) -- (2.000,4.000) -- cycle;
\filldraw[cefill] (3.000,3.000) -- (4.000,3.000) -- (4.000,4.000) -- (3.000,4.000) -- cycle;
\filldraw[iefill] (1.066,1.920) -- (1.816,1.816) -- (1.816,2.184) -- (1.066,2.080) -- cycle;
\filldraw[iefill] (2.184,1.816) -- (2.934,1.920) -- (2.934,2.080) -- (2.184,2.184) -- cycle;
\filldraw[iefill] (2.184,2.184) -- (2.080,2.934) -- (1.920,2.934) -- (1.816,2.184) -- cycle;
\filldraw[iefill] (2.080,1.066) -- (2.184,1.816) -- (1.816,1.816) -- (1.920,1.066) -- cycle;
\draw[red, line width=1.1pt] (1.066,2.080) -- (1.066,1.920);
\draw[red, line width=1.1pt] (2.934,1.920) -- (2.934,2.080);
\draw[red, line width=1.1pt] (2.080,2.934) -- (1.920,2.934);
\draw[red, line width=1.1pt] (1.920,1.066) -- (2.080,1.066);
\node[vreg, font=\tiny\bfseries, dot] at (0.000,0.000) {};
\node[vreg, font=\tiny\bfseries, dot] at (1.000,0.000) {};
\node[vreg, font=\tiny\bfseries, dot] at (2.000,0.000) {};
\node[vreg, font=\tiny\bfseries, dot] at (3.000,0.000) {};
\node[vreg, font=\tiny\bfseries, dot] at (4.000,0.000) {};
\node[vreg, font=\tiny\bfseries, dot] at (0.000,1.000) {};
\node[vreg, font=\tiny\bfseries, dot] at (1.000,1.000) {};
\node[vcol, font=\tiny\bfseries, dot] at (2.000,1.000) {};
\node[vreg, font=\tiny\bfseries, dot] at (3.000,1.000) {};
\node[vreg, font=\tiny\bfseries, dot] at (4.000,1.000) {};
\node[vreg, font=\tiny\bfseries, dot] at (0.000,2.000) {};
\node[vcol, font=\tiny\bfseries, dot] at (1.000,2.000) {};
\node[vie, font=\tiny\bfseries, minimum size=0.24cm] at (2.184,2.184) {$V$};
\node[vcol, font=\tiny\bfseries, dot] at (3.000,2.000) {};
\node[vreg, font=\tiny\bfseries, dot] at (4.000,2.000) {};
\node[vreg, font=\tiny\bfseries, dot] at (0.000,3.000) {};
\node[vreg, font=\tiny\bfseries, dot] at (1.000,3.000) {};
\node[vcol, font=\tiny\bfseries, dot] at (2.000,3.000) {};
\node[vreg, font=\tiny\bfseries, dot] at (3.000,3.000) {};
\node[vreg, font=\tiny\bfseries, dot] at (4.000,3.000) {};
\node[vreg, font=\tiny\bfseries, dot] at (0.000,4.000) {};
\node[vreg, font=\tiny\bfseries, dot] at (1.000,4.000) {};
\node[vreg, font=\tiny\bfseries, dot] at (2.000,4.000) {};
\node[vreg, font=\tiny\bfseries, dot] at (3.000,4.000) {};
\node[vreg, font=\tiny\bfseries, dot] at (4.000,4.000) {};
\node[vie, font=\tiny\bfseries, minimum size=0.24cm] at (1.816,1.816) {$V_{1}$};
\node[vie, font=\tiny\bfseries, minimum size=0.24cm] at (1.816,2.184) {$V_{2}$};
\node[vie, font=\tiny\bfseries, minimum size=0.24cm] at (2.184,1.816) {$V_{3}$};

\end{tikzpicture}
}\caption{Aresta 4: Open Hole.}\end{subfigure}
\caption{Quatro arestas chegando ao vértice interno $V$ (detalhe do centro de uma malha $4\times4$; a próxima aresta a
inserir está em vermelho grosso). (a) A primeira deixa $V$ colapsado. (b) A segunda, alinhada, divide $V$ em duas cópias
coladas pelos lados curtos. (c) A terceira descola esses lados: buraco triangular (hachurado). (d) A quarta aumenta o
buraco para um quadrilátero. Com $k$ arestas, o buraco tem $k$ vértices.}
\label{fig:hole}
\end{figure}

\subsection{Relações entre os operadores}
\label{sec:relacoes}

Olhando os quatro operadores lado a lado, todos fazem a mesma coisa em essência: \emph{dividem um leque da estrela de $v$
entre duas ``paredes'', dão um vértice novo a um dos lados e decidem o que fazer com os lados curtos}. O que muda é a
parede onde a rotação começa e onde ela para:

\begin{center}\small
\begin{tabular}{llll}
\toprule
Operador & começa em & para em & lados curtos\\
\midrule
Canônico (borda) & aresta nova & borda da malha & novo lado na borda\\
Expand Edge & aresta colapsada $c$ & aresta nova & $c$ é colada ao novo lado\\
Split Element & aresta colapsada $c$ & aresta nova (após 1 elemento) & $c$ é colada ao novo lado\\
Open Hole & aresta nova & primeiro elemento de interface (ou a borda) & descolados (buraco)\\
\bottomrule
\end{tabular}
\end{center}

Uma consequência interessante: o \textbf{Split Element é um caso particular do Expand Edge}. Na hipótese $x_3$, a rotação
do Expand a partir de $c$ visita exatamente o canto de $c$ e o canto $\nxt(b)$ de $S$ e para em $he$ --- as mesmas
alterações do Split. Na hipótese $x_1$, a rotação do Expand vai pelo outro lado e reindexa todos os elementos da estrela
\emph{exceto} $S$; o resultado é a mesma divisão da estrela, só com os nomes $v$ e $n$ trocados. Verificamos isso
experimentalmente: substituindo o Split pelo Expand em 791 ocorrências (400 malhas aleatórias de triângulos,
quadriláteros, híbridas e com polígonos), a topologia final foi idêntica em todos os casos (\cref{sec:verificacao}).
Desde a v0.5 essa equivalência é um teste automatizado, em malhas aleatórias dos cinco tipos (\cref{sec:testes}). O
Split Element é, portanto, uma \emph{otimização}: custo $O(\log n_c)$ em vez de $O(d\log n_c)$, movendo sempre o lado
menor. Isso não é um defeito; mas é uma informação útil para a extensão 3D,
onde talvez valha começar pela versão geral (rotação) e só depois especializar.

\subsection{Contagem prévia de vértices}
\label{sec:contagem}

Para reservar memória de uma vez, \code{computeNumberOfNewVertices} calcula, antes de alterar a malha, quantos vértices
a chamada vai criar. A contagem tem duas etapas:
\begin{enumerate}
\item \textbf{Seleção.} Percorre as arestas do lote e descarta as mesmas que \code{insertInterfaceElements} vai
rejeitar, com os mesmos testes e na mesma ordem: índice que não é half-edge da malha ($h\ge Of[n_c]$), aresta de borda,
aresta de elemento de interface e aresta repetida no lote (inclusive quando vem uma vez por cada uma de suas
half-edges; a aresta é identificada pela menor das duas). Para cada vértice tocado pelas arestas que sobram, guarda o
número $n_v$ dessas arestas e uma half-edge que parte dele.
\item \textbf{Regra por vértice.} Aplica, uma vez por vértice,
\[
\text{novos}(v)=
\begin{cases}
n_v-1, & v \text{ fora de EI e interno (o primeiro contato só colapsa)},\\
n_v, & v \text{ fora de EI e na borda (o primeiro contato já duplica)},\\
n_v, & v \text{ já em algum EI (todo contato duplica)}.
\end{cases}
\]
\end{enumerate}
Como a regra é aplicada por vértice, e não por aresta, o teste de borda (\code{isBorder}, que gira em torno do vértice)
roda uma vez por vértice distinto, e um vértice tocado por várias arestas do lote é contado corretamente. Por descartar
as arestas de elementos de interface antes do teste de borda, a contagem também deixou de chegar ao caminho em que
\code{isBorder} não termina (I2). O resultado é exatamente o número de vértices criados, em todos os testes
(\cref{sec:verificacao}); até a v0.3 a contagem era um limite superior (defeito I5, corrigido). Num vértice de borda com
duas ou mais arestas, ela dá o valor correto ($n_v$), uma cópia por aresta (Open Hole na borda). Os elementos
são reservados à razão de um por aresta do lote.

\section{O algoritmo de travessia}
\label{sec:travessia}

\subsection{Inserção de um lote de arestas}

O ponto de entrada é \code{insertInterfaceElements}, que recebe uma lista de half-edges (uma por aresta; qualquer uma
das duas serve) e as processa na ordem dada (\cref{alg:insert}). Antes de alterar a malha, o método conta os vértices
novos e reserva espaço para eles e para um elemento de quatro half-edges por aresta. Depois, para cada aresta, valida a
entrada, abre o novo elemento a partir do gabarito da \cref{fig:gabarito}, escolhe e aplica o operador em $A$ e em $B$
e, por fim, liga as arestas do meio. A \cref{sec:ibhm-insercao} mostra o efeito de cada passo nos vetores da IBHM.

\begin{algorithm}[ht]
\caption{\code{insertInterfaceElements}$(E)$ \quad (classe \code{IBHMInterfaceElementBuilder})}
\label{alg:insert}
\begin{algorithmic}[1]
\Require lista $E$ de half-edges, uma por aresta a abrir
\State $N\gets$ \code{numberOfHalfEdges}$()$ \Comment{$Of[n_c]$ antes da chamada}
\State $n_v\gets$ \code{computeNumberOfNewVertices}$(E)$
\State \code{reserveSpaceForElements}$(4|E|, |E|)$; \code{reserveSpaceForNodes}$(n_v)$; aumenta
\code{\_isInterfaceElement} e \code{\_inInterfaceElement}
\ForAll{$a\in E$}
  \If{$a\ge N$} \textbf{continue} \Comment{não é half-edge da malha: rejeitada}
  \EndIf
  \State $b\gets\Op[a]$
  \If{$b=\BORDER$ \textbf{ou} $b=\COLLAPSED$} \textbf{continue} \Comment{aresta de borda: rejeitada}
  \EndIf
  \If{$\mathrm{elem}(a)$ \textbf{ou} $\mathrm{elem}(b)$ é elemento de interface} \textbf{continue} \Comment{rejeitada}
  \EndIf
  \State $h_0\gets$ \code{nextAvailableHE}$()$; \quad $h_1,h_2,h_3\gets h_0+1,h_0+2,h_0+3$
  \State $e\gets$ \code{commitElement}$(4)$ \Comment{antes dos operadores: $\nxt/\prv$ dependem de $Of[\,]$}
  \State $\mathit{op}\gets$ \code{retrieveOperator}$(a)$; aplica $\mathit{op}$ com $(a,h_3)$ \Comment{vértice $A$}
  \State $\mathit{op}\gets$ \code{retrieveOperator}$(b)$; aplica $\mathit{op}$ com $(b,h_1)$ \Comment{vértice $B$}
  \State $\Op[h_0]\leftrightarrow b$; \quad $\Op[h_2]\leftrightarrow a$ \Comment{arestas do meio}
  \State \code{\_isInterfaceElement}$[e]\gets$ \textbf{true}
\EndFor
\end{algorithmic}
\end{algorithm}

Alguns detalhes de implementação merecem atenção:
\begin{itemize}
\item \textbf{Rejeições.} Índices que não são half-edges da malha, arestas de borda e arestas que já pertencem a um
elemento de interface (por exemplo, a mesma aresta pedida duas vezes, no mesmo lote ou em chamadas diferentes) são
ignorados com uma mensagem no console. Nada é devolvido a quem chamou (I7), e o espaço reservado para eles fica sem uso.
O teste do índice vem primeiro, porque as consultas seguintes só valem para half-edges de elementos validados. O limite
$N$ é lido no início da chamada: os elementos criados pela própria chamada não contam, e uma half-edge deles na entrada
também é rejeitada. A pré-contagem aplica os mesmos testes (\cref{sec:contagem}), por isso a contagem é exata.
\item \textbf{Validação antecipada do elemento.} \code{commitElement(4)} grava o offset do novo elemento antes de os
operadores rodarem, porque na IBHM $\nxt$ e $\prv$ das half-edges $h_0,\dots,h_3$ dependem dele. Durante os operadores,
as posições $\Op[h_0]$ e $\Op[h_2]$ ainda contêm o \BORDER{} da reserva; nenhuma rotação chega a elas porque nenhuma
half-edge antiga aponta para o novo elemento até o fim da iteração.
\item \textbf{Ordem $A$, depois $B$.} O operador de $B$ é escolhido depois de o operador de $A$ ter sido aplicado. Como
os operadores só mexem na estrela do próprio vértice, isso não altera a decisão em $B$.
\item \textbf{$a$ e $b$ ainda opostas.} Enquanto os operadores trabalham, $a$ e $b$ continuam opostas entre si; é isso
que permite que as rotações em torno de $A$ e de $B$ atravessem a aresta que está sendo aberta (por exemplo, o laço
horário do Open Hole parte de $he$ e passa por $\Op[he]$). Só no final $a$ e $b$ são religadas a $h_2$ e $h_0$.
\item \textbf{Lote ou uma a uma.} Chamar o método uma vez com todas as arestas ou uma vez por aresta produz a mesma
topologia (\cref{sec:verificacao}); a diferença é só a quantidade de reservas de memória.
\end{itemize}

\subsection{O exemplo completo, passo a passo}
\label{sec:exemplo}

A \cref{tab:exemplo} lista o operador escolhido pelo código em cada extremidade, para as 12 arestas do exemplo da
apresentação, na ordem dos slides. Os 24 operadores coincidem exatamente com os anunciados nos slides 15 a 26 da
apresentação~\cite{slides}. A \cref{fig:passos} mostra o estado da malha depois de cada inserção, e a
\cref{fig:final} o resultado final com todos os vértices nomeados.

\begin{table}[ht]
\centering\small
\begin{tabular}{clllc}
\toprule
Aresta & $A$--$B$ (no código) & Operador em $A$ & Operador em $B$ & vértices novos\\
\midrule
1 & $K$--$Q$ & Canônico & Canônico & 1 \\
2 & $K$--$L$ & Split Element & Canônico & 1 \\
3 & $G$--$C$ & Canônico & Canônico & 1 \\
4 & $M$--$H$ & Canônico & Canônico & 1 \\
5 & $S$--$M$ & Canônico & Expand Edge & 2 \\
6 & $L$--$M$ & Expand Edge & Open Hole & 2 \\
7 & $L_{1}$--$G$ & Open Hole & Expand Edge & 2 \\
8 & $K$--$F$ & Open Hole & Canônico & 1 \\
9 & $J$--$K$ & Canônico & Open Hole & 2 \\
10 & $M_{1}$--$N$ & Open Hole & Canônico & 2 \\
11 & $M_{1}$--$I$ & Open Hole & Canônico & 2 \\
12 & $M_{3}$--$T$ & Open Hole & Canônico & 2 \\
\midrule
\multicolumn{4}{r}{total} & 19\\
\bottomrule
\end{tabular}

\caption{Operadores escolhidos pelo código no exemplo (IBHM, malha \code{malha2Hyb.txt}). $X_i$ é a $i$-ésima cópia
criada do vértice $X$. A coluna $A$--$B$ mostra a orientação da half-edge usada (por isso a aresta 1, ``$QK$'' nos
slides, aparece como $K$--$Q$).}
\label{tab:exemplo}
\end{table}

\begin{figure}[p]
\centering
\newcommand{\passo}[2]{\begin{subfigure}[t]{0.32\linewidth}\centering\input{figs/example_step#1.tex}\caption*{\scriptsize #2}\end{subfigure}}
\passo{01}{1: $K$ Canônico, $Q$ Canônico (borda)}\hfill
\passo{02}{2: $K$ Split, $L$ Canônico}\hfill
\passo{03}{3: $G$ Canônico, $C$ Canônico (borda)}\\[3pt]
\passo{04}{4: $M$ Canônico, $H$ Canônico (borda)}\hfill
\passo{05}{5: $S$ Canônico (borda), $M$ Expand}\hfill
\passo{06}{6: $L$ Expand, $M$ Open Hole}\\[3pt]
\passo{07}{7: $L$ Open Hole, $G$ Expand}\hfill
\passo{08}{8: $K$ Open Hole, $F$ Canônico}\hfill
\passo{09}{9: $J$ Canônico (borda), $K$ Open Hole}\\[3pt]
\passo{10}{10: $M$ Open Hole, $N$ Canônico (borda)}\hfill
\passo{11}{11: $M$ Open Hole, $I$ Canônico (borda)}\hfill
\passo{12}{12: $M$ Open Hole, $T$ Canônico (borda)}
\caption{Estado da malha após cada uma das 12 inserções do exemplo (vista explodida). As arestas ainda não inseridas
estão em vermelho e a próxima, em vermelho grosso. Observe a sequência em $M$: colapsado (4), dividido por Expand (5),
buraco triangular (6) que cresce até um hexágono (12). Em $K$ a segunda aresta já chega por Split Element (2), porque
$KL$ e $KQ$ compartilham o elemento $KLRQ$.}
\label{fig:passos}
\end{figure}

\begin{figure}[ht]
\centering
\begin{tikzpicture}[scale=3.200, every node/.style={font=\scriptsize}]
\fill[holefill] (0.880,2.120) -- (1.120,2.120) -- (1.120,1.880) -- (0.880,1.880) -- cycle;
\fill[holefill] (1.880,1.880) -- (2.000,2.170) -- (2.120,1.880) -- cycle;
\fill[holefill] (3.076,1.848) -- (2.880,1.880) -- (2.880,2.120) -- (3.076,2.152) -- (3.152,2.076) -- (3.152,1.924) -- cycle;
\filldraw[cefill] (0.000,0.000) -- (1.000,0.000) -- (1.000,1.000) -- (0.000,1.000) -- cycle;
\filldraw[cefill] (1.000,0.000) -- (1.880,0.120) -- (1.830,1.000) -- (1.000,1.000) -- cycle;
\filldraw[cefill] (2.120,0.120) -- (3.000,0.000) -- (2.830,1.000) -- (2.170,1.000) -- cycle;
\filldraw[cefill] (0.000,1.000) -- (1.000,1.000) -- (0.880,1.880) -- (0.120,1.880) -- cycle;
\filldraw[cefill] (1.000,1.000) -- (1.830,1.000) -- (1.880,1.880) -- (1.120,1.880) -- cycle;
\filldraw[cefill] (2.170,1.000) -- (2.830,1.000) -- (2.880,1.880) -- (2.120,1.880) -- cycle;
\filldraw[cefill] (3.120,1.120) -- (3.848,1.076) -- (3.076,1.848) -- cycle;
\filldraw[cefill] (3.924,1.152) -- (3.880,1.880) -- (3.152,1.924) -- cycle;
\filldraw[cefill] (0.120,2.120) -- (0.880,2.120) -- (0.880,2.880) -- (0.000,3.000) -- cycle;
\filldraw[cefill] (1.120,2.120) -- (2.000,2.170) -- (2.000,3.000) -- (1.120,2.880) -- cycle;
\filldraw[cefill] (2.000,2.170) -- (2.880,2.120) -- (2.880,2.880) -- (2.000,3.000) -- cycle;
\filldraw[cefill] (3.152,2.076) -- (3.880,2.120) -- (3.924,2.848) -- cycle;
\filldraw[cefill] (3.076,2.152) -- (3.848,2.924) -- (3.120,2.880) -- cycle;
\filldraw[iefill] (1.120,2.120) -- (1.120,2.880) -- (0.880,2.880) -- (0.880,2.120) -- cycle;
\filldraw[iefill] (1.120,1.880) -- (1.880,1.880) -- (2.000,2.170) -- (1.120,2.120) -- cycle;
\filldraw[iefill] (1.830,1.000) -- (1.880,0.120) -- (2.120,0.120) -- (2.170,1.000) -- cycle;
\filldraw[iefill] (2.880,1.880) -- (2.830,1.000) -- (3.120,1.120) -- (3.076,1.848) -- cycle;
\filldraw[iefill] (2.880,2.880) -- (2.880,2.120) -- (3.076,2.152) -- (3.120,2.880) -- cycle;
\filldraw[iefill] (2.120,1.880) -- (2.880,1.880) -- (2.880,2.120) -- (2.000,2.170) -- cycle;
\filldraw[iefill] (1.880,1.880) -- (1.830,1.000) -- (2.170,1.000) -- (2.120,1.880) -- cycle;
\filldraw[iefill] (0.880,1.880) -- (0.947,1.046) -- (1.050,1.049) -- (1.120,1.880) -- cycle;
\filldraw[iefill] (0.120,1.880) -- (0.880,1.880) -- (0.880,2.120) -- (0.120,2.120) -- cycle;
\filldraw[iefill] (3.152,1.924) -- (3.880,1.880) -- (3.880,2.120) -- (3.152,2.076) -- cycle;
\filldraw[iefill] (3.076,1.848) -- (3.848,1.076) -- (3.924,1.152) -- (3.152,1.924) -- cycle;
\filldraw[iefill] (3.152,2.076) -- (3.924,2.848) -- (3.848,2.924) -- (3.076,2.152) -- cycle;
\draw[red, line width=1.1pt] (0.947,1.046) -- (1.050,1.049);
\node[vreg, font=\tiny\bfseries, minimum size=0.29cm] at (0.000,0.000) {$A$};
\node[vreg, font=\tiny\bfseries, minimum size=0.29cm] at (1.000,0.000) {$B$};
\node[vie, font=\tiny\bfseries, minimum size=0.29cm] at (2.120,0.120) {$C$};
\node[vreg, font=\tiny\bfseries, minimum size=0.29cm] at (3.000,0.000) {$D$};
\node[vreg, font=\tiny\bfseries, minimum size=0.29cm] at (0.000,1.000) {$E$};
\node[vcol, font=\tiny\bfseries, minimum size=0.29cm] at (1.000,1.000) {$F$};
\node[vie, font=\tiny\bfseries, minimum size=0.29cm] at (1.830,1.000) {$G$};
\node[vie, font=\tiny\bfseries, minimum size=0.29cm] at (3.120,1.120) {$H$};
\node[vie, font=\tiny\bfseries, minimum size=0.29cm] at (3.924,1.152) {$I$};
\node[vie, font=\tiny\bfseries, minimum size=0.29cm] at (0.120,1.880) {$J$};
\node[vie, font=\tiny\bfseries, minimum size=0.29cm] at (0.880,2.120) {$K$};
\node[vie, font=\tiny\bfseries, minimum size=0.29cm] at (2.000,2.170) {$L$};
\node[vie, font=\tiny\bfseries, minimum size=0.29cm] at (2.880,2.120) {$M$};
\node[vie, font=\tiny\bfseries, minimum size=0.29cm] at (3.880,2.120) {$N$};
\node[vreg, font=\tiny\bfseries, minimum size=0.29cm] at (0.000,3.000) {$P$};
\node[vie, font=\tiny\bfseries, minimum size=0.29cm] at (0.880,2.880) {$Q$};
\node[vreg, font=\tiny\bfseries, minimum size=0.29cm] at (2.000,3.000) {$R$};
\node[vie, font=\tiny\bfseries, minimum size=0.29cm] at (2.880,2.880) {$S$};
\node[vie, font=\tiny\bfseries, minimum size=0.29cm] at (3.848,2.924) {$T$};
\node[vie, font=\tiny\bfseries, minimum size=0.29cm] at (1.120,2.880) {$Q_{1}$};
\node[vie, font=\tiny\bfseries, minimum size=0.29cm] at (1.120,2.120) {$K_{1}$};
\node[vie, font=\tiny\bfseries, minimum size=0.29cm] at (1.880,0.120) {$C_{1}$};
\node[vie, font=\tiny\bfseries, minimum size=0.29cm] at (2.830,1.000) {$H_{1}$};
\node[vie, font=\tiny\bfseries, minimum size=0.29cm] at (3.120,2.880) {$S_{1}$};
\node[vie, font=\tiny\bfseries, minimum size=0.29cm] at (3.076,1.848) {$M_{1}$};
\node[vie, font=\tiny\bfseries, minimum size=0.29cm] at (1.880,1.880) {$L_{1}$};
\node[vie, font=\tiny\bfseries, minimum size=0.29cm] at (2.880,1.880) {$M_{2}$};
\node[vie, font=\tiny\bfseries, minimum size=0.29cm] at (2.120,1.880) {$L_{2}$};
\node[vie, font=\tiny\bfseries, minimum size=0.29cm] at (2.170,1.000) {$G_{1}$};
\node[vie, font=\tiny\bfseries, minimum size=0.29cm] at (1.120,1.880) {$K_{2}$};
\node[vie, font=\tiny\bfseries, minimum size=0.29cm] at (0.120,2.120) {$J_{1}$};
\node[vie, font=\tiny\bfseries, minimum size=0.29cm] at (0.880,1.880) {$K_{3}$};
\node[vie, font=\tiny\bfseries, minimum size=0.29cm] at (3.152,2.076) {$M_{3}$};
\node[vie, font=\tiny\bfseries, minimum size=0.29cm] at (3.880,1.880) {$N_{1}$};
\node[vie, font=\tiny\bfseries, minimum size=0.29cm] at (3.152,1.924) {$M_{4}$};
\node[vie, font=\tiny\bfseries, minimum size=0.29cm] at (3.848,1.076) {$I_{1}$};
\node[vie, font=\tiny\bfseries, minimum size=0.29cm] at (3.076,2.152) {$M_{5}$};
\node[vie, font=\tiny\bfseries, minimum size=0.29cm] at (3.924,2.848) {$T_{1}$};

\end{tikzpicture}

\caption{Resultado final do exemplo, com os nomes de todos os vértices. Os vértices $K$, $L$ e $M$ terminam com 4, 3 e 6
cópias e com buracos de 4, 3 e 6 lados; $G$ (duas arestas) termina com duas cópias coladas; $F$ (uma aresta, interno)
continua com uma aresta colapsada (contorno vermelho); os vértices de borda $C$, $H$, $I$, $J$, $N$, $Q$, $S$ e $T$
foram duplicados uma vez. A posição das cópias é apenas ilustrativa: no código cada cópia tem exatamente as
coordenadas do seu vértice de origem.}
\label{fig:final}
\end{figure}

\subsection{Propriedades observadas}

\paragraph{Independência da ordem.} O resultado final não depende da ordem de inserção, da half-edge escolhida para
cada aresta, nem de as arestas virem em lote ou uma a uma. A razão é que a topologia final é determinada pelo
\emph{conjunto} de arestas: as arestas selecionadas cortam a estrela de cada vértice em setores, e cada setor termina com
uma cópia própria (salvo no vértice interno com uma única aresta, que fica colapsado). Os operadores sempre dividem a
estrela exatamente ao longo das arestas selecionadas, de modo que qualquer ordem leva à mesma partição. Verificamos isso
para três ordens do exemplo (a dos slides, a de \code{main.cpp} e o lote inteiro em uma única chamada), para outras
21 ordens nos testes automatizados (\cref{sec:testes}) e, nos 659 casos aleatórios que terminaram
(\cref{sec:verificacao}), para três ordens aleatórias e uma inserção em lote.

\paragraph{Custo.} Cada aresta custa $O\big((\deg A+\deg B)\log n_c\big)$, dominado pelas rotações (\code{isBorder} no
Canônico, reindexação no Expand e no Open Hole); o fator $\log n_c$ vem da busca binária em $Of[\,]$ que toda consulta
$\nxt$/$\prv$ faz. A pré-contagem também chama \code{isBorder}, uma vez por vértice distinto. O custo total é proporcional ao
número de arestas inseridas para malhas de valência limitada, e só depende do tamanho da malha pelo fator logarítmico.

\paragraph{Memória.} O número de vértices novos é exatamente o previsto, inclusive com arestas rejeitadas ou repetidas
no lote. É reservado um elemento por aresta do lote, e o espaço das arestas rejeitadas fica sem uso.

\paragraph{O exemplo no \code{main.cpp}.} Hoje o \code{main.cpp} executa este construtor sobre \code{malha2Hyb.txt},
com as mesmas 12 arestas dos slides em outra ordem (9, 1, 2, 8, 6, 7, 3, 10, 4, 5, 11, 12, na numeração dos slides); o
resultado é o mesmo da \cref{fig:final}.

\section{Verificação experimental}
\label{sec:verificacao}

Para confirmar o entendimento descrito nas seções anteriores, escrevemos um programa de teste
(\code{tools/driver.cpp}, \cref{app:repro}) que usa os arquivos \code{IBHM.cpp} e \code{IBHMInterfaceElementBuilder.cpp}
como estão, insere arestas (uma a uma ou em lote), registra o operador escolhido em cada extremidade e, após cada
inserção, verifica:
\begin{enumerate}[label=\textbf{V\arabic*.}]
\item simetria das opostas: $\Op[\Op[h]]=h$ para toda oposta válida;
\item consistência de vértices entre opostas: $\V[h]=\V[\nxt(\Op[h])]$ e $\V[\nxt(h)]=\V[\Op[h]]$;
\item toda half-edge \COLLAPSED{} tem comprimento zero: $\V[h]=\V[\nxt(h)]$;
\item variedade: os cantos de cada vértice formam um único leque conexo pela relação $h\sim\Op[\prv(h)]$ (nenhum vértice
``estrangulado'' por dois leques separados);
\item todo vértice criado é usado por algum elemento;
\item o número de vértices criados é igual ao previsto por \code{computeNumberOfNewVertices};
\item toda cópia tem exatamente as coordenadas do seu vértice de origem.
\end{enumerate}
As verificações V1--V5 e V7 ficam em \code{tests/support/MeshValidation} e são as mesmas usadas pelos testes
automatizados (\cref{sec:testes}).
Para comparar resultados de ordens diferentes, usamos a partição dos cantos dos elementos \emph{originais} em classes de
vértices: duas execuções produzem a mesma topologia (a menos de renumeração) se e somente se essas partições coincidem.

\subsection{Malhas aleatórias}

O teste com malhas aleatórias (\code{tools/stress.py}) gera grades de $3$ a $7$ por $3$ a $7$ células de quatro tipos
--- só quadriláteros, só triângulos, quadriláteros e triângulos misturados, e misturadas com hexágonos (dois
quadriláteros fundidos) ---, sorteia até metade das arestas internas e insere esse conjunto quatro vezes: em três ordens
aleatórias, escolhendo ao acaso qual half-edge representa cada aresta, e uma vez em lote, numa única chamada. As quatro
execuções precisam ser válidas e produzir a mesma topologia. Foram feitas duas rodadas de 400 malhas
(\cref{tab:estresse}), refeitas com o código desta versão: os números são os mesmos da v0.3, agora também com a V7.

\begin{table}[ht]
\centering\small
\begin{tabularx}{\linewidth}{>{\raggedright}X>{\centering}p{3.3cm}>{\centering\arraybackslash}p{3.3cm}}
\toprule
& Rodada 1: sem restrição & Rodada 2: sem vértice de borda com 2+ arestas\\
\midrule
Malhas (quad / tri / híbrida / polígonos) & 100 / 100 / 100 / 100 & 100 / 100 / 100 / 100\\
Concluídas e válidas (V1--V7) & 259 & \textbf{400}\\
\quad por tipo & 100 / 37 / 56 / 66 & 100 / 100 / 100 / 100\\
Abortadas & 141 & 0\\
\quad das quais com vértice de borda com 2+ arestas & \textbf{141} & ---\\
\quad asserção em \code{openHole} / em \code{heNext} & 72 / 69 & ---\\
Ordens diferentes ou lote com resultado diferente & 0 & 0\\
Contagem de vértices diferente da prevista & 0 & 0\\
Operadores aplicados (Canônico / Split / Expand / Hole) & 3124 / 553 / 384 / 321 & 5130 / 791 / 708 / 535\\
\bottomrule
\end{tabularx}
\caption{Teste com malhas aleatórias. Todas as falhas da rodada 1 são exatamente os casos em que algum vértice de borda
recebe duas ou mais arestas. Nas malhas só de quadriláteros essa configuração não aparece (em uma grade, cada vértice de
borda tem uma única aresta interna); nas de triângulos ela é frequente. As asserções que interrompem a execução são a
do primeiro laço do \code{openHole} e a de \code{heNext}, chamada pelo segundo laço com \BORDER{}
(\cref{sec:falha-borda}).}
\label{tab:estresse}
\end{table}

\subsection{Casos específicos}

A \cref{tab:verificacao} reúne os cenários construídos à mão: o exemplo dos slides em várias ordens, uma trinca reta e
um laço fechado, o experimento que troca Split Element por Expand Edge e um conjunto de entradas inválidas. As linhas
marcadas com \textdagger{} mudaram com as correções desde a v0.3 (I3, I4 e I5).

\begin{table}[ht]
\centering\small
\begin{tabularx}{\linewidth}{>{\raggedright}p{5.6cm}X}
\toprule
Cenário & Resultado\\
\midrule
Exemplo dos slides (ordem dos slides) & válido em todos os passos; 24/24 operadores iguais aos slides; 19 vértices novos
(previsto: 19); buracos em $K$, $L$ e $M$.\\
Exemplo, ordem do \code{main.cpp}, lote único e mais 21 ordens & mesma topologia da ordem dos slides.\\
Coordenadas das cópias (V7)\,\textdagger & toda cópia tem as coordenadas do original, no exemplo e nos 659 casos
aleatórios válidos (até a v0.3 as cópias nasciam em $(0,0)$).\\
Trinca reta $J$--$K$--$L$--$M$--$N$ em 3 ordens & válida; mesma topologia nas 3 ordens.\\
Laço fechado em torno do elemento $FGLK$ & válido; Split Element nos quatro cantos.\\
Split Element substituído por Expand Edge\textsuperscript{a} & 791 ocorrências nas 400 malhas da rodada 2: topologia
idêntica nas 1600 execuções.\\
Mesma aresta duas vezes (em chamadas separadas ou no mesmo lote)\,\textdagger & a segunda é rejeitada; a pré-contagem é
exata (antes reservava vértices que não eram usados).\\
Aresta de borda na entrada\,\textdagger & rejeitada; contada como 0 vértices (antes, 2).\\
Half-edge de um elemento de interface na entrada\,\textdagger & rejeitada, sem alterar a malha (antes, \textbf{laço
infinito} em \code{isBorder}, dentro de \code{computeNumberOfNewVertices}).\\
Índice de half-edge maior que o vetor\,\textdagger & rejeitado, sem alterar a malha (antes, asserção em
\code{heOpposite}; em \emph{release}, leitura fora do vetor).\\
Índice no espaço reservado e ainda não validado\textsuperscript{b}\,\textdagger & rejeitado; a posição reservada de
$\Op$ vale \BORDER{} (antes valia 0, a aresta era aceita e a inserção terminava numa asserção do Open Hole).\\
Malha com furo interno (\code{meshIBHM.txt}): aresta $0$--$4$ e depois $4$--$8$ & a primeira é válida (Canônico de
borda nas duas pontas); a segunda falha no Open Hole, porque o vértice 4 está na borda do furo (I1).\\
\bottomrule
\end{tabularx}
\caption{Cenários específicos. \textsuperscript{a}Variante compilada à parte em que \code{retrieveOperator} nunca
escolhe Split Element. \textsuperscript{b}Por exemplo, a half-edge 48 de \code{malha2Hyb.txt} depois de uma chamada
cuja única aresta foi rejeitada: a reserva criou as posições 48--51, mas nenhum elemento foi validado.}
\label{tab:verificacao}
\end{table}

Em resumo: \emph{para todas as configurações suportadas, a implementação produz topologias válidas, independentes da
ordem e do modo de inserção, com contagem exata de vértices e cópias com as coordenadas do original}. Desde a v0.3, as
entradas inválidas passaram a ser rejeitadas sem alterar a malha. A falha que resta é o caso de borda ainda não tratado,
discutido no fim desta seção e na \cref{sec:defeitos}.

\subsection{Testes automatizados}
\label{sec:testes}

Os cenários acima viraram testes automatizados do repositório (\code{tests/}), escritos com o Google Test e executados
pelo CTest (\cref{app:repro}). Cada teste roda em um processo próprio, com limite de tempo, de modo que uma asserção da
biblioteca ou um laço infinito faz falhar só aquele teste. O núcleo é compilado de novo para os testes com as asserções
ligadas, qualquer que seja o tipo de build, porque alguns defeitos só aparecem como asserção. A \cref{tab:suites} lista
as suítes; o teste com malhas aleatórias também é registrado no CTest.

\begin{table}[ht]
\centering\small
\begin{tabularx}{\linewidth}{>{\ttfamily\small}p{3.6cm}c X}
\toprule
\normalfont Suíte & Testes & O que verifica\\
\midrule
test\_ibhm & 13 & reserva exata e inclusão de pontos, \code{duplicatePoint} (inclusive com realocação durante a cópia),
3 coordenadas, um teste de morte (\code{addPoint} com número errado de coordenadas dispara a asserção),
\code{numberOfHalfEdges} e $\Op[\,]$ reservado igual a \BORDER.\\
test\_vertex\_count & 6 & a pré-contagem: exemplo em lote, aresta repetida, aresta de borda, vértice já em elemento de
interface, vértice de borda com duas arestas (só a contagem).\\
test\_base\_mesh & 5 & o exemplo dos slides: os 24 operadores e os vértices novos de cada inserção, com V1--V7 a cada
passo; malha final; lote igual a uma por vez; duas chamadas iguais a uma; 21 ordens com a half-edge sorteada.\\
test\_invalid\_input & 7 & rejeições sem alterar a malha: índice fora do vetor, índice no espaço reservado, índice de
elemento criado na mesma chamada, aresta de borda, aresta de elemento de interface, aresta já aberta, aresta repetida.\\
test\_known\_defects & 3 & desativados: I1 (uma aresta por chamada e em lote) e I2 (\code{isBorder}). Descrevem o
comportamento correto; hoje, rodados à força, o I1 aborta no Open Hole e o I2 não termina.\\
stress\_supported & 1 & \code{stress.py} com a semente 22 e 400 malhas, sem a configuração do I1 (rodada 2).\\
stress\_all & 1 & rodada 1 (semente 21, todas as configurações); desativado até o I1 ser corrigido.\\
\bottomrule
\end{tabularx}
\caption{Testes automatizados. Os 31 testes ativos das suítes do Google Test levam cerca de 0,2\,s; o
\code{stress\_supported}, cerca de 7\,s.}
\label{tab:suites}
\end{table}

Para saber se os testes detectam os defeitos que já foram corrigidos, as mesmas suítes foram compiladas contra o núcleo
de versões anteriores (\cref{tab:controle}). Cada correção é detectada por mais de um teste, e o núcleo atual passa em
todos.

\begin{table}[ht]
\centering\small
\begin{tabularx}{\linewidth}{>{\raggedright}p{4.6cm}c X}
\toprule
Núcleo & Falhas & Testes que falham\\
\midrule
atual & 0 de 31 & ---\\
antes do I3 (commit \code{3b991e3}) & 4 de 31 & $\Op[\,]$ reservado; índice fora do vetor e índice de elemento da mesma
chamada (asserção em \code{heOpposite}); índice no espaço reservado (asserção no Open Hole).\\
antes do I5 (\code{2a4eb0e}) & 12 de 31 & os anteriores, cinco testes de contagem e três de entrada inválida, sendo
três por laço infinito.\\
antes do I4 (\code{28dd85d}) & 15 de 18\textsuperscript{*} & os mesmos testes de contagem e de entrada inválida e os
quatro testes do exemplo que verificam a V7.\\
\bottomrule
\end{tabularx}
\caption{Controle negativo: as suítes compiladas contra o núcleo de versões anteriores. \textsuperscript{*}Sem
\code{test\_ibhm}, que usa \code{duplicatePoint}, inexistente nessa versão.}
\label{tab:controle}
\end{table}

\subsection{O caso não suportado: vértice de borda com duas arestas}
\label{sec:falha-borda}

Quando um vértice de borda já foi aberto pelo Canônico e recebe uma segunda aresta, ele pertence a um elemento de
interface e não tem aresta colapsada; \code{retrieveOperator} escolhe então Open Hole. Mas um dos dois laços de
\code{openHole} chega à borda da malha antes de encontrar um elemento de interface (\cref{fig:borda}). Em modo
\emph{debug}, a asserção \code{currentHE != IBHM::BORDER} interrompe o programa (os comentários do código perguntam
``\emph{@todo Is it possible?}'': sim, é possível). Em modo \emph{release} o resultado é pior: o primeiro laço escreve
na posição \code{0xFFFFFFFF} do vetor de vértices, e o segundo laço não detecta a borda, porque testa $\nxt(\Op[h])$ em
vez de $\Op[h]$ --- e, na IBHM, $\nxt(\BORDER)$ é uma busca binária com um índice inválido, cujo resultado é
imprevisível.

Esse caso não é raro: aparece sempre que uma trinca chega à borda por um vértice que já tem outra aresta selecionada, o
que é comum em malhas de triângulos (63 das 100 malhas de triângulos da rodada 1) e nos vértices da borda de furos
internos.

\begin{figure}[ht]
\centering
\begin{tikzpicture}[scale=2.200, every node/.style={font=\scriptsize}]
\filldraw[cefill] (0.130,0.075) -- (1.000,0.000) -- (1.200,0.700) -- (0.552,0.725) -- cycle;
\filldraw[cefill] (-0.075,0.130) -- (0.352,0.841) -- (0.000,1.400) -- (-0.500,0.866) -- cycle;
\filldraw[cefill] (-0.075,0.130) -- (-0.500,0.866) -- (-1.200,0.700) -- (-1.000,0.000) -- cycle;
\filldraw[iefill] (0.130,0.075) -- (0.552,0.725) -- (0.352,0.841) -- (-0.075,0.130) -- cycle;
\draw[seledge] (-0.075,0.130) -- (-0.500,0.866);
\draw[hearrow, red!70!black] (-0.306,0.267) -- (-0.544,0.679);
\node[helabel, text=red!70!black] at (-0.549,0.450) {$a$};
\node[vie, minimum size=0.29cm] at (0.130,0.075) {$Q$};
\node[vreg, dot] at (1.000,0.000) {};
\node[vie, minimum size=0.29cm] at (0.352,0.841) {$R$};
\node[vreg, minimum size=0.29cm] at (-0.500,0.866) {$P$};
\node[vreg, dot] at (-1.000,0.000) {};
\node[vreg, dot] at (1.200,0.700) {};
\node[vreg, dot] at (0.000,1.400) {};
\node[vreg, dot] at (-1.200,0.700) {};
\node[vie, minimum size=0.29cm] at (-0.075,0.130) {$Q_{1}$};
\node[vie, minimum size=0.29cm] at (0.552,0.725) {$R_{1}$};

\end{tikzpicture}

\caption{Configuração que falha. $Q$ está na borda e já foi aberto pela aresta $QR$ (cópias $Q$ e $Q_1$; os dois
elementos à esquerda de $QR$ passaram a usar $Q_1$). Ao inserir $QP$, o Open Hole em $Q_1$ gira no sentido
anti-horário a partir de $he$ e alcança imediatamente a borda (à esquerda), sem encontrar elemento de interface. Nesse
caso não há buraco a abrir: bastaria duplicar $Q_1$ e mover para a cópia nova os elementos entre $QP$ e a borda (como no
Canônico de borda), deixando o lado curto do novo elemento na borda.}
\label{fig:borda}
\end{figure}

\section{Documentação \texorpdfstring{$\times$}{x} implementação}
\label{sec:documentacao}

Esta seção confronta as três fontes de documentação --- a apresentação de projeto, os comentários do código e o artigo
da IBHM --- com o que o código faz. Os comentários analisados são os de \code{IBHMInterfaceElementBuilder.h} e
\code{IBHM.h}. A apresentação foi preparada \emph{antes} da implementação, de modo que parte das diferenças é
simplesmente desatualização; ainda assim, vale registrá-las, porque a apresentação é hoje o principal documento do
algoritmo. Cada item tem uma etiqueta: \textbf{[Erro]} (afirma algo falso sobre o código), \textbf{[Imprecisão]}
(incompleto ou ambíguo), \textbf{[Desatualizado]} e \textbf{[Menor]} (redação, digitação).

\subsection{Apresentação de projeto (\code{Operators\_Template})}

\begin{nota}[O que confere]
A decomposição em quatro operadores, as regras de decisão do slide 13, o comportamento do Canônico (aresta colapsada
no interior, vértice novo na borda), o tratamento das arestas do meio (slide 14) e o exemplo completo (slides 15--27)
coincidem com a implementação. Em particular, os 24 operadores anunciados no exemplo são exatamente os escolhidos pelo
código, e o resultado final tem os mesmos buracos ($K$: 4 lados, $L$: 3, $M$: 6) e a mesma aresta colapsada em $F$.
\end{nota}

\begin{enumerate}[label=\textbf{A\arabic*.}, leftmargin=2.3em]
\item \textbf{[Desatualizado]} Nomes: os slides usam \emph{Expand Collapsed Edge} (\code{expandCollapsedEdge}) e
\emph{Hole} (\code{hole(he)}); no código são \code{EXPAND\_EDGE}/\code{expandEdge} e \code{OPEN\_HOLE}/\code{openHole}.
\item \textbf{[Imprecisão]} Slide 6, condição do Split Element: ``\emph{shares an outcome half-edge}''. O código testa
duas hipóteses: um lado do elemento de interface que \emph{sai} do vértice ($x_1$) ou que \emph{chega} nele ($x_3$),
conforme o elemento compartilhado contenha $he$ ou $b$ (\cref{sec:split}). O slide 7 fala em ``\emph{income/outcome}'',
que é o correto. Também vale explicitar que a hipótese depende da posição da aresta nova em relação ao elemento de
interface, não da half-edge escolhida pelo usuário.
\item \textbf{[Imprecisão]} Slide 7: ``\emph{This operator replaces the algorithm that rotates around a vertex}''. Na
implementação os dois coexistem, e o Split Element é um atalho $O(1)$ do Expand Edge que produz a mesma topologia
(\cref{sec:relacoes}). Seria mais preciso dizer que ele \emph{evita} a rotação quando o elemento compartilhado existe.
\item \textbf{[Desatualizado]} Slide 9, Expand: ``\emph{We need to perform a traversal around the vertex starting from
both element's side}''. O código faz uma única rotação anti-horária, que começa na aresta colapsada e para na aresta
nova. A travessia pelos dois lados existe, mas no Open Hole. A preocupação do slide com o produto vetorial também não se
aplica à solução adotada: a rotação é combinatória ($h\leftarrow\Op[\prv(h)]$) e não usa geometria.
\item \textbf{[Imprecisão]} Slides 11--12, Hole: ``\emph{the edge in the interface element is also duplicated for the
new two vertices}'' e ``\emph{both edge's vertices should be handled as in Expand}''. O código não duplica aresta
alguma: ele descola os lados curtos de dois elementos de interface que estavam colados ($\Op\leftarrow\BORDER$), move um
canto de um deles para o vértice novo, e usa duas travessias que param em elementos de interface --- um procedimento
diferente do Expand.
\item \textbf{[Imprecisão]} Slide 13, regra 3 (``\emph{part of an interface element, but not duplicated: this
condition will open a hole}''): não distingue vértices internos de vértices de borda. O código segue a regra ao pé da
letra e aplica Open Hole também na borda, onde não há buraco a abrir e a execução falha (\cref{sec:falha-borda}). A
especificação precisa de um caso para a borda.
\item \textbf{[Imprecisão]} Slide 5: ``\emph{the opposite half-edge for these edges is not defined}''. No código ela é
definida, com o valor sentinela \COLLAPSED{} ($-2$), que as consultas reconhecem (\code{isBorder} pula essas arestas).
Esse estado é uma extensão da estrutura de dados e merece ser documentado como tal.
\item \textbf{[Menor]} Slide 2: ``IBVHM'' em vez de ``IBHM''.
\item \textbf{[Menor]} O slide 27 tem o título ``\emph{Mesh Example -- Edge 11}'', igual ao do slide 26, mas mostra o
estado após a aresta 12.
\item \textbf{[Observação]} Slide 14: ``\emph{all the operators are designed to work on vertices. This property is
desired in order to be easily translated to 3D}''. Em 3D a duplicação acontece também em \emph{arestas} (\cref{sec:3d});
a afirmação vale como objetivo, mas é um ponto de risco.
\end{enumerate}

\subsection{Comentários do código}

\begin{enumerate}[label=\textbf{B\arabic*.}, leftmargin=2.3em]
\item \textbf{[Erro]} \code{openHole}, parâmetro \code{he}: ``\emph{This vertex should be part of a collapsed edge}''.
É o contrário: Open Hole é escolhido justamente quando o vértice \emph{não} tem aresta colapsada (texto copiado de
\code{expandEdge}).
\item \textbf{[Erro]} \code{IBHM::addPoint}: ``\emph{The function supposes that the required space is already
allocated}''. Mas \code{reserveSpaceForNodes} usa \code{resize}, e \code{addPoint} usa \code{push\_back}: as coordenadas
novas vão para depois do bloco reservado, o vetor de coordenadas cresce o dobro do necessário, e
\code{coordinatesFromIndex} de um vértice novo devolve o valor zerado da reserva, não o que foi passado a
\code{addPoint}. Hoje passa despercebido porque todos os pontos novos são criados na origem.
\item \textbf{[Erro]} \code{IBHM::\_halfEdgeVertex}: ``\emph{For each \_numberVerticesByElement vector values, an
element is represented}''. A IBHM não tem esse membro; os elementos são delimitados por \code{\_offset} (texto copiado da
CHE).
\item \textbf{[Erro]} Documentação da classe \code{IBHMInterfaceElementBuilder}: ``\emph{There are only 4 operators to
3D}''. A implementação é 2D; provavelmente a intenção era ``apenas 4 operadores, pensados para serem estendidos a 3D''.
\item \textbf{[Imprecisão]} \code{splitElement}: ``\emph{The interface element with this collapsed edge should have an
edge in the element of he}''. Também pode ser no elemento de $\Op[he]$ (hipótese $x_3$, \cref{sec:split}).
\item \textbf{[Imprecisão]} \code{heOpposite} e \code{\_oppositeHalfEdge} documentam apenas o valor $-1$ (borda); falta
\COLLAPSED{} ($-2$).
\item \textbf{[Imprecisão]} \code{computeNumberOfNewVertices}: ``\emph{n if it is in the border}''. O código soma 1
\emph{por aresta} que toca um vértice de borda, então um vértice de borda com duas arestas no mesmo lote é contado como
3 em vez de 2; arestas que depois serão rejeitadas também são contadas. O efeito é só reservar memória a mais.
\item \textbf{[Imprecisão]} \code{reindexElementsCCW} não documenta a pré-condição de que a rotação não pode atravessar
uma aresta colapsada: se $\Op[\prv(h)]=\COLLAPSED$, o laço usa o sentinela como índice. Nos usos atuais isso não ocorre,
mas é um invariante implícito.
\item \textbf{[Imprecisão]} \code{insertInterfaceElements} não documenta que arestas de borda e de elementos de interface
são rejeitadas, nem que nada é informado a quem chamou; também não diz que as half-edges precisam ser índices válidos de
elementos já existentes.
\item \textbf{[Menor]} Os operadores documentam ``\emph{@return True if the operation was performed successfully and
false otherwise}'', mas sempre devolvem \code{true}, e o valor é ignorado.
\item \textbf{[Menor]} \code{numberReservedElements} diz devolver ``\emph{the number of elements}'', mas devolve o número
de \emph{posições de half-edges} reservadas.
\item \textbf{[Menor]} \code{canonical}: ``\emph{See @reindexElements}'' --- o método se chama
\code{reindexElementsCCW}, e \code{@reindexElements} não é um comando Doxygen (seria \code{@see}).
\item \textbf{[Menor]} \code{\_collapsed}: ``\emph{the value is the half-edge from the first-edge vertex}''. Mais
preciso: o valor é a própria half-edge colapsada ($h_3$ ou $h_1$ do elemento de interface).
\item \textbf{[Menor]} Restos de cópia e digitação: cabeçalho de \code{IBHM.h} com ``\emph{File: CHELevelOne.h}'';
\code{elementsList} cita ``\emph{The CHE class}''; o destrutor está documentado como ``\emph{Default constructor}''; o
construtor cita um parâmetro \code{numberOfVertices} que não existe; \code{buildOppositesLinear} fala em ``mapa do
vértice para os elementos'', mas guarda, em vetores, as half-edges que partem de cada vértice; ``\emph{the bother}''
(border) e ``\emph{the same build should be used}'' (builder).
\end{enumerate}

\subsection{Artigo da IBHM}

\begin{enumerate}[label=\textbf{C\arabic*.}, leftmargin=2.3em]
\item \textbf{[Imprecisão no artigo]} Equação (1): $\mathrm{next}_{he}(he,c):=Of[c]+\mathrm{circ}(he+1,n_v(c))$. Lida
literalmente, com $he$ sendo o índice global, a fórmula só funciona quando $Of[c]$ é múltiplo de $n_v(c)$, o que em
geral não acontece em malhas híbridas. No próprio exemplo da Figura~4 do artigo, a face 2 (quadrilátero) tem $Of=3$ e
half-edges 4 a 7; a fórmula dá $\mathrm{next}(4)=3+\mathrm{circ}(5,4)=4$, quando deveria dar 5. A forma correta aplica
o $\mathrm{circ}$ ao índice local: $Of[c]+\mathrm{circ}(he-Of[c]+1,\,n_v(c))$. A implementação usa a forma correta
(\code{firstHE + (he - firstHE + 1) \% n}).
\item \textbf{[Imprecisão no artigo, menor]} Algoritmo~4 (estrela de um vértice): como está impresso, para um vértice
de borda o laço executa $he\leftarrow\Op[\prv(he)]=0$ e insere $\mathrm{cell}_{he}(0)$ antes de testar a condição de
parada, e a célula de $VH[v]$ só é incluída quando o laço volta a ela (vértices internos). Ao que parece, a intenção é
inserir primeiro a célula de $VH[v]$ e testar $\Op[\prv(he)]=0$ antes de avançar. O algoritmo não é usado pela
implementação.
\item \textbf{[Divergência]} A implementação não tem o vetor $VH[\,]$. Por isso \code{isBorder} precisa girar em torno
do vértice ($O(d\log n_c)$ em vez de $O(1)$), e é nessa rotação que mora o defeito I2 (\cref{sec:defeitos}). Manter $VH[\,]$ atualizado pelos operadores é simples (cada operador sabe quais cantos mudou) e
eliminaria as rotações de teste.
\item \textbf{[Divergência]} Convenções diferentes (índices a partir de zero, $-1$ para borda, $n_c+1$ offsets, estado
\COLLAPSED), resumidas na \cref{tab:estruturas}. Não são erros, mas devem ser registradas para quem comparar código e
artigo.
\item \textbf{[Observação]} O artigo aponta como limitação da IBHM o remalhamento dinâmico. A estratégia
``contar $\to$ reservar $\to$ validar'' da implementação é uma resposta a essa limitação e merece ser destacada em uma
futura publicação.
\end{enumerate}

\subsection{Arquivos de dados}

\begin{enumerate}[label=\textbf{D\arabic*.}, leftmargin=2.3em]
\item \textbf{[Menor]} Em \code{malha2Hyb.txt} o vértice 18 ($T$) tem coordenadas $(3,4)$; nos slides ele está em
$(4,3)$. Com $(3,4)$, o triângulo $M$--$T$--$S$ fica degenerado (três vértices colineares em $x=3$). A topologia não é
afetada, mas qualquer visualização ou verificação geométrica será. As figuras deste relatório usam $(4,3)$.
\item \textbf{[Menor]} \code{meshIBHM.txt} reproduz a Figura~4 do artigo, mas o último triângulo começa em outro
vértice, o que muda a numeração de suas half-edges em relação ao artigo.
\item \textbf{[Menor]} \code{readHybridMesh} (\code{main.cpp}) imprime os offsets a cada elemento lido, o que polui a
saída em malhas grandes.
\end{enumerate}

\subsection{Defeitos e limitações encontrados no código}
\label{sec:defeitos}

A \cref{tab:defeitos} reúne os problemas de \code{IBHMInterfaceElementBuilder} e \code{IBHM} encontrados durante a
análise e confirmados com o programa de teste, com uma sugestão de correção. Como esta é a implementação que será
evoluída, todos devem ser corrigidos; os de gravidade alta impedem o uso em malhas gerais (de triângulos, ou com trincas
que chegam à borda). Nenhuma correção foi aplicada ao código.

\begin{table}[ht]
\centering\small
\setlength{\tabcolsep}{4pt}
\begin{tabularx}{\linewidth}{c c >{\raggedright\arraybackslash}p{3.0cm} X}
\toprule
ID & Grav. & Onde & Defeito e sugestão\\
\midrule
I1 & alta & \code{openHole} & Vértice de borda (externa ou de furo) já aberto recebendo nova aresta: asserção (debug)
ou escrita fora do vetor (release); 141 das 400 malhas da rodada 1. \emph{Sugestão:} permitir que os dois laços parem na
borda; se o laço anti-horário parar na borda, não há canto de elemento de interface a mover; no laço horário, testar
$\Op[h]=\BORDER$ \emph{antes} de aplicar $\nxt$. Nos vértices de borda os lados curtos já são de borda, então nada mais
muda.\\
I2 & alta & \code{isBorder}, \code{insertInterface\allowbreak{}Elements} & Laço infinito quando a half-edge inicial vem
logo após uma aresta colapsada: o desvio ``pula'' exatamente a half-edge de parada. Acontece quando a entrada contém uma
half-edge de elemento de interface, porque a rejeição dessas arestas vem \emph{depois} da pré-contagem. \emph{Sugestão:}
validar a entrada antes de contar e testar a parada também após o desvio.\\
I3 & média & \code{insertInterface\allowbreak{}Elements}, \code{reserveSpace\allowbreak{}ForElements} & Entrada sem
validação de índice: um índice maior que o vetor dispara uma asserção; um índice no espaço reservado e ainda não validado
é aceito, porque as posições reservadas de $\Op[\,]$ começam com 0 (uma half-edge válida). \emph{Sugestão:} aceitar só
$h<Of[n_c]$ e inicializar as posições reservadas de $\Op[\,]$ com \BORDER.\\
I4 & média & \code{addPoint}, \code{reserveSpace\allowbreak{}ForNodes} & Coordenadas novas gravadas fora do lugar (B2).
\emph{Sugestão:} usar \code{reserve} ou escrever na posição do índice; atribuir às cópias a coordenada do vértice de
origem.\\
I5 & baixa & \code{computeNumber\allowbreak{}OfNewVertices} & Contagem em excesso na borda, para arestas rejeitadas e
para arestas repetidas no lote (B7). Só desperdiça memória.\\
I6 & baixa & \code{retrieveAvailable\allowbreak{}Vertex}, \code{openHole}, \code{reindexElements\allowbreak{}CCW} & Os
sentinelas \BORDER{} e \COLLAPSED{} podem chegar a $\nxt$/$\prv$ (B8; segundo laço do Open Hole; \COLLAPSED{} não é
testado em $x_1$/$x_3$). Na IBHM isso vira uma busca binária com índice inválido. \emph{Sugestão:} testes explícitos
antes de cada $\nxt$/$\prv$ aplicado a uma oposta.\\
I7 & baixa & operadores, \code{insertInterface\allowbreak{}Elements} & Não há relato de erro: os operadores sempre
devolvem \code{true} (B10) e o método devolve \code{void}; rejeições só aparecem no console. \emph{Sugestão:} devolver a
lista de arestas inseridas/rejeitadas e um estado de erro.\\
I8 & baixa & \code{insertInterface\allowbreak{}Elements} & A mensagem de operador indefinido para $B$ imprime \code{a}
em vez de \code{b}; \code{\%d} com \code{unsigned}.\\
I9 & manut. & \code{IBHM} & Toda consulta $\nxt$/$\prv$ refaz a busca binária; nas rotações, o elemento corrente já é
conhecido e o intervalo dele poderia ser reaproveitado.\\
I10 & manut. & build & \code{IBHM.cpp} usa \code{memset} e \code{std::upper\_bound} sem incluir \code{<cstring>} e
\code{<algorithm>} (compila por inclusão transitiva, o que não é garantido em outros compiladores);
\code{CMakeLists.txt} exige Qt, que o código atual não usa.\\
\bottomrule
\end{tabularx}
\caption{Defeitos encontrados em \code{IBHMInterfaceElementBuilder} e \code{IBHM}, por gravidade.}
\label{tab:defeitos}
\end{table}

\section{Considerações finais e próximos passos}
\label{sec:conclusao}

O protótipo 2D (\code{IBHMInterfaceElementBuilder}) cumpre o que a apresentação propôs: com quatro operadores locais,
escolhidos apenas pelo estado do vértice, insere elementos de interface em qualquer ordem, sobre malhas híbridas que já
contêm elementos de interface, em uma estrutura compacta baseada em índices e com memória reservada antecipadamente e de
forma exata. A verificação mostrou, em malhas de triângulos, quadriláteros, polígonos e com furos, que o resultado é a
topologia esperada (uma cópia por componente da estrela cortada), válido e independente da ordem e do modo de
inserção, com contagem exata de vértices. Desde a v0.3 foram corrigidos todos os defeitos que afetavam o resultado:
cópias com as coordenadas do original, entrada validada, contagem exata, vértice de borda com várias arestas, laço
infinito latente e sentinelas protegidos por asserções. Os princípios do algoritmo têm testes automatizados próprios, e
o teste de mutação mostrou que eles pegam erros que deixam o resultado válido mas errado. O que resta em aberto é relato
à aplicação (I7), desempenho (I9) e memória temporária na construção (I12).

\subsection{Um invariante que resume o algoritmo}

Uma forma compacta de enunciar o que os operadores constroem, e que pode orientar tanto os testes quanto a extensão 3D,
é a seguinte:

\begin{nota}[Invariante]
Depois de inseridas as arestas de um conjunto $S$, cada vértice original $v$ tem tantas cópias quantos são os
componentes conexos da sua estrela depois de removidas as arestas de $S$. Cada componente usa uma única cópia.
\end{nota}

Para um vértice interno com $k$ arestas selecionadas, isso dá $1$ cópia se $k\le1$ (a ponta de trinca, que fica
colapsada) e $k$ cópias se $k\ge2$; para um vértice de borda, $k+1$ cópias. O buraco que aparece com $k\ge3$ é a
consequência geométrica de ligar $k$ cópias apenas por elementos de espessura nula. Como a partição em componentes não
depende da ordem, qualquer algoritmo correto chega ao mesmo resultado, e é isso que explica a independência da ordem
observada. Os operadores são uma forma \emph{incremental} e local de manter esse invariante.

\subsection{Rumo ao 3D}
\label{sec:3d}

Algumas observações, como hipóteses de trabalho para a próxima etapa:
\begin{itemize}
\item Em 3D a entrada são faces, e o elemento de interface é um prisma (faces triangulares) ou hexaedro (faces
quadrilaterais) de espessura nula. A duplicação deixa de acontecer só nos vértices: as \emph{arestas} das faces
selecionadas também precisam ser duplicadas. Uma aresta compartilhada por $k$ faces selecionadas se comporta como um
vértice 2D com $k$ arestas selecionadas, e girar em torno de uma aresta em 3D (de célula em célula pelas half-faces
opostas) é o análogo exato de girar em torno de um vértice em 2D. Os operadores 2D devem reaparecer, quase sem
mudanças, como operadores \emph{de aresta}.
\item A estrela de um vértice em 3D não é mais um ciclo, e sim uma superfície fechada (ou com borda) de células. Dividir
essa estrela exige uma busca em largura pelas células que não atravessa faces selecionadas, e não uma rotação. O
invariante acima continua valendo e dá o critério: uma cópia por componente.
\item O estado colapsado tem análogos em 3D (faces laterais degeneradas do prisma), e o buraco das junções vira
``tubos'' ao longo das arestas de junção.
\item O sentinela \COLLAPSED, a pré-contagem e a estratégia contar/reservar/validar se transportam diretamente para a
3D-IBHM.
\end{itemize}

\subsection{Próximos passos sugeridos}

\begin{enumerate}
\item \textbf{Extensão 3D} com os testes como base. O oráculo da topologia esperada (\cref{sec:testes}) se transporta
quase sem mudanças: uma cópia por componente da estrela de cada vértice (e de cada aresta) cortada pelas faces
selecionadas. O gerador de malhas aleatórias, as verificações de estrutura (gabarito do elemento, decisão pelo estado,
contrato de cada operador, localidade) e a equivalência entre o operador particular e o geral (Split/Expand) podem ser
reescritos para 3D antes dos operadores, servindo de especificação executável.
\item \textbf{Relatar resultados} (I7): uma entrada por aresta de entrada, com o elemento criado ou o motivo da
rejeição; a biblioteca deixa de imprimir. Retomar junto o B10.
\item \textbf{Transferência de estado}: expor o mapa ``vértice novo $\to$ vértice de origem'', no mesmo formato. Numa
análise adaptativa, os campos nodais (deslocamentos, poropressões etc.) dos vértices novos precisam ser copiados do
vértice de origem; o construtor conhece essa informação no momento da criação (as coordenadas já são copiadas, I4).
\item \textbf{Simplificar os operadores} em torno de uma primitiva comum de ``divisão de leque'' (\cref{sec:relacoes}),
mantendo o Split Element como otimização.
\item \textbf{Desempenho e memória} (I9, I12): manter o vetor $VH[\,]$, reaproveitar o intervalo do elemento corrente
nas rotações e montar as opostas sem um vetor por vértice.
\item \textbf{Integração com o GeMA}.
\end{enumerate}


\begin{thebibliography}{9}\small
\bibitem{ibhm} M.~Lage, L.~F. Martha, J.~P. Moitinho de Almeida, H.~Lopes.
  \emph{IBHM: index-based data structures for 2D and 3D hybrid meshes}.
  Engineering with Computers 33:727--744, 2017. DOI 10.1007/s00366-015-0395-0.
\bibitem{che} M.~Lage, T.~Lewiner, H.~Lopes, L.~Velho. \emph{CHE: a scalable topological data structure for triangular
  meshes}. Relatório técnico, PUC-Rio, 2005.
\bibitem{chf} M.~Lage, T.~Lewiner, H.~Lopes, L.~Velho. \emph{CHF: a scalable topological data structure for
  tetrahedral meshes}. SIBGRAPI 2005, pp.~349--356.
\bibitem{mantyla} M.~Mäntylä. \emph{An Introduction to Solid Modeling}. Computer Science Press, 1988.
\bibitem{tops} W.~Celes, G.~H. Paulino, R.~Espinha. \emph{A compact adjacency-based topological data structure for
  finite element mesh representation}. Int. J. Numer. Methods Eng. 64:1529--1556, 2005.
\bibitem{paulino} G.~H. Paulino, W.~Celes, R.~Espinha, Z.~Zhang. \emph{A general topology-based framework for adaptive
  insertion of cohesive elements in finite element meshes}. Engineering with Computers 24:59--78, 2008.
\bibitem{slides} J.~Coelho. \emph{Adaptive Interface Elements Generation -- 2D Case}. Apresentação (Tecgraf/PUC-Rio),
  versão de 25/02/2026. Arquivos \code{doc/ppt/Operators\_Template.\{pptx,pdf\}}.
\end{thebibliography}

\appendix
\section{Testes, ferramentas e geração das figuras}
\label{app:repro}

Os testes automatizados ficam em \code{tests/} e são compilados pelo CMake; as ferramentas do relatório ficam em
\code{doc/report/tools/}. Todos usam o código do repositório como está e podem ser rodados a partir da raiz do
repositório. Cada arquivo começa com um comentário que descreve seu uso; este apêndice explica o que cada um faz e como
os resultados do relatório foram obtidos.

\begin{center}\small
\begin{tabularx}{\linewidth}{>{\ttfamily\small}p{4.6cm}X}
\toprule
\normalfont Arquivo & Papel\\
\midrule
tests/test\_*.cpp & suítes do Google Test (\cref{tab:suites});\\
tests/support/ & malhas de teste, verificações V1--V7 e acesso de leitura ao estado do construtor, usados pelos testes e
pelo \code{driver};\\
tests/CMakeLists.txt & registra as suítes e o teste com malhas aleatórias no CTest;\\
tools/driver.cpp & programa de teste da \code{IBHMInterfaceElementBuilder}: insere arestas, verifica a malha e grava cada
estado em JSON;\\
tools/compare.py & compara a topologia final de execuções diferentes;\\
tools/stress.py & teste com malhas aleatórias, várias ordens e inserção em lote;\\
tools/gen.py & gera as figuras de malha do relatório (\code{doc/report/figs/}).\\
\bottomrule
\end{tabularx}
\end{center}

\subsection{Como compilar e rodar}

Os comandos abaixo partem da raiz do repositório; \code{\$W} é um diretório de trabalho qualquer, fora do repositório.
O CMake não depende de Qt. O Google Test vem do pacote instalado, se houver; senão, a versão 1.17.0 é baixada na
configuração, com o hash SHA256 fixado no \code{CMakeLists.txt}. (No macOS, o \code{cmake} e o \code{ninja} do CLion
ficam em \code{/Applications/CLion.app/Contents/bin/}.) O tipo de build pode ser \code{Debug} ou \code{Release}: nos
dois, os testes rodam com as asserções ligadas.
\begin{lstlisting}[language=bash]
cmake -S . -B $W/build -G Ninja -DCMAKE_BUILD_TYPE=Debug
cmake --build $W/build
cd $W/build && ctest -j8                      # ou: ctest -L unit; ctest -L stress
D=$W/build/tests/driver
$D malha2Hyb.txt $W/ex.json 31 34 11 26 47 38 23 19 30 42 29 45         # exemplo dos slides
python3 doc/report/tools/stress.py $D 21 400 $W/s1                       # rodada 1
python3 doc/report/tools/stress.py $D 22 400 $W/s2 --skip-border-multi   # rodada 2
python3 doc/report/tools/gen.py $D doc/report/figs $W/work               # figuras
cd doc/report && tectonic -X compile relatorio.tex                       # relatório
\end{lstlisting}
Sem CMake, o \code{driver} também pode ser compilado diretamente:
\begin{lstlisting}[language=bash]
c++ -std=c++17 -Icore/TopologyDataStructures -Icore/InterfaceElementsBuilder -Itests/support \
    doc/report/tools/driver.cpp tests/support/MeshValidation.cpp \
    core/TopologyDataStructures/IBHM.cpp \
    core/InterfaceElementsBuilder/IBHMInterfaceElementBuilder.cpp -o $W/driver
\end{lstlisting}

\subsection{Os testes automatizados (\code{tests/})}

O alvo \code{InterfaceElementsTests} reúne as suítes da \cref{tab:suites}, e \code{gtest\_discover\_tests} registra
cada \code{TEST} como um teste do CTest, com o rótulo \code{unit} e limite de 60\,s. Como cada teste roda num processo
próprio, uma asserção ou um laço infinito faz falhar apenas aquele teste. O teste com malhas aleatórias entra como
\code{stress\_supported} (rótulo \code{stress}), que chama o \code{stress.py} com o \code{driver}; o
\code{stress\_all} está registrado, mas desativado até o I1 ser corrigido.

\textbf{Asserções ligadas.} Alguns defeitos só aparecem como asserção, e o teste de morte de \code{addPoint} espera uma.
Por isso o \code{tests/CMakeLists.txt} retira o \code{NDEBUG} das opções do diretório e compila o núcleo de novo
(\code{InterfaceElementsLibChecked}), qualquer que seja o tipo de build; a biblioteca \code{InterfaceElementsLib}
continua compilada normalmente.

\textbf{Código de apoio} (\code{tests/support/}):
\begin{itemize}
\item \code{TestMeshes}: a malha base (\code{malha2Hyb.txt}, com o vértice $T$ em $(4,3)$) e a lista das 12 arestas do
exemplo, a malha em leque do caso de borda com duas arestas e uma malha de dois triângulos;
\item \code{MeshValidation}: a ``foto'' da malha (lida pela interface pública da \code{IBHM}, com $\nxt$/$\prv$
recalculados), as verificações V1--V5, a origem de cada cópia e a V7, e a partição dos cantos usada para comparar
topologias;
\item \code{BuilderInspector}: acesso de leitura ao estado privado do construtor (o operador que seria escolhido,
\code{\_isInterfaceElement}, \code{\_inInterfaceElement} e \code{\_collapsed}). O construtor a declara como classe
amiga. A alternativa usada até a v0.3, \code{\#define private public}, foi abandonada: no MSVC o nível de acesso faz
parte do nome decorado das funções, e a ligação falha.
\end{itemize}

\textbf{Defeitos em aberto.} Os testes de \code{test\_known\_defects.cpp} têm o prefixo \code{DISABLED\_}: aparecem como
desativados em toda execução e descrevem o comportamento correto, de modo que, corrigido o defeito, basta tirar o
prefixo. Para rodá-los mesmo assim: \code{InterfaceElementsTests --gtest\_also\_run\_disabled\_tests
--gtest\_filter='KnownDefects*'}.

\subsection{O programa de teste (\code{driver.cpp})}

\textbf{Uso:} \code{driver [--batch] <malha> <saida.json> <he0> <he1> ...}. A malha está no formato de
\code{malha2Hyb.txt}: o número de vértices e de elementos, as coordenadas e, para cada elemento, o número de vértices
seguido dos índices. Cada \code{he} é uma half-edge da aresta a abrir (qualquer uma das duas). Sem \code{--batch}, cada
aresta é inserida numa chamada de \code{insertInterfaceElements}; com \code{--batch}, todas vão numa única chamada.

\textbf{O que faz em cada chamada:}
\begin{enumerate}
\item tira uma ``foto'' da malha (vértices de cada elemento, offsets e opostas), lida pela interface pública da
\code{IBHM};
\item se a aresta vai ser aceita, pergunta a \code{retrieveOperator} qual operador será usado em cada extremidade, só para
registrá-lo (a decisão não muda nada na malha);
\item chama \code{computeNumberOfNewVertices} e \code{insertInterfaceElements};
\item descobre o vértice de origem de cada vértice novo: todo vértice novo substitui, em algum canto de um elemento que
já existia, o vértice que estava lá antes; esse vértice antigo (ou a origem dele) é a origem do novo;
\item aplica as verificações V1--V5 e V7 da \cref{sec:verificacao} e compara o número de vértices criados com o
previsto (V6);
\item imprime uma linha e grava o estado no JSON.
\end{enumerate}
As verificações são feitas sobre a ``foto'', com consultas $\nxt$/$\prv$ recalculadas, para não depender do código que
está sendo testado; a foto, as verificações e a origem das cópias vêm de \code{tests/support/MeshValidation}, as mesmas
dos testes automatizados. O estado privado do construtor (\code{\_collapsed}, \code{\_isInterfaceElement} e o operador
escolhido) é lido por \code{BuilderInspector}. Na v0.4 o \code{driver} passou a usar esse código compartilhado; sua saída
(console e JSON) continuou idêntica, byte a byte, no exemplo, nos casos de entrada inválida e em 834 execuções do
estresse (120 malhas de cada rodada, inclusive as que abortam).

\textbf{Saída no console:} uma linha por chamada, como
\begin{lstlisting}[language={}]
step  2 edge he=34 (10-11)  A:SPLIT_ELEMENT B:CANONICAL     new=1 predicted=1 ok
\end{lstlisting}
com a half-edge, os dois vértices da aresta, o operador de cada extremidade, os vértices criados e previstos e o
resultado (\code{ok} ou \code{INVALID}, seguido das verificações que falharam). O código de saída é 0 se tudo passou e
2 se alguma verificação falhou; uma asserção da biblioteca interrompe o programa.

\textbf{JSON:} um objeto com \code{points} (coordenadas originais) e \code{steps}, uma lista em que o passo 0 é a malha de
entrada e o passo $k$ é o estado depois da $k$-ésima chamada. Cada passo tem:

\begin{center}\small
\begin{tabularx}{\linewidth}{>{\ttfamily\small}p{3.6cm}X}
\toprule
\normalfont Campo & Conteúdo\\
\midrule
edge, A, B & a half-edge inserida e os vértices da sua aresta\\
opA, opB & o operador de cada extremidade (vazio no modo em lote ou para arestas rejeitadas)\\
predicted, created & vértices previstos por \code{computeNumberOfNewVertices} e vértices criados\\
errors, msgs & número de problemas e mensagens das verificações\\
elements, first & os vértices de cada elemento e a primeira half-edge de cada elemento (os offsets)\\
opp & a oposta de cada half-edge; $-1$ é \BORDER{} e $-2$ é \COLLAPSED\\
isIE & 1 para os elementos de interface\\
root & para cada vértice, o vértice original de que ele é cópia\\
collapsed & o mapa \code{\_collapsed} do construtor, como pares (vértice, half-edge)\\
\bottomrule
\end{tabularx}
\end{center}
O arquivo é gravado e esvaziado (\emph{flush}) a cada passo; se um passo posterior aborta, os anteriores continuam
legíveis (o \code{gen.py} usa isso para desenhar o caso de borda que falha).

\subsection{Comparação de topologias (\code{compare.py})}

Duas execuções que inserem as mesmas arestas em ordens diferentes criam os vértices novos em ordens diferentes, então os
índices não servem para comparar. O \code{compare.py} olha só para os elementos que já existiam na malha de entrada
(eles mantêm seus índices) e agrupa os cantos desses elementos pelo vértice que usam. O conjunto de grupos é uma partição
dos cantos; duas execuções têm a mesma topologia, a menos de renumeração, se e somente se as partições são iguais. Os
elementos de interface ficam determinados pelos elementos contínuos em volta e não precisam entrar na comparação. Pela
linha de comando, \code{python3 compare.py a.json b.json ...} diz, para cada arquivo, se a partição é igual à do
primeiro.

\subsection{Teste com malhas aleatórias (\code{stress.py})}

\textbf{Uso:} \code{stress.py <driver> <semente> <casos> <diretório> [--skip-border-multi]}. Cada caso:
\begin{enumerate}
\item gera uma malha numa grade de $3$ a $7$ por $3$ a $7$ células, alternando quatro tipos (\cref{fig:estresse-malhas}):
só quadriláteros, só triângulos (cada célula dividida por uma das diagonais, ao acaso), híbrida (60\% das células viram
triângulos) e com polígonos (como a híbrida, mas 25\% das vezes duas células vizinhas viram um hexágono);
\item sorteia até metade das arestas internas;
\item insere o mesmo conjunto quatro vezes com o \code{driver}: três vezes uma aresta por chamada, em ordens aleatórias e
com a half-edge de cada aresta escolhida ao acaso, e uma vez em lote;
\item exige que as quatro execuções sejam válidas, que a contagem de vértices seja sempre a prevista e que as quatro
partições sejam iguais.
\end{enumerate}
Um caso é marcado como ``vértice de borda com 2+ arestas'' quando a seleção tem um vértice de borda tocado por duas ou
mais arestas (a configuração do defeito I1). Com \code{--skip-border-multi}, essas seleções são sorteadas de novo. No
fim, o script imprime os contadores (casos válidos, inválidos e abortados, quantos abortados tinham a configuração do I1,
casos dependentes da ordem, casos em que o lote difere, contagens erradas, operadores aplicados), a asserção que causou
cada aborto e o número de casos válidos por tipo de malha. O código de saída é 0 se todos os casos passaram e 1 se
houve aborto, caso inválido, dependência da ordem ou do lote, ou contagem errada; é ele que o CTest usa. A semente fixa
o gerador aleatório: a mesma semente sempre gera as mesmas malhas, seleções e ordens (o relatório usa as sementes 21 e
22, com 400 casos cada).

\begin{figure}[ht]
\centering
\begin{subfigure}[b]{0.24\linewidth}\centering\resizebox{\linewidth}{!}{\begin{tikzpicture}[scale=0.620, every node/.style={font=\scriptsize}]
\filldraw[cefill] (0.000,0.000) -- (1.000,0.000) -- (1.000,1.000) -- (0.000,1.000) -- cycle;
\filldraw[cefill] (1.000,0.000) -- (2.000,0.000) -- (2.000,1.000) -- (1.000,1.000) -- cycle;
\filldraw[cefill] (2.000,0.000) -- (3.000,0.000) -- (3.000,1.000) -- (2.000,1.000) -- cycle;
\filldraw[cefill] (3.000,0.000) -- (4.000,0.000) -- (4.000,1.000) -- (3.000,1.000) -- cycle;
\filldraw[cefill] (4.000,0.000) -- (5.000,0.000) -- (5.000,1.000) -- (4.000,1.000) -- cycle;
\filldraw[cefill] (5.000,0.000) -- (6.000,0.000) -- (6.000,1.000) -- (5.000,1.000) -- cycle;
\filldraw[cefill] (0.000,1.000) -- (1.000,1.000) -- (1.000,2.000) -- (0.000,2.000) -- cycle;
\filldraw[cefill] (1.000,1.000) -- (2.000,1.000) -- (2.000,2.000) -- (1.000,2.000) -- cycle;
\filldraw[cefill] (2.000,1.000) -- (3.000,1.000) -- (3.000,2.000) -- (2.000,2.000) -- cycle;
\filldraw[cefill] (3.000,1.000) -- (4.000,1.000) -- (4.000,2.000) -- (3.000,2.000) -- cycle;
\filldraw[cefill] (4.000,1.000) -- (5.000,1.000) -- (5.000,2.000) -- (4.000,2.000) -- cycle;
\filldraw[cefill] (5.000,1.000) -- (6.000,1.000) -- (6.000,2.000) -- (5.000,2.000) -- cycle;
\filldraw[cefill] (0.000,2.000) -- (1.000,2.000) -- (1.000,3.000) -- (0.000,3.000) -- cycle;
\filldraw[cefill] (1.000,2.000) -- (2.000,2.000) -- (2.000,3.000) -- (1.000,3.000) -- cycle;
\filldraw[cefill] (2.000,2.000) -- (3.000,2.000) -- (3.000,3.000) -- (2.000,3.000) -- cycle;
\filldraw[cefill] (3.000,2.000) -- (4.000,2.000) -- (4.000,3.000) -- (3.000,3.000) -- cycle;
\filldraw[cefill] (4.000,2.000) -- (5.000,2.000) -- (5.000,3.000) -- (4.000,3.000) -- cycle;
\filldraw[cefill] (5.000,2.000) -- (6.000,2.000) -- (6.000,3.000) -- (5.000,3.000) -- cycle;
\filldraw[cefill] (0.000,3.000) -- (1.000,3.000) -- (1.000,4.000) -- (0.000,4.000) -- cycle;
\filldraw[cefill] (1.000,3.000) -- (2.000,3.000) -- (2.000,4.000) -- (1.000,4.000) -- cycle;
\filldraw[cefill] (2.000,3.000) -- (3.000,3.000) -- (3.000,4.000) -- (2.000,4.000) -- cycle;
\filldraw[cefill] (3.000,3.000) -- (4.000,3.000) -- (4.000,4.000) -- (3.000,4.000) -- cycle;
\filldraw[cefill] (4.000,3.000) -- (5.000,3.000) -- (5.000,4.000) -- (4.000,4.000) -- cycle;
\filldraw[cefill] (5.000,3.000) -- (6.000,3.000) -- (6.000,4.000) -- (5.000,4.000) -- cycle;
\draw[seledge] (2.000,2.000) -- (2.000,3.000);
\draw[seledge] (2.000,0.000) -- (2.000,1.000);
\draw[seledge] (1.000,2.000) -- (1.000,3.000);
\draw[seledge] (2.000,3.000) -- (2.000,4.000);
\draw[seledge] (4.000,2.000) -- (5.000,2.000);
\draw[seledge] (0.000,3.000) -- (1.000,3.000);
\draw[seledge] (5.000,2.000) -- (5.000,3.000);
\draw[seledge] (4.000,2.000) -- (4.000,3.000);
\draw[seledge] (2.000,3.000) -- (3.000,3.000);
\node[vreg, dot] at (0.000,0.000) {};
\node[vreg, dot] at (1.000,0.000) {};
\node[vreg, dot] at (2.000,0.000) {};
\node[vreg, dot] at (3.000,0.000) {};
\node[vreg, dot] at (4.000,0.000) {};
\node[vreg, dot] at (5.000,0.000) {};
\node[vreg, dot] at (6.000,0.000) {};
\node[vreg, dot] at (0.000,1.000) {};
\node[vreg, dot] at (1.000,1.000) {};
\node[vreg, dot] at (2.000,1.000) {};
\node[vreg, dot] at (3.000,1.000) {};
\node[vreg, dot] at (4.000,1.000) {};
\node[vreg, dot] at (5.000,1.000) {};
\node[vreg, dot] at (6.000,1.000) {};
\node[vreg, dot] at (0.000,2.000) {};
\node[vreg, dot] at (1.000,2.000) {};
\node[vreg, dot] at (2.000,2.000) {};
\node[vreg, dot] at (3.000,2.000) {};
\node[vreg, dot] at (4.000,2.000) {};
\node[vreg, dot] at (5.000,2.000) {};
\node[vreg, dot] at (6.000,2.000) {};
\node[vreg, dot] at (0.000,3.000) {};
\node[vreg, dot] at (1.000,3.000) {};
\node[vreg, dot] at (2.000,3.000) {};
\node[vreg, dot] at (3.000,3.000) {};
\node[vreg, dot] at (4.000,3.000) {};
\node[vreg, dot] at (5.000,3.000) {};
\node[vreg, dot] at (6.000,3.000) {};
\node[vreg, dot] at (0.000,4.000) {};
\node[vreg, dot] at (1.000,4.000) {};
\node[vreg, dot] at (2.000,4.000) {};
\node[vreg, dot] at (3.000,4.000) {};
\node[vreg, dot] at (4.000,4.000) {};
\node[vreg, dot] at (5.000,4.000) {};
\node[vreg, dot] at (6.000,4.000) {};

\end{tikzpicture}
}\caption{quad}\end{subfigure}\hfill
\begin{subfigure}[b]{0.24\linewidth}\centering\resizebox{\linewidth}{!}{\begin{tikzpicture}[scale=0.620, every node/.style={font=\scriptsize}]
\filldraw[cefill] (0.000,0.000) -- (1.000,0.000) -- (0.000,1.000) -- cycle;
\filldraw[cefill] (1.000,0.000) -- (1.000,1.000) -- (0.000,1.000) -- cycle;
\filldraw[cefill] (1.000,0.000) -- (2.000,0.000) -- (1.000,1.000) -- cycle;
\filldraw[cefill] (2.000,0.000) -- (2.000,1.000) -- (1.000,1.000) -- cycle;
\filldraw[cefill] (2.000,0.000) -- (3.000,0.000) -- (3.000,1.000) -- cycle;
\filldraw[cefill] (2.000,0.000) -- (3.000,1.000) -- (2.000,1.000) -- cycle;
\filldraw[cefill] (3.000,0.000) -- (4.000,0.000) -- (3.000,1.000) -- cycle;
\filldraw[cefill] (4.000,0.000) -- (4.000,1.000) -- (3.000,1.000) -- cycle;
\filldraw[cefill] (4.000,0.000) -- (5.000,0.000) -- (5.000,1.000) -- cycle;
\filldraw[cefill] (4.000,0.000) -- (5.000,1.000) -- (4.000,1.000) -- cycle;
\filldraw[cefill] (5.000,0.000) -- (6.000,0.000) -- (6.000,1.000) -- cycle;
\filldraw[cefill] (5.000,0.000) -- (6.000,1.000) -- (5.000,1.000) -- cycle;
\filldraw[cefill] (0.000,1.000) -- (1.000,1.000) -- (1.000,2.000) -- cycle;
\filldraw[cefill] (0.000,1.000) -- (1.000,2.000) -- (0.000,2.000) -- cycle;
\filldraw[cefill] (1.000,1.000) -- (2.000,1.000) -- (2.000,2.000) -- cycle;
\filldraw[cefill] (1.000,1.000) -- (2.000,2.000) -- (1.000,2.000) -- cycle;
\filldraw[cefill] (2.000,1.000) -- (3.000,1.000) -- (2.000,2.000) -- cycle;
\filldraw[cefill] (3.000,1.000) -- (3.000,2.000) -- (2.000,2.000) -- cycle;
\filldraw[cefill] (3.000,1.000) -- (4.000,1.000) -- (3.000,2.000) -- cycle;
\filldraw[cefill] (4.000,1.000) -- (4.000,2.000) -- (3.000,2.000) -- cycle;
\filldraw[cefill] (4.000,1.000) -- (5.000,1.000) -- (5.000,2.000) -- cycle;
\filldraw[cefill] (4.000,1.000) -- (5.000,2.000) -- (4.000,2.000) -- cycle;
\filldraw[cefill] (5.000,1.000) -- (6.000,1.000) -- (5.000,2.000) -- cycle;
\filldraw[cefill] (6.000,1.000) -- (6.000,2.000) -- (5.000,2.000) -- cycle;
\filldraw[cefill] (0.000,2.000) -- (1.000,2.000) -- (1.000,3.000) -- cycle;
\filldraw[cefill] (0.000,2.000) -- (1.000,3.000) -- (0.000,3.000) -- cycle;
\filldraw[cefill] (1.000,2.000) -- (2.000,2.000) -- (2.000,3.000) -- cycle;
\filldraw[cefill] (1.000,2.000) -- (2.000,3.000) -- (1.000,3.000) -- cycle;
\filldraw[cefill] (2.000,2.000) -- (3.000,2.000) -- (2.000,3.000) -- cycle;
\filldraw[cefill] (3.000,2.000) -- (3.000,3.000) -- (2.000,3.000) -- cycle;
\filldraw[cefill] (3.000,2.000) -- (4.000,2.000) -- (3.000,3.000) -- cycle;
\filldraw[cefill] (4.000,2.000) -- (4.000,3.000) -- (3.000,3.000) -- cycle;
\filldraw[cefill] (4.000,2.000) -- (5.000,2.000) -- (4.000,3.000) -- cycle;
\filldraw[cefill] (5.000,2.000) -- (5.000,3.000) -- (4.000,3.000) -- cycle;
\filldraw[cefill] (5.000,2.000) -- (6.000,2.000) -- (5.000,3.000) -- cycle;
\filldraw[cefill] (6.000,2.000) -- (6.000,3.000) -- (5.000,3.000) -- cycle;
\filldraw[cefill] (0.000,3.000) -- (1.000,3.000) -- (1.000,4.000) -- cycle;
\filldraw[cefill] (0.000,3.000) -- (1.000,4.000) -- (0.000,4.000) -- cycle;
\filldraw[cefill] (1.000,3.000) -- (2.000,3.000) -- (2.000,4.000) -- cycle;
\filldraw[cefill] (1.000,3.000) -- (2.000,4.000) -- (1.000,4.000) -- cycle;
\filldraw[cefill] (2.000,3.000) -- (3.000,3.000) -- (3.000,4.000) -- cycle;
\filldraw[cefill] (2.000,3.000) -- (3.000,4.000) -- (2.000,4.000) -- cycle;
\filldraw[cefill] (3.000,3.000) -- (4.000,3.000) -- (4.000,4.000) -- cycle;
\filldraw[cefill] (3.000,3.000) -- (4.000,4.000) -- (3.000,4.000) -- cycle;
\filldraw[cefill] (4.000,3.000) -- (5.000,3.000) -- (5.000,4.000) -- cycle;
\filldraw[cefill] (4.000,3.000) -- (5.000,4.000) -- (4.000,4.000) -- cycle;
\filldraw[cefill] (5.000,3.000) -- (6.000,3.000) -- (5.000,4.000) -- cycle;
\filldraw[cefill] (6.000,3.000) -- (6.000,4.000) -- (5.000,4.000) -- cycle;
\draw[seledge] (3.000,2.000) -- (4.000,2.000);
\draw[seledge] (1.000,1.000) -- (2.000,1.000);
\draw[seledge] (4.000,3.000) -- (5.000,3.000);
\draw[seledge] (4.000,1.000) -- (3.000,2.000);
\draw[seledge] (5.000,2.000) -- (6.000,2.000);
\draw[seledge] (1.000,0.000) -- (1.000,1.000);
\draw[seledge] (3.000,3.000) -- (4.000,3.000);
\draw[seledge] (2.000,1.000) -- (3.000,1.000);
\draw[seledge] (4.000,2.000) -- (3.000,3.000);
\draw[seledge] (5.000,2.000) -- (5.000,3.000);
\draw[seledge] (6.000,2.000) -- (5.000,3.000);
\draw[seledge] (5.000,0.000) -- (6.000,1.000);
\draw[seledge] (2.000,3.000) -- (3.000,3.000);
\draw[seledge] (3.000,1.000) -- (4.000,1.000);
\draw[seledge] (2.000,2.000) -- (3.000,2.000);
\node[vreg, dot] at (0.000,0.000) {};
\node[vreg, dot] at (1.000,0.000) {};
\node[vreg, dot] at (2.000,0.000) {};
\node[vreg, dot] at (3.000,0.000) {};
\node[vreg, dot] at (4.000,0.000) {};
\node[vreg, dot] at (5.000,0.000) {};
\node[vreg, dot] at (6.000,0.000) {};
\node[vreg, dot] at (0.000,1.000) {};
\node[vreg, dot] at (1.000,1.000) {};
\node[vreg, dot] at (2.000,1.000) {};
\node[vreg, dot] at (3.000,1.000) {};
\node[vreg, dot] at (4.000,1.000) {};
\node[vreg, dot] at (5.000,1.000) {};
\node[vreg, dot] at (6.000,1.000) {};
\node[vreg, dot] at (0.000,2.000) {};
\node[vreg, dot] at (1.000,2.000) {};
\node[vreg, dot] at (2.000,2.000) {};
\node[vreg, dot] at (3.000,2.000) {};
\node[vreg, dot] at (4.000,2.000) {};
\node[vreg, dot] at (5.000,2.000) {};
\node[vreg, dot] at (6.000,2.000) {};
\node[vreg, dot] at (0.000,3.000) {};
\node[vreg, dot] at (1.000,3.000) {};
\node[vreg, dot] at (2.000,3.000) {};
\node[vreg, dot] at (3.000,3.000) {};
\node[vreg, dot] at (4.000,3.000) {};
\node[vreg, dot] at (5.000,3.000) {};
\node[vreg, dot] at (6.000,3.000) {};
\node[vreg, dot] at (0.000,4.000) {};
\node[vreg, dot] at (1.000,4.000) {};
\node[vreg, dot] at (2.000,4.000) {};
\node[vreg, dot] at (3.000,4.000) {};
\node[vreg, dot] at (4.000,4.000) {};
\node[vreg, dot] at (5.000,4.000) {};
\node[vreg, dot] at (6.000,4.000) {};

\end{tikzpicture}
}\caption{tri}\end{subfigure}\hfill
\begin{subfigure}[b]{0.24\linewidth}\centering\resizebox{\linewidth}{!}{\begin{tikzpicture}[scale=0.620, every node/.style={font=\scriptsize}]
\filldraw[cefill] (0.000,0.000) -- (1.000,0.000) -- (0.000,1.000) -- cycle;
\filldraw[cefill] (1.000,0.000) -- (1.000,1.000) -- (0.000,1.000) -- cycle;
\filldraw[cefill] (1.000,0.000) -- (2.000,0.000) -- (1.000,1.000) -- cycle;
\filldraw[cefill] (2.000,0.000) -- (2.000,1.000) -- (1.000,1.000) -- cycle;
\filldraw[cefill] (2.000,0.000) -- (3.000,0.000) -- (3.000,1.000) -- (2.000,1.000) -- cycle;
\filldraw[cefill] (3.000,0.000) -- (4.000,0.000) -- (4.000,1.000) -- cycle;
\filldraw[cefill] (3.000,0.000) -- (4.000,1.000) -- (3.000,1.000) -- cycle;
\filldraw[cefill] (4.000,0.000) -- (5.000,0.000) -- (5.000,1.000) -- (4.000,1.000) -- cycle;
\filldraw[cefill] (5.000,0.000) -- (6.000,0.000) -- (6.000,1.000) -- cycle;
\filldraw[cefill] (5.000,0.000) -- (6.000,1.000) -- (5.000,1.000) -- cycle;
\filldraw[cefill] (0.000,1.000) -- (1.000,1.000) -- (1.000,2.000) -- (0.000,2.000) -- cycle;
\filldraw[cefill] (1.000,1.000) -- (2.000,1.000) -- (1.000,2.000) -- cycle;
\filldraw[cefill] (2.000,1.000) -- (2.000,2.000) -- (1.000,2.000) -- cycle;
\filldraw[cefill] (2.000,1.000) -- (3.000,1.000) -- (2.000,2.000) -- cycle;
\filldraw[cefill] (3.000,1.000) -- (3.000,2.000) -- (2.000,2.000) -- cycle;
\filldraw[cefill] (3.000,1.000) -- (4.000,1.000) -- (3.000,2.000) -- cycle;
\filldraw[cefill] (4.000,1.000) -- (4.000,2.000) -- (3.000,2.000) -- cycle;
\filldraw[cefill] (4.000,1.000) -- (5.000,1.000) -- (5.000,2.000) -- (4.000,2.000) -- cycle;
\filldraw[cefill] (5.000,1.000) -- (6.000,1.000) -- (5.000,2.000) -- cycle;
\filldraw[cefill] (6.000,1.000) -- (6.000,2.000) -- (5.000,2.000) -- cycle;
\filldraw[cefill] (0.000,2.000) -- (1.000,2.000) -- (1.000,3.000) -- (0.000,3.000) -- cycle;
\filldraw[cefill] (1.000,2.000) -- (2.000,2.000) -- (1.000,3.000) -- cycle;
\filldraw[cefill] (2.000,2.000) -- (2.000,3.000) -- (1.000,3.000) -- cycle;
\filldraw[cefill] (2.000,2.000) -- (3.000,2.000) -- (3.000,3.000) -- cycle;
\filldraw[cefill] (2.000,2.000) -- (3.000,3.000) -- (2.000,3.000) -- cycle;
\filldraw[cefill] (3.000,2.000) -- (4.000,2.000) -- (3.000,3.000) -- cycle;
\filldraw[cefill] (4.000,2.000) -- (4.000,3.000) -- (3.000,3.000) -- cycle;
\filldraw[cefill] (4.000,2.000) -- (5.000,2.000) -- (4.000,3.000) -- cycle;
\filldraw[cefill] (5.000,2.000) -- (5.000,3.000) -- (4.000,3.000) -- cycle;
\filldraw[cefill] (5.000,2.000) -- (6.000,2.000) -- (6.000,3.000) -- (5.000,3.000) -- cycle;
\filldraw[cefill] (0.000,3.000) -- (1.000,3.000) -- (1.000,4.000) -- (0.000,4.000) -- cycle;
\filldraw[cefill] (1.000,3.000) -- (2.000,3.000) -- (2.000,4.000) -- (1.000,4.000) -- cycle;
\filldraw[cefill] (2.000,3.000) -- (3.000,3.000) -- (2.000,4.000) -- cycle;
\filldraw[cefill] (3.000,3.000) -- (3.000,4.000) -- (2.000,4.000) -- cycle;
\filldraw[cefill] (3.000,3.000) -- (4.000,3.000) -- (3.000,4.000) -- cycle;
\filldraw[cefill] (4.000,3.000) -- (4.000,4.000) -- (3.000,4.000) -- cycle;
\filldraw[cefill] (4.000,3.000) -- (5.000,3.000) -- (5.000,4.000) -- (4.000,4.000) -- cycle;
\filldraw[cefill] (5.000,3.000) -- (6.000,3.000) -- (6.000,4.000) -- cycle;
\filldraw[cefill] (5.000,3.000) -- (6.000,4.000) -- (5.000,4.000) -- cycle;
\draw[seledge] (2.000,1.000) -- (2.000,2.000);
\draw[seledge] (3.000,1.000) -- (3.000,2.000);
\draw[seledge] (2.000,3.000) -- (3.000,3.000);
\draw[seledge] (2.000,2.000) -- (1.000,3.000);
\draw[seledge] (1.000,3.000) -- (1.000,4.000);
\draw[seledge] (4.000,1.000) -- (4.000,2.000);
\draw[seledge] (1.000,2.000) -- (2.000,2.000);
\draw[seledge] (3.000,2.000) -- (4.000,2.000);
\draw[seledge] (0.000,2.000) -- (1.000,2.000);
\draw[seledge] (5.000,2.000) -- (4.000,3.000);
\draw[seledge] (5.000,1.000) -- (6.000,1.000);
\draw[seledge] (4.000,2.000) -- (4.000,3.000);
\draw[seledge] (5.000,2.000) -- (5.000,3.000);
\node[vreg, dot] at (0.000,0.000) {};
\node[vreg, dot] at (1.000,0.000) {};
\node[vreg, dot] at (2.000,0.000) {};
\node[vreg, dot] at (3.000,0.000) {};
\node[vreg, dot] at (4.000,0.000) {};
\node[vreg, dot] at (5.000,0.000) {};
\node[vreg, dot] at (6.000,0.000) {};
\node[vreg, dot] at (0.000,1.000) {};
\node[vreg, dot] at (1.000,1.000) {};
\node[vreg, dot] at (2.000,1.000) {};
\node[vreg, dot] at (3.000,1.000) {};
\node[vreg, dot] at (4.000,1.000) {};
\node[vreg, dot] at (5.000,1.000) {};
\node[vreg, dot] at (6.000,1.000) {};
\node[vreg, dot] at (0.000,2.000) {};
\node[vreg, dot] at (1.000,2.000) {};
\node[vreg, dot] at (2.000,2.000) {};
\node[vreg, dot] at (3.000,2.000) {};
\node[vreg, dot] at (4.000,2.000) {};
\node[vreg, dot] at (5.000,2.000) {};
\node[vreg, dot] at (6.000,2.000) {};
\node[vreg, dot] at (0.000,3.000) {};
\node[vreg, dot] at (1.000,3.000) {};
\node[vreg, dot] at (2.000,3.000) {};
\node[vreg, dot] at (3.000,3.000) {};
\node[vreg, dot] at (4.000,3.000) {};
\node[vreg, dot] at (5.000,3.000) {};
\node[vreg, dot] at (6.000,3.000) {};
\node[vreg, dot] at (0.000,4.000) {};
\node[vreg, dot] at (1.000,4.000) {};
\node[vreg, dot] at (2.000,4.000) {};
\node[vreg, dot] at (3.000,4.000) {};
\node[vreg, dot] at (4.000,4.000) {};
\node[vreg, dot] at (5.000,4.000) {};
\node[vreg, dot] at (6.000,4.000) {};

\end{tikzpicture}
}\caption{hybrid}\end{subfigure}\hfill
\begin{subfigure}[b]{0.24\linewidth}\centering\resizebox{\linewidth}{!}{\begin{tikzpicture}[scale=0.620, every node/.style={font=\scriptsize}]
\filldraw[cefill] (0.000,0.000) -- (1.000,0.000) -- (2.000,0.000) -- (2.000,1.000) -- (1.000,1.000) -- (0.000,1.000) -- cycle;
\filldraw[cefill] (2.000,0.000) -- (3.000,0.000) -- (3.000,1.000) -- cycle;
\filldraw[cefill] (2.000,0.000) -- (3.000,1.000) -- (2.000,1.000) -- cycle;
\filldraw[cefill] (3.000,0.000) -- (4.000,0.000) -- (4.000,1.000) -- (3.000,1.000) -- cycle;
\filldraw[cefill] (4.000,0.000) -- (5.000,0.000) -- (5.000,1.000) -- (4.000,1.000) -- cycle;
\filldraw[cefill] (5.000,0.000) -- (6.000,0.000) -- (6.000,1.000) -- cycle;
\filldraw[cefill] (5.000,0.000) -- (6.000,1.000) -- (5.000,1.000) -- cycle;
\filldraw[cefill] (0.000,1.000) -- (1.000,1.000) -- (1.000,2.000) -- (0.000,2.000) -- cycle;
\filldraw[cefill] (1.000,1.000) -- (2.000,1.000) -- (2.000,2.000) -- cycle;
\filldraw[cefill] (1.000,1.000) -- (2.000,2.000) -- (1.000,2.000) -- cycle;
\filldraw[cefill] (2.000,1.000) -- (3.000,1.000) -- (3.000,2.000) -- (2.000,2.000) -- cycle;
\filldraw[cefill] (3.000,1.000) -- (4.000,1.000) -- (3.000,2.000) -- cycle;
\filldraw[cefill] (4.000,1.000) -- (4.000,2.000) -- (3.000,2.000) -- cycle;
\filldraw[cefill] (4.000,1.000) -- (5.000,1.000) -- (4.000,2.000) -- cycle;
\filldraw[cefill] (5.000,1.000) -- (5.000,2.000) -- (4.000,2.000) -- cycle;
\filldraw[cefill] (5.000,1.000) -- (6.000,1.000) -- (5.000,2.000) -- cycle;
\filldraw[cefill] (6.000,1.000) -- (6.000,2.000) -- (5.000,2.000) -- cycle;
\filldraw[cefill] (0.000,2.000) -- (1.000,2.000) -- (1.000,3.000) -- (0.000,3.000) -- cycle;
\filldraw[cefill] (1.000,2.000) -- (2.000,2.000) -- (1.000,3.000) -- cycle;
\filldraw[cefill] (2.000,2.000) -- (2.000,3.000) -- (1.000,3.000) -- cycle;
\filldraw[cefill] (2.000,2.000) -- (3.000,2.000) -- (3.000,3.000) -- (2.000,3.000) -- cycle;
\filldraw[cefill] (3.000,2.000) -- (4.000,2.000) -- (5.000,2.000) -- (5.000,3.000) -- (4.000,3.000) -- (3.000,3.000) -- cycle;
\filldraw[cefill] (5.000,2.000) -- (6.000,2.000) -- (6.000,3.000) -- (5.000,3.000) -- cycle;
\filldraw[cefill] (0.000,3.000) -- (1.000,3.000) -- (1.000,4.000) -- cycle;
\filldraw[cefill] (0.000,3.000) -- (1.000,4.000) -- (0.000,4.000) -- cycle;
\filldraw[cefill] (1.000,3.000) -- (2.000,3.000) -- (3.000,3.000) -- (3.000,4.000) -- (2.000,4.000) -- (1.000,4.000) -- cycle;
\filldraw[cefill] (3.000,3.000) -- (4.000,3.000) -- (4.000,4.000) -- (3.000,4.000) -- cycle;
\filldraw[cefill] (4.000,3.000) -- (5.000,3.000) -- (4.000,4.000) -- cycle;
\filldraw[cefill] (5.000,3.000) -- (5.000,4.000) -- (4.000,4.000) -- cycle;
\filldraw[cefill] (5.000,3.000) -- (6.000,3.000) -- (6.000,4.000) -- (5.000,4.000) -- cycle;
\draw[seledge] (3.000,2.000) -- (3.000,3.000);
\draw[seledge] (2.000,2.000) -- (2.000,3.000);
\draw[seledge] (2.000,3.000) -- (3.000,3.000);
\draw[seledge] (3.000,2.000) -- (4.000,2.000);
\draw[seledge] (1.000,1.000) -- (2.000,1.000);
\draw[seledge] (1.000,2.000) -- (2.000,2.000);
\draw[seledge] (0.000,1.000) -- (1.000,1.000);
\draw[seledge] (2.000,0.000) -- (3.000,1.000);
\draw[seledge] (4.000,3.000) -- (4.000,4.000);
\draw[seledge] (5.000,2.000) -- (6.000,2.000);
\draw[seledge] (3.000,1.000) -- (4.000,1.000);
\node[vreg, dot] at (0.000,0.000) {};
\node[vreg, dot] at (1.000,0.000) {};
\node[vreg, dot] at (2.000,0.000) {};
\node[vreg, dot] at (3.000,0.000) {};
\node[vreg, dot] at (4.000,0.000) {};
\node[vreg, dot] at (5.000,0.000) {};
\node[vreg, dot] at (6.000,0.000) {};
\node[vreg, dot] at (0.000,1.000) {};
\node[vreg, dot] at (1.000,1.000) {};
\node[vreg, dot] at (2.000,1.000) {};
\node[vreg, dot] at (3.000,1.000) {};
\node[vreg, dot] at (4.000,1.000) {};
\node[vreg, dot] at (5.000,1.000) {};
\node[vreg, dot] at (6.000,1.000) {};
\node[vreg, dot] at (0.000,2.000) {};
\node[vreg, dot] at (1.000,2.000) {};
\node[vreg, dot] at (2.000,2.000) {};
\node[vreg, dot] at (3.000,2.000) {};
\node[vreg, dot] at (4.000,2.000) {};
\node[vreg, dot] at (5.000,2.000) {};
\node[vreg, dot] at (6.000,2.000) {};
\node[vreg, dot] at (0.000,3.000) {};
\node[vreg, dot] at (1.000,3.000) {};
\node[vreg, dot] at (2.000,3.000) {};
\node[vreg, dot] at (3.000,3.000) {};
\node[vreg, dot] at (4.000,3.000) {};
\node[vreg, dot] at (5.000,3.000) {};
\node[vreg, dot] at (6.000,3.000) {};
\node[vreg, dot] at (0.000,4.000) {};
\node[vreg, dot] at (1.000,4.000) {};
\node[vreg, dot] at (2.000,4.000) {};
\node[vreg, dot] at (3.000,4.000) {};
\node[vreg, dot] at (4.000,4.000) {};
\node[vreg, dot] at (5.000,4.000) {};
\node[vreg, dot] at (6.000,4.000) {};

\end{tikzpicture}
}\caption{poly}\end{subfigure}
\caption{Um exemplo de cada tipo de malha gerada pelo \code{stress.py} (grade de $6\times4$ células), com uma seleção
de arestas em vermelho. As malhas e as seleções do teste usam o mesmo gerador, com tamanhos e sementes variados.}
\label{fig:estresse-malhas}
\end{figure}

\subsection{Geração das figuras (\code{gen.py})}

\textbf{Uso:} \code{gen.py <driver> <diretório-das-figuras> <diretório-de-trabalho>}. Para cada figura, o script monta
uma malha pequena (ou lê uma do repositório), roda o \code{driver}, carrega o estado desejado e o converte num
\code{tikzpicture}, gravado em \code{figs/<nome>.tex}. Assim, toda figura de malha mostra o que o código realmente
produz; as figuras não devem ser editadas à mão, e sim regeneradas. Três detalhes do desenho:
\begin{itemize}
\item \textbf{Vista explodida.} As cópias de um vértice coincidem por definição (o código lhes dá as coordenadas do
original). Para a topologia ficar visível, cada cópia de um vértice dividido é desenhada um pouco afastada da posição
original, na direção média dos centróides dos elementos contínuos que a usam; como cada cópia é usada pelos elementos
de um setor, as cópias se afastam, cada uma para o seu setor.
\item \textbf{Arestas colapsadas.} Os dois cantos de um elemento de interface sobre uma aresta colapsada são afastados na
direção do elemento vizinho, e a aresta é desenhada como um pequeno segmento vermelho com espessura apenas visual.
\item \textbf{Buracos.} As half-edges de borda formam laços; um laço percorrido no sentido horário (área com sinal
negativa) é um buraco e é hachurado.
\end{itemize}
Os nomes dos vértices seguem a regra das figuras: o vértice original mantém o nome e suas cópias recebem índices
($K$, $K_1$, $K_2$, \dots) na ordem de criação. As figuras desenhadas à mão (gabarito, árvore de decisão, estados,
hipóteses $x_1$/$x_3$) ficam no próprio texto, em TikZ, e usam os estilos de \code{styles.tex}.


\end{document}
